% Lecture notes on the classification of homogeneous plane continua,
% following the formal (Lean 4) proof in this repository.
% Volume 2: Uniqueness and homogeneity of the pseudo-arc.
%
% Compile with:  pdflatex main && pdflatex main   (two runs for the contents and references)
\documentclass[11pt,oneside]{book}

\usepackage[T1]{fontenc}
\usepackage{lmodern}
\usepackage[margin=1.15in]{geometry}
\usepackage{amsmath,amssymb,amsthm}
\usepackage{enumitem}
\usepackage{booktabs}
\usepackage{longtable}
\usepackage{tikz}
\usepackage[hidelinks]{hyperref}

% ---------------------------------------------------------------- theorem styles
\theoremstyle{plain}
\newtheorem{theorem}{Theorem}[chapter]
\newtheorem{lemma}[theorem]{Lemma}
\newtheorem{proposition}[theorem]{Proposition}
\newtheorem{corollary}[theorem]{Corollary}
\newtheorem{claim}[theorem]{Claim}
\theoremstyle{definition}
\newtheorem{definition}[theorem]{Definition}
\newtheorem{example}[theorem]{Example}
\newtheorem{exercise}[theorem]{Exercise}
\theoremstyle{remark}
\newtheorem{remark}[theorem]{Remark}

% ---------------------------------------------------------------- notation
\newcommand{\N}{\mathbb{N}}
\newcommand{\Z}{\mathbb{Z}}
\newcommand{\R}{\mathbb{R}}
\newcommand{\diam}{\operatorname{diam}}
\newcommand{\dist}{\operatorname{dist}}
\newcommand{\cl}[1]{\overline{#1}}
\newcommand{\abs}[1]{\left\lvert #1\right\rvert}
\newcommand{\set}[1]{\left\{ #1 \right\}}
\newcommand{\fl}[1]{\left\lfloor #1\right\rfloor}
% the composant of p in M
\newcommand{\cmp}[2]{\operatorname{cmp}_{#1}(#2)}
% the union of the links a..b of a chain D
\newcommand{\spn}[3]{#1[#2,#3]}
% least and greatest coarse index of a fine link
\newcommand{\lo}{\rho}
\newcommand{\hi}{\sigma}
% envelope and arrival level of a leg
\newcommand{\env}{\operatorname{Env}}
\newcommand{\arr}{\operatorname{arr}}

% Names of declarations in the Lean development (typewriter, underscores allowed).
\newcommand{\lean}[1]{\texttt{\detokenize{#1}}}
% A Lean name ending in a subscript digit (such as Tracks with subscript 1).
\newcommand{\leansub}[2]{\texttt{\detokenize{#1}}\ensuremath{_{#2}}}
% A paragraph pointing to the formal counterpart of a result.
\newcommand{\formal}[1]{\par\smallskip\noindent{\small\textit{Formal counterpart:} #1}\par\smallskip}
% A reference to a numbered result of Volume 1.
\newcommand{\vone}[1]{Volume~1, #1}

\setlist{itemsep=2pt,topsep=4pt}

\title{\textbf{Homogeneous Plane Continua}\\[6pt]
\Large Lecture notes following the formal proof\\[18pt]
\large Volume 2\\[4pt]
\Large Uniqueness and homogeneity of the pseudo-arc}
\author{}
\date{}

\begin{document}

\frontmatter
\maketitle
\tableofcontents

\mainmatter
% The introduction is Chapter 0, as in Volume 1.
\setcounter{chapter}{-1}
%======================================================================
\chapter{Introduction}\label{ch:intro}
%======================================================================

\section{What this volume proves}

Volume~1 constructed a pseudo-arc in the plane: a nondegenerate continuum that is snake-like
(it can be covered by chains with arbitrarily small links) and hereditarily indecomposable
(every subcontinuum is indecomposable). This volume proves that there is essentially only one
pseudo-arc, and that it is homogeneous. These are the two facts about the pseudo-arc that the
classification of homogeneous plane continua needs.

Recall that the \emph{composant} of a point $p$ in a continuum $M$ is the union of all proper
subcontinua of $M$ that contain $p$ (Chapter~\ref{ch:composants}).

\begin{theorem}[Uniqueness of the pseudo-arc]\label{thm:intro-unique}
Any two pseudo-arcs are homeomorphic. More precisely, let $M$ and $M'$ be pseudo-arcs (in
possibly different metric spaces), let $p,q\in M$ lie in different composants of $M$, and let
$p',q'\in M'$ lie in different composants of $M'$. Then there is a homeomorphism $h\colon M\to
M'$ with $h(p)=p'$ and $h(q)=q'$.
\end{theorem}

\begin{theorem}[Homogeneity of the pseudo-arc]\label{thm:intro-homog}
For any two points $p,q$ of a pseudo-arc $M$ there is a homeomorphism $h\colon M\to M$ with
$h(p)=q$.
\end{theorem}

\begin{corollary}\label{cor:intro}
\begin{enumerate}[label=(\alph*)]
\item Every pseudo-arc is homeomorphic to the pseudo-arc in the plane constructed in
Volume~1.
\item Every nondegenerate subcontinuum of a pseudo-arc is a pseudo-arc, and hence is
homeomorphic to the whole pseudo-arc.
\end{enumerate}
\end{corollary}

These results are due to R.~H.~Bing (1948 and 1951). Theorem~\ref{thm:intro-unique} for
the continua described by sequences of crooked chains is Theorem~12 of his 1948 paper; the
fact that every pseudo-arc has such a description, so that the theorem applies to all
pseudo-arcs, is Theorem~1 of his 1951 paper ``Concerning hereditarily indecomposable
continua''. As in Volume~1, the notes follow the formal proof rather than the papers; the main
differences are listed in Section~\ref{sec:intro-diff}.

\section{The shape of the proof}

The proof has three parts.

\paragraph{Part I: two descriptions give a homeomorphism (Chapters~\ref{ch:patterns}--\ref{ch:two}).}
A \emph{description} from $P$ to $Q$ is a sequence of chains $D^{(0)},D^{(1)},\dots$ from $P$ to
$Q$, each crooked in the one before, with closures of links nested and mesh tending to zero
(\vone{Definition~3.1}). Its \emph{described set} is $M=\bigcap_i\bigcup D^{(i)}$. We prove:

\begin{quote}
\emph{If two descriptions, from $P_1$ to $Q_1\ne P_1$ and from $P_2$ to $Q_2\neq P_2$, live in
compact metric spaces, then there is a homeomorphism between the described sets taking $P_1$
to $P_2$ and $Q_1$ to $Q_2$.}
\end{quote}

The idea is to build, alternately in the two spaces, two sequences of chains with the same
numbers of links which ``follow the same patterns'' in each other. Two ingredients are
needed.
\begin{itemize}
\item \emph{Pattern following} (Chapter~\ref{ch:patterns}): given any pattern, that is, any walk
$x$ through the links of a level $D^{(i)}$ from its first link to its last, some deeper level
can be grouped into a chain that follows $x$. The proof is a purely combinatorial statement
about integer sequences, proved by induction, and a compactness argument.
\item \emph{Matching covers} (Chapter~\ref{ch:matching}): if two compact sets carry sequences of
finite open covers, indexed in the same way, such that ``the $j$-th set of one level meets the
$k$-th set of the level before'' on one side forces the corresponding sets on the other side to
be close, then the two compact sets are homeomorphic. This is a general metric statement with
no chains in it.
\end{itemize}
Chapter~\ref{ch:two} combines them.

\paragraph{Part II: every pseudo-arc has a description (Chapters~\ref{ch:folding}--\ref{ch:descriptions}).}
Let $M$ be a pseudo-arc and let $p,q\in M$ lie in different composants. We prove:

\begin{quote}
\emph{In the compact space $M$ there is a description from $p$ to $q$ whose described set is
all of $M$.}
\end{quote}

The proof has two halves. First (Chapter~\ref{ch:descriptions}, following Bing's 1951
argument) hereditary indecomposability forces fine chains covering $M$ to be crooked in coarse
ones; this gives a sequence of chains, each crooked in the previous one, but not yet going from
$p$ to $q$. Second, each chain must be replaced by one that goes from $p$ to $q$. For this one
needs a theorem about crooked chains alone:

\begin{itemize}
\item \emph{The folding theorem} (Chapter~\ref{ch:folding}, Bing's Theorem~4): if a chain $E$ is
crooked in a chain $D$, and a stretch of $E$ from a link $r$ to a link $s$ reaches both end links
of $D$, then the links of $E$ can be grouped into consecutive blocks forming a chain, with $r$ in
the first block and $s$ in the last, and every block inside three consecutive links of $D$.
\end{itemize}

This is the longest proof of the volume. It is again purely combinatorial: it is a statement
about an integer sequence (the indices of the links of $D$ containing the successive links of
$E$) that moves by at most one at each step and is ``crooked''.

\paragraph{Part III: composants and the main theorems (Chapters~\ref{ch:composants} and
\ref{ch:main}).}
To apply Parts~I and~II one needs two points in different composants. Chapter~\ref{ch:composants}
proves that in an indecomposable continuum no composant is the whole continuum; so every point
$p$ of a pseudo-arc has a point $q$ outside its composant. Chapter~\ref{ch:main} assembles the
proofs of Theorems~\ref{thm:intro-unique} and~\ref{thm:intro-homog} and of
Corollary~\ref{cor:intro}.

\medskip
The logical dependence of the chapters is:
\begin{center}
\begin{tikzpicture}[x=2.6cm,y=1.1cm,every node/.style={draw,rounded corners,inner sep=3pt,font=\small}]
\node (c1) at (0,2) {Ch.~\ref{ch:composants} Composants};
\node (c2) at (2,2) {Ch.~\ref{ch:consolidations} Consolidations};
\node (c3) at (1,1) {Ch.~\ref{ch:patterns} Patterns};
\node (c4) at (0,1) {Ch.~\ref{ch:matching} Matching};
\node (c5) at (0.5,0) {Ch.~\ref{ch:two} Two descriptions};
\node (c6) at (3,1) {Ch.~\ref{ch:folding} Folding};
\node (c7) at (2.5,0) {Ch.~\ref{ch:descriptions} Descriptions};
\node (c8) at (1.5,-1) {Ch.~\ref{ch:main} Main theorems};
\draw[->] (c2) -- (c3); \draw[->] (c2) -- (c6); \draw[->] (c3) -- (c5); \draw[->] (c4) -- (c5);
\draw[->] (c6) -- (c7); \draw[->] (c2) -- (c7); \draw[->] (c1) -- (c7);
\draw[->] (c5) -- (c8); \draw[->] (c7) -- (c8); \draw[->] (c1) -- (c8);
\end{tikzpicture}
\end{center}

\section{Where the formal proof differs from the papers}\label{sec:intro-diff}

The formal proof follows Bing's architecture. It departs from his text in the following
places, and the notes follow the formal proof.

\begin{enumerate}
\item \textbf{Pattern following needs $P\neq Q$.} Bing's Theorem~6 (pattern following) is false
for a description of a one-point space, where $P=Q$ (Remark~\ref{rem:pf-false}); the formal
proof adds the hypothesis $P\ne Q$, which is harmless since the uniqueness theorem assumes it
anyway.
\item \textbf{Pattern following is proved directly}, from a combinatorial statement about
integer sequences (Theorem~\ref{thm:realize}) whose induction is designed to be checkable. Bing
derives Theorem~6 from a more general Theorem~5, whose proof leaves a step to the reader.
\item \textbf{Compactness of the ambient space.} Bing's uniqueness theorem asks only that the
described sets be compact. That is not enough: there is a description, in a non-complete
space, whose described set is compact but not connected (\vone{Remark~3.5}). The formal
statement asks that the links lie in a compact set (equivalently, after restricting, that the
ambient space be compact).
\item \textbf{The folding theorem with three links.} Bing's Theorem~4 asks that every block lie
in the union of \emph{two} adjacent links of $D$, and derives it from his Theorem~3. The formal
proof proves the version with \emph{three} consecutive links directly, by a ``greedy zigzag''
construction (Chapter~\ref{ch:folding}), which is all that the later arguments use: they only
need the blocks to be small. It also proves the version, needed for chains from a point $R$ to
a point $S$, in which $R$ may lie in two adjacent links.
\item \textbf{Crookedness of fine chains} (Bing 1951, first paragraph of the proof of
Theorem~1) is proved by a limit argument whose witnesses are chosen so that the two limit
continua are visibly not nested (Section~\ref{sec:crooked-eventually}).
\end{enumerate}

\section{Prerequisites and conventions}

The reader should know Volume~1, in particular Chapters~1--3: chains, refinements, crookedness,
descriptions, and the theorem that a described set is a pseudo-arc. References to Volume~1 are
written ``\vone{Lemma~2.5}''. The numbering of Volume~1 starts with Chapter~0 (the
introduction), as here. We use the conventions of Volume~1:

\begin{itemize}
\item $\N=\set{0,1,2,\dots}$, and $[n]=\set{0,1,\dots,n-1}$. Chains are indexed from $0$.
\item A \emph{chain} of length $n$ is a sequence $D=(D_0,\dots,D_{n-1})$ of open sets (links)
with $D_i\cap D_j\neq\emptyset$ if and only if $\abs{i-j}\le1$ (\vone{Definition~2.1}). A chain
\emph{goes from $P$ to $Q$} if $P$ lies in the first link and in no other, and $Q$ lies in the
last link and in no other.
\item ``Crooked'' always means crooked in the sense of \vone{Definition~2.11}: $E$ is crooked
in $D$ if $E$ refines $D$ and, whenever $k>h+2$ and $E_i$ meets $D_h$ and $E_j$ meets $D_k$, there
are $r,s$ strictly between $i$ and $j$, with $r$ nearer to $i$, such that $E_r\subseteq D_{k-1}$
and $E_s\subseteq D_{h+1}$.
\item A function $g$ defined on a set of consecutive integers has \emph{unit steps} if
$\abs{g(t+1)-g(t)}\le1$ whenever both sides are defined.
\end{itemize}

\paragraph{Formal counterparts.} As in Volume~1, many results are followed by a line
``Formal counterpart:'' naming the corresponding declaration of the Lean development (all names
begin with \lean{HPC.}). These lines may be skipped. Appendix~\ref{app:formal} collects them.

\paragraph{Exercises.} Each chapter ends with exercises. None of them is used later.

%======================================================================
\chapter{Composants}\label{ch:composants}
%======================================================================

The uniqueness theorem compares two pseudo-arcs through two chosen points of each, and these
points must lie in different composants. This chapter proves that such points exist: in an
indecomposable continuum no composant is the whole continuum. It also records how the
notions of Volume~1 behave when a continuum $M\subseteq X$ is regarded as a metric space in
its own right; Part~II of the proof takes place inside that space. Throughout, $X$ is a metric
space.

\section{Composants}

\begin{definition}\label{def:composant}
Let $M\subseteq X$ and $p\in X$. The \emph{composant} of $p$ in $M$ is
\[
\cmp{M}{p}=\bigcup\set{K : K\text{ is a continuum},\ K\subseteq M,\ K\neq M,\ p\in K}.
\]
\end{definition}

\formal{\lean{HPC.composant} (in \lean{HPC/PseudoArc/Statements79.lean}).}

If $M$ is a nondegenerate continuum and $p\in M$, then $\set p$ is a proper subcontinuum of $M$,
so $p\in\cmp{M}{p}$. In particular:

\begin{lemma}\label{lem:ne-of-composant}
Let $M$ be a set with at least two points, $p\in M$, and $q\notin\cmp{M}{p}$. Then $p\neq q$.
\end{lemma}

\begin{proof}
$\set p$ is compact and connected, lies in $M$, contains $p$, and is not $M$ because $M$ has
two points. So $p\in\cmp Mp$, and $q\notin\cmp Mp$ gives $q\neq p$.
\end{proof}

\formal{\lean{HPC.Plane.PA51.ne_of_not_mem_composant}.}

\begin{example}
Let $M=[0,1]$. For $0<p<1$ every $y\in[0,1]$ lies in one of the proper subcontinua $[0,p]$ and
$[p,1]$, both containing $p$; so $\cmp Mp=M$. For $p=0$, every $y<1$ lies in the proper subcontinuum
$[0,y]$, while a subcontinuum containing $0$ and $1$ is all of $[0,1]$; so $\cmp M0=[0,1)$.
\end{example}

The example shows that for decomposable continua a composant may be everything. The goal of
this chapter is:

\begin{theorem}\label{thm:composant-proper}
Let $M$ be an indecomposable continuum and $p\in M$. Then $\cmp Mp\neq M$; that is, there is a
point $q\in M$ which lies in no proper subcontinuum of $M$ containing $p$.
\end{theorem}

The proof (Section~\ref{sec:baire}) shows that $\cmp Mp$ is contained in a countable union of
proper subcontinua of $M$, and that a proper subcontinuum of an indecomposable continuum is
``thin'': it contains no nonempty relatively open subset of $M$. The Baire category theorem
then finishes the proof.

\section{Boundary bumping}

We need two facts from point-set topology. In both, a subset $W$ of a set $L$ is called
\emph{closed and relatively open in $L$} if $W$ is closed in $X$ and $W=L\cap O$ for some open
set $O\subseteq X$.

\begin{lemma}[Boundary bumping, clopen form]\label{lem:bumping}
Let $M$ be preconnected, $L\subseteq M$ with $L\neq M$, and let $W\subseteq L$ be nonempty, closed,
and relatively open in $L$. Then $W$ meets $\cl{M\setminus L}$.
\end{lemma}

\begin{proof}
Write $W=L\cap O$ with $O$ open, and suppose $W\cap\cl{M\setminus L}=\emptyset$. Consider the open
sets
\[
U=O\setminus\cl{M\setminus L},\qquad V=X\setminus W .
\]
They cover $M$: a point $y\in M$ either lies outside $W$, so in $V$, or lies in $W$, so in $O$ and
(by assumption) outside $\cl{M\setminus L}$, so in $U$. The set $M\cap U$ contains any point of the
nonempty set $W$. The set $M\cap V$ contains any point $x\in M\setminus L$ (there is one since
$L\neq M$), because $W\subseteq L$. Since $M$ is preconnected, some $z\in M$ lies in $U\cap V$. As
$z\in U$, $z\notin\cl{M\setminus L}$, so $z\notin M\setminus L$, that is, $z\in L$. Also $z\in O$,
so $z\in L\cap O=W$, contradicting $z\in V$.
\end{proof}

\formal{\lean{HPC.T13.clopen_meets_closure_diff}.}

\begin{lemma}\label{lem:union-clopen}
Let $K$ be closed and preconnected, let $L_1\subseteq L$ be closed and relatively open in $L$, and
suppose that every nonempty subset of $L$ that is closed and relatively open in $L$ meets $K$.
Then $K\cup L_1$ is preconnected.
\end{lemma}

\begin{proof}
$K\cup L_1$ is closed, so by \vone{(K2)} it suffices to show: if $K\cup L_1\subseteq F\cup G$ with
$F,G$ closed and disjoint, then $K\cup L_1$ lies in $F$ or in $G$. Since $K$ is preconnected, $K$
lies in one of them, say $K\subseteq F$ (the other case is symmetric). We show $L_1\subseteq F$.
Write $L_1=L\cap O_1$ with $O_1$ open, and put $W=L_1\cap G$. Then $W$ is closed, and $W\subseteq L$
is relatively open in $L$: indeed $W=L\cap(O_1\setminus F)$, because a point of $L_1$ lies in $G$
exactly when it does not lie in $F$. If $W$ were nonempty it would meet $K\subseteq F$, which is
impossible as $W\subseteq G$ and $F\cap G=\emptyset$. So $W=\emptyset$, that is, $L_1\subseteq F$.
\end{proof}

\formal{\lean{HPC.T13.isPreconnected_union_clopen}.}

\section{Proper subcontinua of indecomposable continua are thin}

\begin{theorem}\label{thm:thin}
Let $M$ be an indecomposable continuum and $K$ a subcontinuum of $M$ with $K\neq M$. Then no open
set $O\subseteq X$ satisfies $O\cap M\neq\emptyset$ and $O\cap M\subseteq K$.
\end{theorem}

\begin{proof}
Suppose $O$ is such a set. Let $L=\cl{M\setminus K}$. We collect some facts.
\begin{enumerate}[label=(\roman*)]
\item $L\subseteq M$ (as $M$ is closed), $K\cup L=M$, and $L$ is nonempty: there is $x\in
M\setminus K$, and $x\in L$.
\item $L\neq M$. Take $z\in O\cap M$. If $z\in L$, the neighbourhood $O$ of $z$ meets $M\setminus K$,
contradicting $O\cap M\subseteq K$.
\item $\cl{M\setminus L}\subseteq K$: a point of $M\setminus L$ is not in $M\setminus K$, so it is
in $K$; and $K$ is closed.
\item Every nonempty subset $W$ of $L$ that is closed and relatively open in $L$ meets $K$: by
Lemma~\ref{lem:bumping} (applied to $M$, $L$, $W$; $M$ is connected) $W$ meets $\cl{M\setminus
L}\subseteq K$.
\end{enumerate}
\emph{Case 1: $L$ is connected.} Then $L$ is a continuum (closed in the compact set $M$), and
$M=K\cup L$ with $K\ne M$ and $L\neq M$, contradicting indecomposability.

\emph{Case 2: $L$ is not connected.} Since $L$ is nonempty and closed, \vone{(K2)} gives closed sets
$F,G$ with $L=(L\cap F)\cup(L\cap G)$, $L\cap F\cap G=\emptyset$, and both $L\cap F$ and $L\cap G$
nonempty. Then $L\cap F=L\cap(X\setminus G)$ and $L\cap G=L\cap(X\setminus F)$, so both are closed and
relatively open in $L$. Put
\[
A=K\cup(L\cap F),\qquad B=K\cup(L\cap G).
\]
By Lemma~\ref{lem:union-clopen} and~(iv) both are preconnected; they are closed subsets of $M$,
hence compact, and nonempty. So $A$ and $B$ are subcontinua of $M$, and $A\cup B=K\cup L=M$. We show
$A\neq M$; similarly $B\neq M$, and this contradicts indecomposability. Pick $b\in L\cap G$. Then
$b\notin F$, and $X\setminus F$ is an open neighbourhood of $b\in\cl{M\setminus K}$, so it contains a
point $y\in M\setminus K$. This $y$ is not in $K$, and not in $L\cap F$ because $y\notin F$. So
$y\in M\setminus A$.
\end{proof}

\formal{\lean{HPC.T13.false_of_interior}.}

\section{The Baire argument}\label{sec:baire}

\begin{lemma}[Baire category theorem for compact spaces]\label{lem:baire}
Let $Y$ be a nonempty compact metric space and $Y=\bigcup_{k\in\N}F_k$ with every $F_k$ closed.
Then some $F_k$ contains a nonempty open subset of $Y$.
\end{lemma}

\begin{proof}
Suppose not. We choose nonempty open sets $Y\supseteq V_0\supseteq V_1\supseteq\cdots$ with
$\cl{V_{k}}\cap F_k=\emptyset$ and $\cl{V_{k+1}}\subseteq V_k$. Since $F_0\neq Y$, the open set
$Y\setminus F_0$ is nonempty; choose a point $y_0$ in it and a small open ball $V_0$ around $y_0$
whose closure misses the closed set $F_0$. Given $V_k$, the open set $V_k$ is not contained in
$F_{k+1}$, so there is $y_{k+1}\in V_k\setminus F_{k+1}$; choose a small open ball $V_{k+1}$ around
$y_{k+1}$ whose closure lies in $V_k$ and misses $F_{k+1}$. By Cantor's intersection theorem
\vone{(C3)} some point lies in every $\cl{V_k}$; it lies in no $F_k$, contradicting
$Y=\bigcup_kF_k$.
\end{proof}

We also use the \emph{component} of a point $p$ in a set $F$: the union $C$ of all preconnected
subsets of $F$ containing $p$ (and $C=\emptyset$ if $p\notin F$). It is preconnected (a union of
preconnected sets with a common point is preconnected), so it is the largest preconnected subset
of $F$ containing $p$. If $F$ is closed then $C$ is closed: the closure of a preconnected set is
preconnected, and $\cl C\subseteq F$ contains $p$, so $\cl C\subseteq C$ by maximality.

Finally, a compact metric space $Y$ has a \emph{countable basis}: a countable family $\mathcal B$
of open sets such that every open set is a union of members of $\mathcal B$. (Cover $Y$ by
finitely many balls of radius $1/m$ for each $m\ge1$, and take all these balls.)

\begin{proof}[Proof of Theorem~\ref{thm:composant-proper}]
We work in the compact metric space $Y=M$. (A subset of $M$ is a continuum in $M$ if and only if it
is a continuum in $X$, so composants are the same.) Let $\mathcal B$ be a countable basis of $M$,
and let $\mathcal B'$ be the family of nonempty members of $\mathcal B$. For $b\in\mathcal B'$ let
$C_b$ be the component of $p$ in the closed set $M\setminus b$. Each $C_b$ is closed.

Suppose, for a contradiction, that every $y\in M$ lies in a proper subcontinuum $K$ of $M$ containing
$p$. Choose $x\in M\setminus K$. Since $M\setminus K$ is open in $M$, there is $b\in\mathcal B$ with
$x\in b\subseteq M\setminus K$; then $b\in\mathcal B'$, and $K$ is a preconnected subset of
$M\setminus b$ containing $p$, so $K\subseteq C_b$ and $y\in C_b$. Thus
$M=\bigcup_{b\in\mathcal B'}C_b$, a countable union of closed sets. By Lemma~\ref{lem:baire} some
$C_b$ contains a nonempty open subset of $M$. This $C_b$ is nonempty, so it contains $p$; it is
compact (closed in $M$) and connected, so it is a subcontinuum of $M$; and $C_b\neq M$, because
$C_b$ misses the nonempty set $b$. This contradicts Theorem~\ref{thm:thin}.
\end{proof}

\formal{\lean{HPC.T13.exists_not_mem_proper_subcontinua} (in a compact ambient space) and
\lean{HPC.composant_ne_of_indecomposable}.}

\begin{corollary}\label{cor:other-composant}
Let $M$ be a pseudo-arc and $p\in M$. Then there is $q\in M$ with $q\notin\cmp Mp$, and then
$q\neq p$.
\end{corollary}

\begin{proof}
$M$ is an indecomposable continuum (it is hereditarily indecomposable, and a subcontinuum of itself).
Theorem~\ref{thm:composant-proper} gives $q$, and Lemma~\ref{lem:ne-of-composant} gives $q\neq p$,
since a pseudo-arc is nondegenerate.
\end{proof}

\formal{\lean{HPC.Plane.exists_not_mem_composant}. Formally the theorem is applied in the compact
space $M$, and the subcontinuum found there is transported back to $X$.}

\section{A continuum as a space in its own right}\label{sec:subspace}

Let $M\subseteq X$ be a continuum. Part~II of the proof constructs a description \emph{inside the
compact metric space $M$} (with the metric of $X$). This section checks that the hypotheses
transfer. The open sets of the space $M$ are the sets $U\cap M$ with $U$ open in $X$, and diameters
are computed with the same metric.

\begin{lemma}\label{lem:sub-HI}
\begin{enumerate}[label=(\alph*)]
\item The subcontinua of the space $M$ are exactly the subcontinua of $M$ in $X$. Hence if $M$ is
hereditarily indecomposable in $X$, every subcontinuum of the space $M$ is indecomposable.
\item If $p,q\in M$ and $q\notin\cmp Mp$ (composant in $X$), then $q$ is not in the composant of $p$
in the space $M$.
\end{enumerate}
\end{lemma}

\begin{proof}
Compactness and connectedness of a subset $K\subseteq M$ do not depend on whether $K$ is regarded as
a subset of $M$ or of $X$ (the inclusion $M\to X$ is continuous, and so is its inverse on its
image). Both statements follow at once.
\end{proof}

\formal{\lean{HPC.Plane.PA51.hI_sub}, \lean{HPC.Plane.PA51.not_mem_composant_sub}.}

Snake-likeness needs an argument, because the links of an $\varepsilon$-chain covering $M$ need not
all meet $M$, and consecutive links need not meet inside $M$.

\begin{lemma}\label{lem:cover-meeting}
Let $M$ be preconnected and nonempty, and let $D$ be a chain of length $n$ with $M\subseteq\bigcup D$.
\begin{enumerate}[label=(\alph*)]
\item The indices $i$ with $D_i\cap M\neq\emptyset$ form an interval $a\le i\le b$ of integers, and
$(D_a,\dots,D_b)$ is a chain covering $M$ all of whose links meet $M$.
\item If every link of $D$ meets $M$, then $D_i\cap D_{i+1}\cap M\neq\emptyset$ for every $i<n-1$, and
$(D_0\cap M,\dots,D_{n-1}\cap M)$ is a chain in the space $M$.
\end{enumerate}
\end{lemma}

\begin{proof}
(a) Suppose $i<j<k$ with $D_i$ and $D_k$ meeting $M$ but $D_j\cap M=\emptyset$. Let $U$ be the union of
the links with index $<j$ and $V$ the union of the links with index $>j$. These are open, they cover
$M$ (as $D_j$ contributes nothing to $M$), both meet $M$, and they are disjoint, because a link with
index $<j$ and a link with index $>j$ have indices differing by at least $2$. This contradicts the
preconnectedness of $M$. So the set of indices of links meeting $M$ (nonempty, as $M\ne\emptyset$) is an
interval $[a,b]$. A run of consecutive links of a chain is a chain, and it covers $M$ since the other
links do not meet $M$.

(b) Fix $i<n-1$. Let $U$ be the union of the links with index $\le i$ and $V$ the union of those with
index $\ge i+1$. Both are open, they cover $M$, $U$ meets $M$ (in $D_i\cap M$) and $V$ meets $M$ (in
$D_{i+1}\cap M$). Since $M$ is preconnected there is $z\in M\cap U\cap V$: $z\in D_l\cap D_{l'}$ with
$l\le i<i+1\le l'$. As $D_l$ and $D_{l'}$ meet, $l'\le l+1$, which forces $l=i$ and $l'=i+1$. So
$z\in D_i\cap D_{i+1}\cap M$. Now the sets $D_i\cap M$ are open in $M$; if $\abs{i-j}\ge2$ they are
disjoint because $D_i\cap D_j=\emptyset$; if $j=i\pm1$ they meet by what we just proved; and each is
nonempty. So they form a chain in $M$.
\end{proof}

\formal{(a) is \lean{HPC.Plane.exists_cover_meeting}; (b) is \lean{HPC.Plane.PA51.inter_inter_nonempty}
and \lean{HPC.Plane.PA51.subChain}.}

\begin{corollary}\label{cor:sub-snake}
If $M\subseteq X$ is a snake-like continuum, then for every $\varepsilon>0$ the space $M$ is covered
by an $\varepsilon$-chain of the space $M$. Consequently, if $M$ is a pseudo-arc in $X$, then the
space $M$ is a pseudo-arc in itself, and points $p,q\in M$ with $q\notin\cmp Mp$ lie in different
composants of the space $M$.
\end{corollary}

\begin{proof}
Take an $\varepsilon$-chain covering $M$, cut it down by Lemma~\ref{lem:cover-meeting}(a), and restrict
it by Lemma~\ref{lem:cover-meeting}(b). Diameters do not increase on intersecting with $M$. The rest
is Lemma~\ref{lem:sub-HI}.
\end{proof}

\formal{\lean{HPC.Plane.PA51.snake_sub}.}

\section*{Exercises}
\addcontentsline{toc}{section}{Exercises}

\begin{exercise}
Show that two composants of an indecomposable continuum are disjoint or equal. (Hint: if $K\ni p$ and
$K'\ni q$ are proper subcontinua with a common point, then $K\cup K'$ is a continuum, and it is not
$M$ by indecomposability.)
\end{exercise}

\begin{exercise}
Show that every composant of an indecomposable continuum $M$ is dense in $M$. (Hint: use
Theorem~\ref{thm:thin} on the component of $p$ in the complement of a small open set.)
\end{exercise}

\begin{exercise}
Show that if a continuum $M$ is decomposable, $M=A\cup B$ with $A,B$ proper subcontinua, then
$A\cap B\neq\emptyset$ and $\cmp Mp=M$ for every $p\in A\cap B$.
\end{exercise}

\begin{exercise}
Show by an example that Lemma~\ref{lem:cover-meeting}(b) fails if $M$ is not assumed preconnected.
\end{exercise}

%======================================================================
\chapter{Consolidations and other tools for chains}\label{ch:consolidations}
%======================================================================

This chapter collects the operations on chains used in the rest of the volume. The most
important is \emph{consolidation}: grouping consecutive links of a fine chain into blocks, each
block becoming one link of a coarser chain. Pattern following (Chapter~\ref{ch:patterns}) and
the folding theorem (Chapter~\ref{ch:folding}) both produce their chains in this way, and
Bing's Theorem~1 (Section~\ref{sec:bing-thm1}) says that crookedness survives consolidation.
Throughout, $X$ is a metric space, and chains are as in \vone{Definition~2.1}.

\section{Consolidations and block maps}

\begin{definition}\label{def:consolidation}
A chain $D$ is a \emph{consolidation} of a chain $E$ if $E$ refines $D$ and every link of $D$ is a
union of links of $E$.
\end{definition}

\formal{\lean{HPC.Chain.IsConsolidationOf}.}

The consolidations we construct all come from grouping consecutive links.

\begin{definition}\label{def:block-map}
Let $E$ be a chain of length $m$. A \emph{block map} is a surjective function $\beta\colon[m]\to[p]$
with unit steps: $\abs{\beta(t+1)-\beta(t)}\le1$ for $t+1<m$. The \emph{block chain} $E^\beta$ has
the $p$ links
\[
E^\beta_u=\bigcup\set{E_t : \beta(t)=u}\qquad(u\in[p]).
\]
The set $\set{t:\beta(t)=u}$ is the \emph{block} $u$.
\end{definition}

A block need not consist of consecutive indices: $\beta$ may return to a value it has left. All
that is required is that $\beta$ moves by at most one at each step.

\begin{lemma}\label{lem:block-chain}
If $\beta$ is a block map for $E$, then $E^\beta$ is a chain of length $p$, and it is a
consolidation of $E$. Moreover $E_t\subseteq E^\beta_{\beta(t)}$ for every $t$.
\end{lemma}

\begin{proof}
Each $E^\beta_u$ is open, as a union of open sets, and the last statement is clear; it shows that $E$
refines $E^\beta$, and by definition each link of $E^\beta$ is a union of links of $E$. It remains to
check the chain condition $E^\beta_u\cap E^\beta_{u'}\ne\emptyset\iff\abs{u-u'}\le1$.

Suppose $z\in E^\beta_u\cap E^\beta_{u'}$, say $z\in E_t\cap E_{t'}$ with $\beta(t)=u$, $\beta(t')=u'$.
As $E_t$ and $E_{t'}$ meet, $\abs{t-t'}\le1$, and unit steps give $\abs{u-u'}\le1$.

Conversely let $\abs{u-u'}\le 1$. If $u=u'$, the link $E^\beta_u$ is nonempty because some $t$ has
$\beta(t)=u$ (surjectivity) and links of a chain are nonempty. If $u'=u+1$ (the case $u=u'+1$ is
symmetric), choose $t_1,t_2$ with $\beta(t_1)=u$, $\beta(t_2)=u+1$, and walk from $t_1$ to $t_2$. Along
the walk $\beta$ changes by at most one at each step and goes from $u$ to $u+1$. Let $t$ be the last
index of the walk at which $\beta$ equals $u$, and $t'$ the next index of the walk. Then $\beta(t')\ne u$;
and $\beta(t')\neq u-1$, since otherwise $\beta$ would have to pass through $u$ again later. So
$\beta(t')=u+1$, and since $\abs{t-t'}=1$, the links $E_t\subseteq E^\beta_u$ and $E_{t'}\subseteq
E^\beta_{u+1}$ meet.
\end{proof}

\formal{\lean{HPC.Chain.IsBlockMap}, \lean{HPC.Chain.ofBlocks},
\lean{HPC.Chain.blockLink_adjacent}, \lean{HPC.Chain.ofBlocks_isConsolidationOf}.}

\begin{lemma}\label{lem:block-from}
Let $\beta\colon[m]\to[p]$ be a block map for $E$, and $R,S\in X$. Then $E^\beta$ goes from $R$ to $S$
if and only if $R$ lies in some link of $E$ and every link $E_t$ containing $R$ has $\beta(t)=0$, and $S$
lies in some link of $E$ and every link $E_t$ containing $S$ has $\beta(t)=p-1$.
\end{lemma}

\begin{proof}
$R\in E^\beta_u$ means that $R\in E_t$ for some $t$ with $\beta(t)=u$. So ``$R$ lies in $E^\beta_0$ and in no
other link of $E^\beta$'' says exactly that $R$ lies in some $E_t$, and that every $E_t$ containing $R$
has $\beta(t)=0$. The same holds for $S$ and the last link.
\end{proof}

\formal{\lean{HPC.Fold.Q78.isFrom_ofBlocks_iff}.}

Often it is more natural to produce a function into $\N$ than onto some $[p]$.

\begin{definition}\label{def:schedule}
A \emph{schedule} for a chain of length $m$ is a function $\beta\colon[m]\to\N$ with unit steps
which takes the value $0$.
\end{definition}

\begin{lemma}\label{lem:schedule}
A schedule $\beta$ is a block map onto $[p]$, where $p=1+\max_t\beta(t)$.
\end{lemma}

\begin{proof}
Its values lie in $[p]$. Every $v$ with $0\le v\le p-1$ lies between the value $0$ and the maximal value,
so it is attained, by the discrete intermediate value theorem \vone{Lemma~2.5}.
\end{proof}

\formal{\lean{HPC.IsSchedule}, \lean{HPC.IsSchedule.exists_blockMap}.}

\begin{lemma}\label{lem:consolidate-length}
Let $D$ be a chain of length $n$ from $P$ to $Q$, and $2\le m\le n$. Then $D$ has a consolidation of
length $m$ which goes from $P$ to $Q$.
\end{lemma}

\begin{proof}
Use the block map $\beta(i)=\min(i,m-1)$: the blocks are $\set0,\set1,\dots,\set{m-2}$ and
$\set{m-1,\dots,n-1}$. It is surjective with unit steps, so $D^\beta$ is a consolidation by
Lemma~\ref{lem:block-chain}. The only link of $D$ containing $P$ is $D_0$, and $\beta(0)=0$; the only link
containing $Q$ is $D_{n-1}$, and $\beta(n-1)=m-1$. Lemma~\ref{lem:block-from} shows that $D^\beta$ goes from
$P$ to $Q$.
\end{proof}

\formal{\lean{HPC.Chain.exists_consolidation_isFrom}.}

\section{Crookedness passes to consolidations}\label{sec:bing-thm1}

\vone{Lemma~2.12} says that a refinement of a chain crooked in $D$ is again crooked in $D$. The next
theorem says the same for the coarse chain: if $F$ is crooked in $E$, it is crooked in every
consolidation of $E$.

\begin{lemma}\label{lem:index-dist}
Let $E$ refine $D$, and let $E_a\subseteq D_h$, $E_b\subseteq D_k$. If $\abs{a-b}\le2$, then
$\abs{h-k}\le2$.
\end{lemma}

\begin{proof}
If $\abs{a-b}\le1$, then $E_a$ and $E_b$ meet, so $D_h$ and $D_k$ meet and $\abs{h-k}\le1$. If
$\abs{a-b}=2$, let $c$ be the index between them and $E_c\subseteq D_l$. Then $E_c$ meets $E_a$ and $E_b$,
so $D_l$ meets $D_h$ and $D_k$, and $\abs{h-k}\le\abs{h-l}+\abs{l-k}\le2$.
\end{proof}

\formal{\lean{HPC.Chain.index_dist_le_of_subset}.}

\begin{theorem}[Bing 1948, Theorem~1]\label{thm:bing1}
Let $D$ be a consolidation of $E$, and let $F$ be crooked in $E$. Then $F$ is crooked in $D$.
\end{theorem}

\begin{proof}
$F$ refines $E$, which refines $D$, so $F$ refines $D$. Let $k>h+2$, and let $F_i$ meet $D_h$ and $F_j$ meet
$D_k$.

\emph{Step 1: links of $E$ in between.} Let $z\in F_i\cap D_h$. Since $D_h$ is a union of links of $E$,
there is a link $E_a\subseteq D_h$ containing $z$; so $F_i$ meets $E_a$. Similarly there is $E_b\subseteq D_k$
which meets $F_j$. By Lemma~\ref{lem:index-dist}, $\abs{a-b}>2$, since $\abs{h-k}>2$.

\emph{Step 2: crookedness in $E$.} Suppose $a<b$, so $b>a+2$. Crookedness of $F$ in $E$ gives $r,s$ with
$\mathrm{Btw}(i,r,s,j)$, $F_r\subseteq E_{b-1}$ and $F_s\subseteq E_{a+1}$. If $b<a$, crookedness applied to
the pair $(b,a)$ and $(j,i)$ gives $s,r$ with $\mathrm{Btw}(j,s,r,i)$, $F_s\subseteq E_{a-1}$, $F_r\subseteq
E_{b+1}$; and $\mathrm{Btw}(j,s,r,i)$ is the same as $\mathrm{Btw}(i,r,s,j)$. In both cases we have
$\mathrm{Btw}(i,r,s,j)$, $F_r\subseteq E_{b'}$ and $F_s\subseteq E_{a'}$ with $\abs{b'-b}\le1$ and
$\abs{a'-a}\le1$.

\emph{Step 3: back to $D$.} Let $E_{b'}\subseteq D_{k'}$ and $E_{a'}\subseteq D_{h'}$. Since $E_{b'}$ meets
$E_b\subseteq D_k$, we get $\abs{k'-k}\le1$, so $k'\ge k-1$; similarly $h'\le h+1$. Choose a refinement map
$\chi$ for $F$ in $D$ with $\chi(r)=k'$ and $\chi(s)=h'$ (possible, as $F_r\subseteq D_{k'}$ and $F_s\subseteq
D_{h'}$, and $r\neq s$). By \vone{Lemma~2.4} $\chi$ has unit steps, and $\chi(r)\ge k-1>h+1\ge\chi(s)$.

\emph{Step 4: interpolation.} Suppose $r<s$ (if $r>s$, run the same argument backwards). Let $y$ be the
first index in $[r,s]$ with $\chi(y)\le h+1$. Then $y>r$, so $\chi(y-1)>h+1$, and unit steps force
$\chi(y)=h+1$. Next let $x$ be the last index in $[r,y]$ with $\chi(x)\ge k-1$. Then $x<y$ (as $\chi(y)=h+1<k-1$),
so $\chi(x+1)<k-1$, and unit steps force $\chi(x)=k-1$. Thus $r\le x<y\le s$, $F_x\subseteq D_{k-1}$ and
$F_y\subseteq D_{h+1}$. Since $\mathrm{Btw}(i,r,s,j)$ and $x,y\in[r,s]$ in the same order as $r,s$, also
$\mathrm{Btw}(i,x,y,j)$.
\end{proof}

\formal{\lean{HPC.Chain.isCrookedIn'_of_isConsolidationOf}; the interpolation of Step~4 is
\lean{HPC.Chain.exists_value_between}.}

\section{Reversal and least indices}

\begin{definition}
The \emph{reverse} of a chain $D$ of length $n$ is the chain $\overline D$ with
$\overline D_i=D_{n-1-i}$.
\end{definition}

\begin{lemma}\label{lem:reverse}
If $E$ is crooked in $D$, then $\overline E$ is crooked in $D$ and $E$ is crooked in $\overline D$.
\end{lemma}

\begin{proof}
Reversing $E$ reverses the order of its indices, and the condition $\mathrm{Btw}(i,r,s,j)$ (``$r,s$ strictly
between $i$ and $j$, with $r$ nearer to $i$'') is invariant under reversing the order. For the second
statement, let $k>h+2$, let $E_i$ meet $\overline D_h=D_{n-1-h}$ and $E_j$ meet $\overline D_k=D_{n-1-k}$.
Put $H=n-1-k$ and $K=n-1-h$, so $K>H+2$, $E_j$ meets $D_H$ and $E_i$ meets $D_K$. Crookedness of $E$ in $D$
gives $r,s$ with $\mathrm{Btw}(j,r,s,i)$, $E_r\subseteq D_{K-1}=\overline D_{h+1}$ and $E_s\subseteq
D_{H+1}=\overline D_{k-1}$. Then $\mathrm{Btw}(i,s,r,j)$, and $s,r$ are the required links.
\end{proof}

\formal{\lean{HPC.Chain.reverse}, \lean{HPC.Chain.IsCrookedIn'.reverse_left},
\lean{HPC.Chain.IsCrookedIn'.reverse_right}.}

Reversal lets us state many arguments ``from the bottom'' and obtain the version ``from the top'' for free.

\begin{definition}\label{def:minidx}
Let $E$ (of length $m$) refine $D$ (of length $n$). For $t\in[m]$ let $\lo(t)$ be the least and $\hi(t)$ the
greatest index $j$ with $E_t\subseteq D_j$.
\end{definition}

\begin{lemma}\label{lem:minidx}
\begin{enumerate}[label=(\alph*)]
\item $\lo(t)\le\hi(t)\le\lo(t)+1$, and $E_t\subseteq D_j$ holds exactly for $j\in\set{\lo(t),\hi(t)}$.
\item $\lo$ and $\hi$ have unit steps.
\item If $E_t$ meets $D_j$, then $\abs{\lo(t)-j}\le1$.
\end{enumerate}
\end{lemma}

\begin{proof}
(a) If $E_t\subseteq D_j\cap D_{j'}$, then $D_j$ and $D_{j'}$ meet (links are nonempty), so
$\abs{j-j'}\le1$. (b) is \vone{Lemma~2.4} for the refinement maps $\lo$ and $\hi$. (c) $E_t\subseteq
D_{\lo(t)}$ meets $D_j$.
\end{proof}

\formal{\lean{HPC.Chain.Refines.minIdx}, \lean{HPC.Chain.Refines.maxIdx},
\lean{HPC.Chain.subset_adjacent_of_subset}, \lean{HPC.Chain.minIdx_lip},
\lean{HPC.Chain.maxIdx_lip}, \lean{HPC.Chain.minIdx_near_of_inter}.}

\section{Spans and subchains}

\begin{definition}\label{def:span}
For a chain $D$ and indices $a\le b$, the \emph{span} $\spn Dab=D_a\cup D_{a+1}\cup\dots\cup D_b$. The
chain $(D_a,\dots,D_b)$ is the \emph{subchain} from $a$ to $b$.
\end{definition}

\begin{lemma}\label{lem:span-chained}
If the links of $D$ are bounded with diameter $\le\varepsilon$, then every span $\spn Dab$ is
$\varepsilon$-chained (\vone{Definition~1.5}).
\end{lemma}

\begin{proof}
The subchain $(D_a,\dots,D_b)$ is a chain with union $\spn Dab$; apply \vone{Lemma~2.10}.
\end{proof}

\formal{\lean{HPC.Chain.span}, \lean{HPC.Chain.chainedAt_span}.}

\begin{definition}\label{def:subchain-from}
For points $R,S$ and indices $a\le b$, we say that \emph{the subchain from $a$ to $b$ joins $R$ and $S$} if
it goes from $R$ to $S$ or from $S$ to $R$; that is, one of the two points lies in $D_a$ and in no other
$D_t$ with $a\le t\le b$, and the other lies in $D_b$ and in no other such link.
\end{definition}

\formal{\lean{HPC.Chain.IsSubchainFrom}.}

\begin{lemma}\label{lem:subchain-exists}
Let $D$ be a chain and $R,S\in\bigcup D$ such that no link of $D$ contains both $R$ and $S$. Then there are
$a<b$ such that the subchain from $a$ to $b$ joins $R$ and $S$.
\end{lemma}

\begin{proof}
Let $I_R$ (resp.\ $I_S$) be the nonempty set of indices of links containing $R$ (resp.\ $S$). Two links
containing the same point meet, so each of $I_R$, $I_S$ consists of one index or two consecutive indices,
and $I_R\cap I_S=\emptyset$ by hypothesis. Hence either every element of $I_R$ is smaller than every element
of $I_S$, or the reverse. (If $r_1<s_1$ and $s_2<r_2$ with $r_1,r_2\in I_R$, $s_1,s_2\in I_S$, then
$\abs{r_1-r_2}\le1$ and $\abs{s_1-s_2}\le1$ leave no room for four distinct indices in this order.) In the
first case let $a=\max I_R$ and $b=\min I_S$. Then $a<b$; $R\in D_a$, and $R$ is in no $D_t$ with $a<t\le b$
because $a$ is the largest index in $I_R$; $S\in D_b$, and $S$ is in no $D_t$ with $a\le t<b$ because $b$ is
the smallest index in $I_S$. The second case is symmetric.
\end{proof}

\formal{\lean{HPC.Chain.exists_isSubchainFrom}.}

\section*{Exercises}
\addcontentsline{toc}{section}{Exercises}

\begin{exercise}
Show that the converse of Lemma~\ref{lem:block-chain} fails in the following sense: a consolidation in
the sense of Definition~\ref{def:consolidation} need not be a block chain $E^\beta$. (Hint: a link of $E$
may be used in two links of $D$.)
\end{exercise}

\begin{exercise}
Let $\beta$ be a block map for $E$, and let $F$ be crooked in $E$. Show directly from
Theorem~\ref{thm:bing1} that $F$ is crooked in $E^\beta$. Then show that the conclusion can fail if $\beta$ is
only assumed surjective (without unit steps): $E^\beta$ need not even be a chain.
\end{exercise}

\begin{exercise}
Prove the converse of Lemma~\ref{lem:index-dist} fails: give $E$ refining $D$ with $\abs{h-k}\le2$ but
$\abs{a-b}$ arbitrarily large.
\end{exercise}

\begin{exercise}
In Lemma~\ref{lem:subchain-exists}, show that the hypothesis ``no link contains both $R$ and $S$'' holds as
soon as the links have diameter less than $d(R,S)$.
\end{exercise}

%======================================================================
\chapter{Pattern following}\label{ch:patterns}
%======================================================================

Let $D^{(0)},D^{(1)},\dots$ be a description from $P$ to $Q$ (\vone{Definition~3.1}). A
\emph{pattern} in a chain $D$ of length $n$ is a function $x\colon[N]\to[n]$ with unit steps,
$x(0)=0$ and $x(N-1)=n-1$: a walk through the links of $D$ from the first to the last. A chain $E$ of
length $N$ \emph{follows} $x$ in $D$ if $E_u\subseteq D_{x(u)}$ for all $u$ (\vone{Definition~2.3}).
This chapter proves:

\begin{theorem}[Pattern following; Bing 1948, Theorem~6]\label{thm:pattern-following}
Let $D^{(0)},D^{(1)},\dots$ be a description from $P$ to $Q$ in a compact metric space $X$, with
$P\neq Q$. Let $i\in\N$ and let $x\colon[N]\to[n_i]$ be a pattern in $D^{(i)}$. Then there are $j>i$
and a consolidation $E$ of $D^{(j)}$ of length $N$ which goes from $P$ to $Q$ and follows $x$ in $D^{(i)}$.
\end{theorem}

\formal{\lean{HPC.Description.patternFollowing}; the property itself is
\lean{HPC.Description.PatternFollowing}, and patterns are \lean{HPC.Description.IsPattern}.}

\begin{remark}\label{rem:pf-false}
The hypothesis $P\neq Q$ cannot be dropped. Let $X$ be a one-point space $\set y$ and let every level be
the chain with one link $X$; this is a description from $y$ to $y$ (crookedness holds vacuously, as
there are no indices $h,k$ with $k>h+2$). The pattern $x=(0,0)$ of length $2$ cannot be followed by any
chain from $y$ to $y$: in a chain of length $2$ from $y$ to $y$, the point $y$ would lie in the first link
and in the last link, but it may lie in only one link. Bing's paper does not have this hypothesis; it is
harmless for us, since the uniqueness theorem assumes $P\neq Q$ anyway.
\end{remark}

\formal{\lean{HPC.Description.unitDescription}, \lean{HPC.Description.not_patternFollowing_unconditional}.}

The proof has a combinatorial core (Sections~\ref{sec:pf-setting}--\ref{sec:pf-induction}), a statement
about integer sequences proved by induction, and a geometric part (Section~\ref{sec:pf-geometry}) which
applies the core to the indices of links. The idea is this. Let $F=D^{(m+1)}$ be a very deep level and
$R=D^{(m)}$ the level above it. Reading off, for each link of $F$, a link of $R$ containing it gives an
integer sequence $\sigma$ which is crooked, because $F$ is crooked in $R$. Since $m$ is large, the links of
$R$ are tiny compared with those of $D^{(0)}$, and a whole stretch of $w$ consecutive links of $R$ lies in a
single link of $D^{(0)}$. We must group the links of $F$ into $N$ blocks, the $u$-th block lying in
$D^{(0)}_{x(u)}$. Crookedness of $\sigma$ is what allows the blocks to go back and forth as the pattern
demands.

\section{The combinatorial setting}\label{sec:pf-setting}

Throughout Sections~\ref{sec:pf-setting}--\ref{sec:pf-induction} we fix the following data.

\begin{itemize}
\item A \emph{fine reading} $\sigma\colon\N\to\Z$ and a time $T\in\N$ such that
\begin{enumerate}[label=(S\arabic*)]
\item\label{S:lip} $\abs{\sigma(f+1)-\sigma(f)}\le1$ for all $f$;
\item\label{S:crook} \emph{crookedness:} if $f<f'\le T$ and $\abs{\sigma(f)-\sigma(f')}\ge3$, then there are
$r,s$ with $f<r<s<f'$, $\abs{\sigma(r)-\sigma(f')}\le2$ and $\abs{\sigma(s)-\sigma(f)}\le2$.
\end{enumerate}
We call $f\in\N$ a \emph{time} and the values of $\sigma$ \emph{positions}.
\item A width $w\ge1$ and a \emph{coarse reading} $\tau\colon\Z\to\Z$ which is \emph{$w$-slow}:
\[
\abs{p-p'}\le w\implies\abs{\tau(p)-\tau(p')}\le1 .
\]
\end{itemize}

The \emph{canonical reading} at time $f$ is $\ell(f)=\tau(\sigma(f))$. Since $w\ge1$, $\ell$ has unit steps.

\begin{definition}\label{def:allowed}
Let $d\in\N$. A value $v\in\Z$ is \emph{allowed at time $f$ with radius $d$} if there is a position $p$
with $\abs{p-\sigma(f)}\le d$ and $\tau(p)=v$. We write $\mathrm{Al}_d(f,v)$.
\end{definition}

\begin{lemma}\label{lem:allowed}
\begin{enumerate}[label=(\alph*)]
\item $\mathrm{Al}_0(f,\ell(f))$; and $\mathrm{Al}_d(f,v)$ implies $\mathrm{Al}_{d'}(f,v)$ for $d'\ge d$.
\item $\abs{\sigma(f')-\sigma(f)}\le\abs{f'-f}$; hence $\mathrm{Al}_d(f,v)$ implies $\mathrm{Al}_{d+k}(f',v)$ when
$\abs{f'-f}\le k$.
\item If $d\le w$ and $\mathrm{Al}_d(f,v)$, then $\abs{\ell(f)-v}\le1$.
\item If $\mathrm{Al}_d(f,v)$, $\mathrm{Al}_{d'}(f',v')$ and $\abs{v-v'}\ge2$, then
$\abs{\sigma(f)-\sigma(f')}\ge w+1-d-d'$, and so $\abs{f-f'}\ge w+1-d-d'$.
\end{enumerate}
\end{lemma}

\begin{proof}
(a) Take $p=\sigma(f)$. (b) The first statement follows from~\ref{S:lip} by the triangle inequality; then a
position within $d$ of $\sigma(f)$ is within $d+k$ of $\sigma(f')$. (c) $\abs{p-\sigma(f)}\le w$ and
slowness give $\abs{\tau(p)-\tau(\sigma(f))}\le1$. (d) Let $\tau(p)=v$, $\tau(p')=v'$ with
$\abs{p-\sigma(f)}\le d$, $\abs{p'-\sigma(f')}\le d'$. Slowness and $\abs{v-v'}\ge2$ force $\abs{p-p'}\ge w+1$,
and the triangle inequality gives the first claim; the second follows from~(b).
\end{proof}

\formal{\lean{HPC.PatCore.Setting}, \lean{HPC.PatCore.Slow}, \lean{HPC.PatCore.Al},
\lean{HPC.PatCore.al_canon}, \lean{HPC.PatCore.al_move}, \lean{HPC.PatCore.al_near},
\lean{HPC.PatCore.al_far}, \lean{HPC.PatCore.time_ge}.}

We also use the discrete intermediate value theorem \vone{Lemma~2.5} for $\tau$ along positions: since $\tau$ has
unit steps on consecutive positions, if $u$ lies between $\tau(p)$ and $\tau(p')$ then $\tau(p'')=u$ for some
$p''$ between $p$ and $p'$; in particular $\abs{p''-p}\le\abs{p'-p}$.

\formal{\lean{HPC.PatCore.ivt_between}.}

\section{Lifts}

Let $x\colon\N\to\Z$ (a pattern; only finitely many values matter). For index bounds $p\le q$, times
$a\le b$ and indices $i,j$:

\begin{definition}\label{def:lift}
A \emph{lift of $x$ over $[a,b]$ in $[p,q]$ from $i$ to $j$} is a function $\varphi\colon[a,b]\to\N$ with
$\varphi(a)=i$, $\varphi(b)=j$, unit steps, and, for every $f\in[a,b]$,
\[
p\le\varphi(f)\le q\qquad\text{and}\qquad \mathrm{Al}_w\bigl(f,x(\varphi(f))\bigr).
\]
\end{definition}

So a lift walks through the indices of the pattern, one step at a time at most, in such a way that at each
time $f$ the pattern value $x(\varphi(f))$ is read by $\tau$ at a position within $w$ of $\sigma(f)$.

\formal{\lean{HPC.PatCore.Real}.}

\begin{lemma}[Operations on lifts]\label{lem:lift-ops}
\begin{enumerate}[label=(\alph*)]
\item \emph{Gluing.} A lift over $[a,m]$ from $i$ to $j$ and a lift over $[m,b]$ from $j$ to $k$ (in the same
$[p,q]$) combine into a lift over $[a,b]$ from $i$ to $k$.
\item \emph{Widening.} A lift in $[p',q']$ is a lift in $[p,q]$ when $p\le p'$ and $q'\le q$.
\item \emph{Negation.} A lift of $-x$ for the coarse reading $-\tau$ is a lift of $x$ for $\tau$ (and $-\tau$ is
$w$-slow when $\tau$ is).
\item \emph{Reflection.} If $\varphi$ is a lift of $k\mapsto x(p+q-k)$ in $[p,q]$ from $i$ to $j$, then
$p+q-\varphi$ is a lift of $x$ in $[p,q]$ from $p+q-i$ to $p+q-j$.
\item \emph{Step.} If $\mathrm{Al}_w(a,x(i))$ and $\mathrm{Al}_w(a+1,x(i+1))$, with $p\le i<q$, then there is a lift
over $[a,a+1]$ from $i$ to $i+1$.
\item \emph{Ramp.} If $q-p\le b-a$ and $\mathrm{Al}_w(f,x(i))$ for all $f\in[a,b]$ and $i\in[p,q]$, then
$\varphi(f)=\min(p+(f-a),q)$ is a lift over $[a,b]$ from $p$ to $q$.
\end{enumerate}
\end{lemma}

\begin{proof}
All are immediate from the definition. In (a) use the first lift on $[a,m]$ and the second on $[m,b]$; they agree at
$m$. In (c), $\tau(p)=v$ if and only if $-\tau(p)=-v$.
\end{proof}

\formal{\lean{HPC.PatCore.real_glue}, \lean{HPC.PatCore.real_widen}, \lean{HPC.PatCore.real_of_neg},
\lean{HPC.PatCore.real_reflect}, \lean{HPC.PatCore.real_step}, \lean{HPC.PatCore.real_ramp}.}

\section{The induction}\label{sec:pf-induction}

Fix a \emph{budget} $B\in\N$ with $2B+10\le w$.

\begin{definition}\label{def:good}
A tuple $(\tau,x,p,q,a,b,e)$ (with $\tau$ $w$-slow) is \emph{good} if
\begin{enumerate}[label=(G\arabic*)]
\item\label{G:order} $p\le q$, $a\le b\le T$, and $q-p\le b-a$;
\item\label{G:x} $x$ has unit steps on $[p,q]$, and $x(p)\le x(i)\le x(q)$ for $p\le i\le q$;
\item\label{G:read} $x(p)\le\ell(f)\le x(q)$ for $a\le f\le b$;
\item\label{G:anchor} $\mathrm{Al}_e(a,x(p))$ and $\mathrm{Al}_e(b,x(q))$ (the \emph{anchors});
\item\label{G:budget} $e+2(q-p)\le B$.
\end{enumerate}
\end{definition}

So the pattern goes from its minimum to its maximum, the canonical readings stay in that range, and the end values
are allowed at the end times with radius $e$. Note $e\le B\le w$, so by Lemma~\ref{lem:allowed}(c),
$\ell(a)\le x(p)+1$ and $\ell(b)\ge x(q)-1$.

\formal{\lean{HPC.PatCore.Good}.}

\begin{theorem}[Main induction]\label{thm:pf-main}
If $(\tau,x,p,q,a,b,e)$ is good, there is a lift of $x$ over $[a,b]$ in $[p,q]$ from $p$ to $q$.
\end{theorem}

\begin{proof}
By strong induction on $b-a$: assume the statement for all good tuples with a shorter time interval. Write
$\mu=x(p)$ and $M=x(q)$. We distinguish cases; they are exhaustive, as explained at the end.

\paragraph{Case 1: $M=\mu$ (constant pattern).} By~\ref{G:x} and~\ref{G:read}, $x(i)=\mu=\ell(f)$ for all $i\in[p,q]$
and $f\in[a,b]$, so every $x(i)$ is allowed at every $f$ with radius $0$. Use the ramp.

\paragraph{Case 2: $M\neq\mu$ and $x(p+1)=\mu$ (a repeated first value).} Then $p<q$ and $a<b$. By
Lemma~\ref{lem:allowed}(b), $\mathrm{Al}_{e+1}(a+1,\mu)$. The tuple $(\tau,x,p+1,q,a+1,b,e+1)$ is good: the pattern on
$[p+1,q]$ still starts at its minimum $x(p+1)=\mu$, and the budget is $e+1+2(q-p-1)<e+2(q-p)$. By induction there is a
lift over $[a+1,b]$ from $p+1$ to $q$. In front of it put the step from $p$ at time $a$ to $p+1$ at time $a+1$ (allowed:
$x(p)=\mu$ with radius $e$, $x(p+1)=\mu$ with radius $e+1$, both $\le w$), and glue.

\paragraph{Case 3: $M\neq\mu$ and $x(q-1)=M$ (a repeated last value).} Symmetric to Case~2, removing the last time and
the last index.

\paragraph{Case 4: $M=\mu+1$, $x(p+1)\neq\mu$, $x(q-1)\ne M$.} Then every $x(i)$ ($p\le i\le q$) and every $\ell(f)$
($a\le f\le b$) is $\mu$ or $\mu+1$; and $x(p+1)=\mu+1$, $x(q-1)=\mu$. Let $W=q-p$.

\emph{Short intervals.} If $b-a+e\le w$, every value $\mu$, $\mu+1$ is allowed at every time of $[a,b]$ with radius $w$
(move the anchors by Lemma~\ref{lem:allowed}(b)); use the ramp.

\emph{Long intervals.} Otherwise, there is a time $f_0\in[a,b]$ at which both $\mu$ and $\mu+1$ are allowed with radius
$e+1$. Indeed, if $\ell(a)=\mu+1$ take $f_0=a$; if $\ell(b)=\mu$ take $f_0=b$. Otherwise $\ell(a)=\mu$ and $\ell(b)=\mu+1$;
let $f_1$ be the first time with $\ell(f_1)=\mu+1$; then $f_1>a$, and $f_0=f_1-1$ has $\ell(f_0)=\mu$ while $\mu+1$ is
allowed at $f_0$ with radius $1$. Now let $u=\max\bigl(a+1,\min(f_0,b-1-W)\bigr)$. Then $a+1\le u$, $u+W+1\le b$ (the
interval is long), and every time $f\in[u,u+W]$ satisfies $\abs{f-f_0}\le W+1$; so both values are allowed at every such
$f$ with radius $e+W+2\le w$. Let $\theta(f)=1$ if $\ell(f)=\mu+1$ and $\theta(f)=0$ otherwise, and define
\[
\varphi(f)=\begin{cases}
p & f=a,\\
p+\theta(f) & a<f<u,\\
\min(p+1+(f-u),\,q-1) & u\le f<u+W,\\
q-1+\theta(f) & u+W\le f<b,\\
q & f=b.
\end{cases}
\]
One checks directly that $\varphi$ has unit steps. It is a lift: at $f=a$ and $f=b$ use the anchors; for $a<f<u$, the
value $x(p+\theta(f))$ is $\mu+1=\ell(f)$ if $\theta(f)=1$ and $\mu=\ell(f)$ otherwise, so it is allowed with radius $0$;
similarly for $u+W\le f<b$, using $x(q-1)=\mu$, $x(q)=\mu+1$; and for $u\le f<u+W$ both values are allowed.

\paragraph{Case 5: $M\ge\mu+2$, $x(p+1)=\mu+1$, and the pattern returns to its minimum:} $x(i)=\mu$ for some $i\in(p,q]$.
Let $i_r$ be the least such $i$; then $i_r\ge p+2$. Let $i_c\in[p,i_r]$ be an index where $x$ is largest on $[p,i_r]$
and $c=x(i_c)$; then $c\ge\mu+1$ and $p<i_c<i_r$.

Let $f_s$ be the first time $f\ge a$ with $\ell(f)\ge c+1$, or $f_s=b$ if there is none. Then $\ell(f)\le c$ for
$a\le f<f_s$, and $f_s>a$ (as $\ell(a)\le\mu+1\le c$, and $b>a$ since $q-p\ge2$). There is a value $v\ge\max(c,\mu+2)$ with
$\mathrm{Al}_e(f_s,v)$: if $\ell(f_s)\ge c+1$ then $\ell(f_s)=c+1$ (unit steps) and we take $v=c+1$; otherwise $f_s=b$ and we
take $v=M\ge c$.

\emph{Crookedness.} By Lemma~\ref{lem:allowed}(d) applied to $\mathrm{Al}_e(a,\mu)$ and $\mathrm{Al}_e(f_s,v)$,
$\abs{\sigma(a)-\sigma(f_s)}\ge w+1-2e\ge3$. By~\ref{S:crook} there are times $a<r<s<f_s$ with
$\abs{\sigma(r)-\sigma(f_s)}\le2$ and $\abs{\sigma(s)-\sigma(a)}\le2$. Then:
\begin{itemize}
\item $\mathrm{Al}_{e+2}(r,c)$: a position $p_v$ with $\tau(p_v)=v\ge c$ lies within $e+2$ of $\sigma(r)$, and
$\ell(r)=\tau(\sigma(r))\le c$; by the intermediate value theorem for $\tau$, the value $c$ is read at a position between
$\sigma(r)$ and $p_v$.
\item $\mathrm{Al}_{e+2}(s,\mu)$: move the anchor at $a$.
\item The three pieces are long: by Lemma~\ref{lem:allowed}(d) and the triangle inequality, $\abs{\sigma(a)-\sigma(r)}$,
$\abs{\sigma(r)-\sigma(s)}$ and $\abs{\sigma(s)-\sigma(b)}$ are all at least $w+1-2e-4\ge q-p$ (using $w\ge2B+10\ge
2e+4(q-p)+10$), and time differences dominate position differences, so $r-a$, $s-r$, $b-s\ge q-p$.
\end{itemize}

\emph{Three subproblems.} The following tuples are good, each with radius $e+2$ and budget at most $e+2+2(q-p-1)\le B$:
\begin{enumerate}[label=(\roman*)]
\item the ascent: pattern on $[p,i_c]$ (from its minimum $\mu$ to its maximum $c$), times $[a,r]$ (readings in $[\mu,c]$
since $r<f_s$; anchors $\mathrm{Al}_e(a,\mu)$, $\mathrm{Al}_{e+2}(r,c)$);
\item the descent: pattern on $[i_c,i_r]$, which goes from its maximum $c$ to its minimum $\mu$, with $x$ and $\tau$
negated so that it goes from minimum to maximum; times $[r,s]$;
\item the rest: pattern on $[i_r,q]$ (from $x(i_r)=\mu$, its minimum, to $M$), times $[s,b]$.
\end{enumerate}
Each time interval is shorter than $[a,b]$. By induction (and negation, Lemma~\ref{lem:lift-ops}(c), for (ii)) we get
lifts from $p$ to $i_c$, from $i_c$ to $i_r$, and from $i_r$ to $q$; widen them to $[p,q]$ and glue.

\paragraph{Case 6: $M\ge\mu+2$, $x(p+1)=\mu+1$, and $x(i)>\mu$ for all $i\in(p,q]$.}
Let $g$ be the last time in $[a,b]$ which is either $a$ or has $\ell(g)=\mu$. Then $\ell(f)\ge\mu+1$ for $g<f\le b$, and
$\mathrm{Al}_e(g,\mu)$. By Lemma~\ref{lem:allowed}(d), $\abs{\sigma(g)-\sigma(b)}\ge w+1-2e$, so $g<b$ and indeed $b-g\ge
q-p$.

\emph{After $g$.} $\mathrm{Al}_{e+1}(g+1,\mu+1)$: if $\ell(g)=\mu$, then $\ell(g+1)=\mu+1$; otherwise $g=a$ and
$\ell(a)=\mu+1$, and we move by one time. The tuple $(\tau,x,p+1,q,g+1,b,e+1)$ is good (the pattern on $[p+1,q]$ has minimum
$x(p+1)=\mu+1$ since it never returns to $\mu$; the readings after $g$ are $\ge\mu+1$), so by induction there is a lift over
$[g+1,b]$ from $p+1$ to $q$. Before it, put the step from $p$ at time $g$ to $p+1$ at time $g+1$.

\emph{Up to $g$: a lift over $[a,g]$ from $p$ to $p$.} If $g=a$ this is a single time. Otherwise $\ell(g)=\mu$. Let $f_k$
be a time in $[a,g]$ where $\ell$ is largest, and $E=\ell(f_k)$.
\begin{itemize}
\item If $E\le\mu+1$, let $\varphi(f)=p+1$ if $f\notin\set{a,g}$ and $\ell(f)=\mu+1$, and $\varphi(f)=p$ otherwise. This has
unit steps, and it is a lift: $x(p+1)=\mu+1=\ell(f)$, resp.\ $x(p)=\mu=\ell(f)$ (for $a<f<g$ the reading is $\mu$ or $\mu+1$),
and at $a$ use the anchor, at $g$ the reading $\ell(g)=\mu$.
\item If $E\ge\mu+2$, then $E\le M$ and, by the intermediate value theorem, the pattern takes the value $E$; let $i_e$ be
the least index with $x(i_e)=E$. Then $p<i_e$ and $x\le E$ on $[p,i_e]$. By Lemma~\ref{lem:allowed}(d),
$f_k-a\ge q-p$ and $g-f_k\ge q-p$. The ascent (pattern on $[p,i_e]$, times $[a,f_k]$, radius $e$, upper anchor
$\mathrm{Al}_0(f_k,E)$) and the descent (pattern $k\mapsto x(p+i_e-k)$ on $[p,i_e]$, negated, times $[f_k,g]$, radius $0$,
anchors $\mathrm{Al}_0(f_k,E)$ and $\mathrm{Al}_0(g,\mu)$) are good tuples with shorter time intervals. Induction, negation
and reflection give lifts over $[a,f_k]$ from $p$ to $i_e$ and over $[f_k,g]$ from $i_e$ to $p$; glue them.
\end{itemize}
Gluing the three pieces gives the lift over $[a,b]$ from $p$ to $q$.

\paragraph{Exhaustiveness.} If neither Case~1 nor Case~2 applies, then $x(p+1)\ne\mu$, and since $x(p)$ is the minimum and
$x$ has unit steps, $x(p+1)=\mu+1$. If $M=\mu+1$, Case~3 or Case~4 applies. Otherwise $M\ge\mu+2$, and Case~5 or Case~6
applies.
\end{proof}

\formal{\lean{HPC.PatCore.main_step} with the cases \lean{HPC.PatCore.case_const},
\lean{HPC.PatCore.case_plateau}, \lean{HPC.PatCore.case_plateau_end}, \lean{HPC.PatCore.case_amp1},
\lean{HPC.PatCore.case_return}, \lean{HPC.PatCore.case_noreturn}; the induction is \lean{HPC.PatCore.ih_all}.}

\begin{theorem}[Combinatorial core of pattern following]\label{thm:realize}
Let $\sigma$, $T$ satisfy~\ref{S:lip} and~\ref{S:crook}, let $N\ge1$ and $w\ge4N+10$, and let $\tau$ be $w$-slow. Let
$x\colon\N\to\Z$ have unit steps on $[0,N-1]$ with $x(0)\le x(i)\le x(N-1)$ for $i<N$. Suppose $x(0)\le\ell(f)\le x(N-1)$
for $f\le T$, $\ell(0)=x(0)$, $\ell(T)=x(N-1)$, and $N-1\le T$. Then there is $\varphi\colon[0,T]\to[0,N-1]$ with unit
steps, $\varphi(0)=0$, $\varphi(T)=N-1$, and $\mathrm{Al}_w(f,x(\varphi(f)))$ for all $f\le T$.
\end{theorem}

\begin{proof}
Take $B=2(N-1)$; then $2B+10\le w$. The tuple $(\tau,x,0,N-1,0,T,0)$ is good: the anchors hold with radius $0$ because
$\ell(0)=x(0)$ and $\ell(T)=x(N-1)$. Apply Theorem~\ref{thm:pf-main}.
\end{proof}

\formal{\lean{HPC.PatCore.realize}.}

\section{Back to chains}\label{sec:pf-geometry}

Let $D^{(0)},D^{(1)},\dots$ be a description from $P$ to $Q$ in a compact metric space $X$, with mesh bounds
$\varepsilon_i\to0$ and lengths $n_i$.

\begin{lemma}\label{lem:n-tendsto}
If $P\neq Q$, then $n_i\to\infty$.
\end{lemma}

\begin{proof}
Both $P$ and $Q$ lie in $\bigcup D^{(i)}$. Let $P\in D^{(i)}_a$ and $Q\in D^{(i)}_b$. By \vone{Lemma~2.9},
$d(P,Q)\le(\abs{a-b}+1)\varepsilon_i\le n_i\varepsilon_i$. As $\varepsilon_i\to0$, $n_i\ge d(P,Q)/\varepsilon_i\to\infty$.
\end{proof}

\formal{\lean{HPC.Description.n_tendsto_atTop} (via the more general \lean{HPC.Chain.diam_le_length_mul}).}

\begin{lemma}[Thick readings]\label{lem:thick}
For every $w\in\N$ there is $m_0$ such that for all $m\ge m_0$ and every index $p$ of $D^{(m)}$ there is an index $s$ of
$D^{(0)}$ with $D^{(m)}_{p'}\subseteq D^{(0)}_s$ for every $p'$ with $\abs{p-p'}\le w$.
\end{lemma}

\begin{proof}
The set $K=\cl{\bigcup D^{(1)}}$ is compact, and $K\subseteq\bigcup D^{(0)}$ because the closure of each link of $D^{(1)}$
lies in a link of $D^{(0)}$. By the Lebesgue number lemma there is $\delta>0$ such that every ball of radius $\delta$ about a
point of $K$ lies in a link of $D^{(0)}$. Choose $m_0\ge1$ with $\varepsilon_m(w+2)<\delta$ for $m\ge m_0$. Given $p$, pick
$y\in D^{(m)}_p$; then $y\in K$ (as $D^{(m)}$ refines $D^{(1)}$), so the ball $B(y,\delta)$ lies in some $D^{(0)}_s$. If
$\abs{p-p'}\le w$ and $z\in D^{(m)}_{p'}$, then $d(z,y)\le(w+1)\varepsilon_m<\delta$ by \vone{Lemma~2.9}, so $z\in
D^{(0)}_s$.
\end{proof}

\formal{\lean{HPC.Description.eventually_thick}.}

\begin{lemma}\label{lem:setting-of-crooked}
Let $F$ (of length $n_F\ge1$) be crooked in $R$, and let $\sigma_0\colon[n_F]\to[n_R]$ satisfy $F_t\subseteq R_{\sigma_0(t)}$.
Then $\sigma(f)=\sigma_0(\min(f,n_F-1))$ satisfies~\ref{S:lip} and~\ref{S:crook} with $T=n_F-1$.
\end{lemma}

\begin{proof}
\ref{S:lip} is \vone{Lemma~2.4} (and $\sigma$ is constant after $n_F-1$). For~\ref{S:crook} let $f<f'\le T$ with
$\abs{\sigma(f)-\sigma(f')}\ge3$; say $\sigma(f')>\sigma(f)+2$ (the other case is symmetric). The link $F_f$ meets
$R_{\sigma(f)}$ and $F_{f'}$ meets $R_{\sigma(f')}$, so crookedness gives $r,s$ with $\mathrm{Btw}(f,r,s,f')$, i.e.
$f<r<s<f'$, and $F_r\subseteq R_{\sigma(f')-1}$, $F_s\subseteq R_{\sigma(f)+1}$. Since $F_r$ also lies in $R_{\sigma(r)}$, these
two links of $R$ meet and $\abs{\sigma(r)-(\sigma(f')-1)}\le1$; so $\abs{\sigma(r)-\sigma(f')}\le2$. Likewise
$\abs{\sigma(s)-\sigma(f)}\le2$.
\end{proof}

\formal{\lean{HPC.Chain.setting_of_crooked}.}

\begin{lemma}\label{lem:consolidation-of-lift}
Let $F$ be a chain of length $n_F$ from $P$ to $Q$, $D$ a chain, and $x\colon[N]\to[n]$. Let $\varphi\colon[n_F]\to[N]$ have
unit steps, $\varphi(0)=0$, $\varphi(n_F-1)=N-1$, and $F_f\subseteq D_{x(\varphi(f))}$ for all $f$. Then there is a
consolidation $E$ of $F$ of length $N$ which goes from $P$ to $Q$ and follows $x$ in $D$.
\end{lemma}

\begin{proof}
$\varphi$ is surjective by the discrete intermediate value theorem, so it is a block map; let $E=F^\varphi$
(Lemma~\ref{lem:block-chain}). The only link of $F$ containing $P$ is $F_0$, in block $0$; the only one containing $Q$ is
$F_{n_F-1}$, in block $N-1$; so $E$ goes from $P$ to $Q$ by Lemma~\ref{lem:block-from}. Finally
$E_u=\bigcup_{\varphi(f)=u}F_f\subseteq D_{x(u)}$.
\end{proof}

\formal{\lean{HPC.Chain.exists_consolidation_of_lift}.}

\begin{proof}[Proof of Theorem~\ref{thm:pattern-following}]
The levels $D^{(i)},D^{(i+1)},\dots$ form a description from $P$ to $Q$ (with the same properties), so we may assume $i=0$
and must find $j>0$. Let $x\colon[N]\to[n_0]$ be a pattern in $D^{(0)}$ and $w=4N+10$. By Lemmas~\ref{lem:thick}
and~\ref{lem:n-tendsto} choose $m$ such that the conclusion of Lemma~\ref{lem:thick} holds for $D^{(m)}$, and $n_{m+1}\ge N$.
Let $R=D^{(m)}$, $F=D^{(m+1)}$ and $T=n_{m+1}-1$. For each index $p$ of $R$ let $\tau_0(p)$ be an index $s$ as in
Lemma~\ref{lem:thick}, and let $\sigma_0(t)$ be an index with $F_t\subseteq R_{\sigma_0(t)}$.

Extend to $\tau(p)=\tau_0(\text{clamp}(p))$ for $p\in\Z$, where $\text{clamp}(p)$ is the nearest index of $R$ to $p$, and
let $\sigma$ be as in Lemma~\ref{lem:setting-of-crooked}; also extend $x$ to $\N$ by $x(i)=x(\min(i,N-1))$.
\begin{itemize}
\item $\tau$ is $w$-slow: if $\abs{p-p'}\le w$, then $R_{p'}$ (with $p,p'$ clamped) lies in $D^{(0)}_{\tau(p)}$ and in
$D^{(0)}_{\tau(p')}$, so these links meet and $\abs{\tau(p)-\tau(p')}\le1$. (Clamping does not increase distances.)
\item $\sigma$ satisfies~\ref{S:lip} and~\ref{S:crook} with this $T$, by Lemma~\ref{lem:setting-of-crooked}, since $F$ is
crooked in $R$.
\item The pattern conditions hold: $x(0)=0\le x(i)\le n_0-1=x(N-1)$, and every reading $\ell(f)$ lies in $[0,n_0-1]$.
\item $\ell(0)=0$: $P\in F_0\subseteq R_{\sigma(0)}$, and $P$ lies only in $R_0$, so $\sigma(0)=0$; and $P\in R_0\subseteq
D^{(0)}_{\tau_0(0)}$, and $P$ lies only in $D^{(0)}_0$, so $\tau_0(0)=0$. Likewise $\ell(T)=n_0-1$ using $Q$.
\item $N-1\le T$.
\end{itemize}
Theorem~\ref{thm:realize} gives $\varphi\colon[0,T]\to[0,N-1]$ with unit steps from $0$ to $N-1$ such that for every $f$
there is a position $p$ with $\abs{p-\sigma(f)}\le w$ and $\tau(p)=x(\varphi(f))$. Then
\[
F_f\subseteq R_{\sigma(f)}\subseteq D^{(0)}_{\tau_0(p)}=D^{(0)}_{x(\varphi(f))},
\]
the middle inclusion by the choice of $\tau_0$, since $\sigma(f)$ is within $w$ of $p$. Lemma~\ref{lem:consolidation-of-lift}
gives the required consolidation of $F=D^{(m+1)}$; so $j=m+1>0$.
\end{proof}

\formal{\lean{HPC.Description.theorem6Pos_of_ne}, and
\lean{HPC.Description.patternFollowing_of_theorem6Pos} for the passage to arbitrary $i$ (via the shifted description
\lean{HPC.Description.shift}).}

\begin{remark}
Bing derives Theorem~6 from a more general Theorem~5 (patterns between any two links $h$ and $k$, followed by
consolidations of subchains). The formal development also proves Theorem~5 (\lean{HPC.Description.theorem5}), but
the uniqueness theorem needs only Theorem~6, and the direct proof above is simpler.
\end{remark}

\section*{Exercises}
\addcontentsline{toc}{section}{Exercises}

\begin{exercise}
Check the unit-step claim for $\varphi$ in Case~4 of the proof of Theorem~\ref{thm:pf-main} at each of the four
transition times $a+1$, $u$, $u+W$ and $b$.
\end{exercise}

\begin{exercise}
Show that crookedness~\ref{S:crook} cannot be dropped from Theorem~\ref{thm:realize}: take $\sigma(f)=f$, $\tau(p)=\fl{p/(w+1)}$,
and a pattern $x$ which, starting from its minimum, rises by at least $4$, then falls by at least $4$, and then rises to its
maximum at the end. Show that no lift exists. (Hint: $\ell$ is nondecreasing.)
\end{exercise}

\begin{exercise}
In Lemma~\ref{lem:consolidation-of-lift}, show that the blocks of $E$ need not be sets of consecutive indices, and draw an example.
\end{exercise}

\begin{exercise}
Where is the compactness of $X$ used in the proof of Theorem~\ref{thm:pattern-following}? Show that it suffices that the
links of the description lie in a compact set.
\end{exercise}

%======================================================================
\chapter{Matching covers}\label{ch:matching}
%======================================================================

This chapter proves a general criterion for two compact sets to be homeomorphic. It contains no
chains: the input is two sequences of finite open covers, indexed in the same way, such that
``meeting'' on one side forces ``closeness'' on the other. Chapter~\ref{ch:two} produces such covers
from two descriptions. The criterion is Bing's Theorem~11 (1948), with an extra clause that tracks
individual points.

Throughout, $X_1$ and $X_2$ are metric spaces, $M_1\subseteq X_1$ and $M_2\subseteq X_2$ are compact, and
$\varepsilon_0,\varepsilon_1,\dots$ are positive numbers with $\sum_i\varepsilon_i<\infty$.

\section{The statement}

\begin{definition}\label{def:matching}
\emph{Matching covers} of $M_1$ and $M_2$ (for $\varepsilon$) consist of numbers $N_i\ge0$ and open sets
$x^1_{i,k}\subseteq X_1$, $x^2_{i,k}\subseteq X_2$ for $i\in\N$ and $k\in[N_i]$ such that, for $c=1,2$:
\begin{enumerate}[label=(M\arabic*)]
\item\label{M:cover} $M_c\subseteq\bigcup_kx^c_{i,k}$ for every $i$;
\item\label{M:meet} every $x^c_{i,k}$ meets $M_c$;
\item\label{M:diam} every $x^c_{i,k}$ is bounded with $\diam x^c_{i,k}\le\varepsilon_i$;
\item\label{M:match} if $x^1_{i+1,j}$ meets $x^1_{i,k}$, then some point of $x^2_{i+1,j}$ is at distance $<\varepsilon_i$ from
some point of $x^2_{i,k}$; and the same with the sides $1$ and $2$ exchanged.
\end{enumerate}
A sequence of indices $a=(a_0,a_1,\dots)$ with $a_i\in[N_i]$ \emph{tracks} a point $p\in X_c$ if $p\in x^c_{i,a_i}$ for
every $i$.
\end{definition}

\formal{\lean{HPC.MatchingCovers}, \leansub{HPC.MatchingCovers.Tracks}{1}, \leansub{HPC.MatchingCovers.Tracks}{2}.}

\begin{theorem}[Matching covers; Bing 1948, Theorem~11]\label{thm:matching}
If $M_1$ and $M_2$ are compact and carry matching covers, then there is a homeomorphism $T\colon M_1\to M_2$ such that
$T(p)=q$ whenever $p\in M_1$, $q\in M_2$ and one index sequence tracks both $p$ and $q$.
\end{theorem}

\formal{\lean{HPC.MatchingCovers.exists_homeomorph}.}

The definition is symmetric in the two sides, so every statement proved below for one side also holds for the other.

\section{Fattening}

Let $t_i=\sum_{l\ge0}\varepsilon_{i+l}$ (the tail sums) and define the \emph{fattening radius} and the \emph{mesh}
\[
\rho_i=\varepsilon_i+2t_{i+1},\qquad \mu_i=\varepsilon_i+2\rho_i .
\]
Since $t_i\to0$ and $\varepsilon_i\to0$, also $\rho_i\to0$ and $\mu_i\to0$. Moreover $t_{i+1}=\varepsilon_{i+1}+t_{i+2}$, which
gives the identity
\begin{equation}\label{eq:rad}
\rho_{i+1}+\varepsilon_{i+1}+\varepsilon_i=\rho_i .
\end{equation}
For a set $S$ and $\delta>0$ let $S^\delta=\set{z: d(z,s)<\delta\text{ for some }s\in S}$, an open set containing $S$, with
$\diam S^\delta\le\diam S+2\delta$. Put
\[
y^c_{i,k}=\bigl(x^c_{i,k}\bigr)^{\rho_i}.
\]
Then $\diam\cl{y^c_{i,k}}=\diam y^c_{i,k}\le\varepsilon_i+2\rho_i=\mu_i$: any two points of $\cl{y^c_{i,k}}$ are at distance at most
$\mu_i$.

\formal{\lean{HPC.tail}, \lean{HPC.rad}, \lean{HPC.mesh}, \lean{HPC.fat}, \leansub{HPC.MatchingCovers.y}{1},
\leansub{HPC.MatchingCovers.y}{2}, \lean{HPC.diam_fat_le_mesh}, \lean{HPC.mesh_tendsto_zero}.}

\begin{lemma}[Nesting]\label{lem:nesting}
Let $A,B$ be sets with $B$ bounded, $\diam B\le\varepsilon_{i+1}$, and suppose some $p\in B$ and $q\in A$ have
$d(p,q)<\varepsilon_i$. Then $\cl{B^{\rho_{i+1}}}\subseteq A^{\rho_i}$.
\end{lemma}

\begin{proof}
Let $z\in\cl{B^{\rho_{i+1}}}$ and $\eta=\varepsilon_i-d(p,q)>0$. There is $z'\in B^{\rho_{i+1}}$ with $d(z,z')<\eta$, and
$b'\in B$ with $d(z',b')<\rho_{i+1}$; also $d(b',p)\le\diam B\le\varepsilon_{i+1}$. Hence
\[
d(z,q)\le d(z,z')+d(z',b')+d(b',p)+d(p,q)<\eta+\rho_{i+1}+\varepsilon_{i+1}+d(p,q)=\rho_{i+1}+\varepsilon_{i+1}+\varepsilon_i=\rho_i
\]
by~\eqref{eq:rad}. So $z\in A^{\rho_i}$.
\end{proof}

\formal{\lean{HPC.closure_fat_subset_fat}.}

\section{Nested fattened sets}

\begin{lemma}[Step 1]\label{lem:step1}
If $x^1_{i+1,s}$ meets $x^1_{i,r}$, then $\cl{y^2_{i+1,s}}\subseteq y^2_{i,r}$.
\end{lemma}

\begin{proof}
By~\ref{M:match} some point of $x^2_{i+1,s}$ is within $\varepsilon_i$ of some point of $x^2_{i,r}$, and
$\diam x^2_{i+1,s}\le\varepsilon_{i+1}$. Apply Lemma~\ref{lem:nesting}.
\end{proof}

\formal{\leansub{HPC.MatchingCovers.step1}{2}.}

\begin{lemma}[Step 2]\label{lem:step2}
Let $i<j$. If $x^1_{i,r}\cap x^1_{j,s}\cap M_1\ne\emptyset$, then $\cl{y^2_{j,s}}\subseteq y^2_{i,r}$.
\end{lemma}

\begin{proof}
Induction on $j$. Let $c\in x^1_{i,r}\cap x^1_{j,s}\cap M_1$. If $j=i+1$ this is Lemma~\ref{lem:step1}. If $j>i+1$, by~\ref{M:cover} there is
$t$ with $c\in x^1_{j-1,t}$. Then $x^1_{j,s}$ meets $x^1_{j-1,t}$, so $\cl{y^2_{j,s}}\subseteq y^2_{j-1,t}$ by
Lemma~\ref{lem:step1}; and $c\in x^1_{i,r}\cap x^1_{j-1,t}\cap M_1$, so $\cl{y^2_{j-1,t}}\subseteq y^2_{i,r}$ by induction. Combine.
\end{proof}

\formal{\leansub{HPC.MatchingCovers.step2}{2}.}

\section{The map}

\begin{lemma}[Step 3]\label{lem:limit-point}
Let $a$ be an index sequence with $x^1_{i+1,a_{i+1}}\cap x^1_{i,a_i}\neq\emptyset$ for all $i$. Then there is $q\in M_2$ with
$q\in\cl{y^2_{i,a_i}}$ for every $i$.
\end{lemma}

\begin{proof}
The sets $K_i=\cl{y^2_{i,a_i}}\cap M_2$ are compact; they decrease, because $\cl{y^2_{i+1,a_{i+1}}}\subseteq y^2_{i,a_i}$ by
Lemma~\ref{lem:step1}; and they are nonempty, because $x^2_{i,a_i}\subseteq y^2_{i,a_i}$ meets $M_2$ by~\ref{M:meet}. By Cantor's
intersection theorem \vone{(C3)} they have a common point.
\end{proof}

\formal{\lean{HPC.MatchingCovers.exists_limit_point}.}

For $p\in M_1$ choose, for each $i$, an index $a_i(p)$ with $p\in x^1_{i,a_i(p)}$ (by~\ref{M:cover}). The sequence $a(p)$
satisfies the hypothesis of Lemma~\ref{lem:limit-point} (consecutive sets contain $p$). Let $T(p)\in M_2$ be a point given by
that lemma.

\begin{lemma}[Step 4: independence of choices]\label{lem:T-indep}
If $p\in M_1$ and $p\in x^1_{j,r}$, then $T(p)\in y^2_{j,r}$.
\end{lemma}

\begin{proof}
$p\in x^1_{j+1,a_{j+1}(p)}\cap x^1_{j,r}$, so $\cl{y^2_{j+1,a_{j+1}(p)}}\subseteq y^2_{j,r}$ by Lemma~\ref{lem:step1}, and $T(p)$
lies in the left side by construction.
\end{proof}

\formal{\lean{HPC.MatchingCovers.toFun}, \leansub{HPC.MatchingCovers.toFun_mem_y}{2}.}

\begin{lemma}[Step 5]\label{lem:T-cont}
$T\colon M_1\to M_2$ is continuous.
\end{lemma}

\begin{proof}
Let $p\in M_1$ and $\eta>0$. Choose $j$ with $\mu_j<\eta$ and $r$ with $p\in x^1_{j,r}$. As $x^1_{j,r}$ is open there is
$\delta>0$ such that every $p'$ with $d(p,p')<\delta$ lies in $x^1_{j,r}$. For such $p'\in M_1$, both $T(p)$ and $T(p')$ lie in
$y^2_{j,r}$ by Lemma~\ref{lem:T-indep}, so $d(T(p),T(p'))\le\mu_j<\eta$.
\end{proof}

\formal{\lean{HPC.MatchingCovers.continuous_toMap}.}

\begin{lemma}[Step 6]\label{lem:T-inj}
$T$ is injective.
\end{lemma}

\begin{proof}
Suppose $p_1\neq p_2$ in $M_1$ with $T(p_1)=T(p_2)=Q$. Choose $k$ with $\mu_k<d(p_1,p_2)$, and $r$ with $Q\in x^2_{k,r}$ (note
$Q\in M_2$). As $x^2_{k,r}$ is open, there is $\eta>0$ with $B(Q,\eta)\subseteq x^2_{k,r}$. Choose $j>k$ with $\mu_j<\eta$, and $u,v$
with $p_1\in x^1_{j,u}$, $p_2\in x^1_{j,v}$. By Lemma~\ref{lem:T-indep}, $Q\in y^2_{j,u}$ and $Q\in y^2_{j,v}$.

\emph{Claim: if $Q\in y^2_{j,w}$, then $\cl{y^1_{j,w}}\subseteq y^1_{k,r}$.} By~\ref{M:meet} pick $c\in x^2_{j,w}\cap M_2$. Both $c$
and $Q$ lie in $y^2_{j,w}$, so $d(c,Q)\le\mu_j<\eta$ and $c\in x^2_{k,r}$. Thus $x^2_{k,r}\cap x^2_{j,w}\cap M_2\neq\emptyset$, and
Lemma~\ref{lem:step2} with the sides exchanged gives the claim.

By the claim, $p_1\in x^1_{j,u}\subseteq y^1_{j,u}\subseteq y^1_{k,r}$ and likewise $p_2\in y^1_{k,r}$. Hence $d(p_1,p_2)\le\mu_k$,
a contradiction.
\end{proof}

\formal{\lean{HPC.MatchingCovers.injective_toMap}.}

\begin{lemma}[Step 7]\label{lem:T-surj}
$T$ is surjective.
\end{lemma}

\begin{proof}
The image $T(M_1)$ is compact (Lemma~\ref{lem:T-cont}), hence closed in $M_2$, so it suffices to show it is dense. Let $q\in M_2$
and $\eta>0$. Choose $j$ with $\mu_j<\eta$ and $r$ with $q\in x^2_{j,r}$, and by~\ref{M:meet} a point $p\in x^1_{j,r}\cap M_1$.
Then $T(p)\in y^2_{j,r}$ (Lemma~\ref{lem:T-indep}) and $q\in y^2_{j,r}$, so $d(T(p),q)\le\mu_j<\eta$.
\end{proof}

\formal{\lean{HPC.MatchingCovers.surjective_toMap}.}

\begin{proof}[Proof of Theorem~\ref{thm:matching}]
$T$ is a continuous bijection from the compact space $M_1$ to the metric space $M_2$, so it is a homeomorphism (a continuous
bijection from a compact space to a Hausdorff space maps closed sets, which are compact, to compact, hence closed, sets).
For the tracking clause, let $a$ track $p\in M_1$ and $q\in M_2$. For every $i$, $T(p)\in y^2_{i,a_i}$ by
Lemma~\ref{lem:T-indep}, and $q\in x^2_{i,a_i}\subseteq y^2_{i,a_i}$; so $d(T(p),q)\le\mu_i$ for all $i$. Since $\mu_i\to0$,
$T(p)=q$.
\end{proof}

\section*{Exercises}
\addcontentsline{toc}{section}{Exercises}

\begin{exercise}
Show that the summability of $\varepsilon$ cannot simply be replaced by $\varepsilon_i\to0$ in the proof: where is
$t_i<\infty$ used? (It is used to make the fattened sets nest.)
\end{exercise}

\begin{exercise}
Let $M_1=M_2=[0,1]$ and let $x^1_{i,k}=x^2_{i,k}$ be the intervals $\bigl((k-1)2^{-i},(k+1)2^{-i}\bigr)\cap(-1,2)$,
$k\in[2^i+1]$. Find $\varepsilon$ for which these are matching covers, and identify the homeomorphism $T$.
\end{exercise}

\begin{exercise}
Show that if one replaces \ref{M:match} by the one-sided condition (from side $1$ to side $2$ only), the map $T$ is still
defined and continuous, but may fail to be injective.
\end{exercise}

%======================================================================
\chapter{Two descriptions give a homeomorphism}\label{ch:two}
%======================================================================

This chapter proves Part~I of the uniqueness theorem.

\begin{theorem}[Bing 1948, Theorem~12]\label{thm:two-desc}
Let $X_1,X_2$ be compact metric spaces, and for $c=1,2$ let $D^{(c,0)},D^{(c,1)},\dots$ be a description from $P_c$ to $Q_c$ in
$X_c$, with $P_c\neq Q_c$ and described set $M_c$. Then there is a homeomorphism $T\colon M_1\to M_2$ with $T(P_1)=P_2$ and
$T(Q_1)=Q_2$.
\end{theorem}

\formal{\lean{HPC.theorem12}, which combines \lean{HPC.Description.exists_homeomorph} (the theorem assuming pattern following)
with \lean{HPC.Description.patternFollowing}.}

The idea: build, alternately in the two spaces, chains $C^{(1,i)}$ in $X_1$ and $C^{(2,i)}$ in $X_2$ with the same number of links
$N_i$, such that $C^{(1,i+1)}$ and $C^{(2,i+1)}$ follow the \emph{same} pattern in $C^{(1,i)}$ and $C^{(2,i)}$. Then the links of
these chains are matching covers (Chapter~\ref{ch:matching}). At each step one side is refined by a deep level of its own
description, which yields a pattern, and the other side follows this pattern by Theorem~\ref{thm:pattern-following}.

\section{Patterns from refinements, and matching}

\begin{lemma}\label{lem:pattern-of-refines}
Let $Y$ (of length $N$) and $Z$ (of length $m$) be chains from $P$ to $Q$, with $Y$ refining $Z$. Then any refinement map
$x\colon[N]\to[m]$ is a pattern in $Z$, and $Y$ follows $x$ in $Z$.
\end{lemma}

\begin{proof}
$x$ has unit steps by \vone{Lemma~2.4}. $P\in Y_0\subseteq Z_{x(0)}$ and $P$ lies only in $Z_0$, so $x(0)=0$; likewise
$x(N-1)=m-1$ using $Q$. The last statement is the definition of a refinement map.
\end{proof}

\formal{\lean{HPC.Chain.exists_pattern_of_refines}.}

\begin{lemma}[Matching]\label{lem:pattern-matching}
Let $Y_1,Z_1$ be chains in $X_1$ and $Y_2,Z_2$ chains in $X_2$, with $Y_1,Y_2$ of length $N$ and $Z_1,Z_2$ of length $m$. Suppose $Y_1$
follows a function $x\colon[N]\to[m]$ in $Z_1$ and $Y_2$ follows the same $x$ in $Z_2$, and the links of $Z_2$ are bounded with
diameter $<\varepsilon$. If $(Y_1)_j$ meets $(Z_1)_k$, then some point of $(Y_2)_j$ is at distance $<\varepsilon$ from some point of
$(Z_2)_k$.
\end{lemma}

\begin{proof}
A point of $(Y_1)_j\cap(Z_1)_k$ lies in $(Z_1)_{x(j)}\cap(Z_1)_k$, so $\abs{x(j)-k}\le1$, and therefore $(Z_2)_{x(j)}$ meets $(Z_2)_k$,
say at $q$. Take any $p\in(Y_2)_j$. Then $p,q\in(Z_2)_{x(j)}$, so $d(p,q)\le\diam(Z_2)_{x(j)}<\varepsilon$.
\end{proof}

\formal{\lean{HPC.Chain.pattern_matching}.}

\begin{lemma}\label{lem:consolidation-meets}
Let $C$ be a consolidation of a level $D^{(l)}$ of a description in a compact space, with described set $M$. Then
$M\subseteq\bigcup C$, and every link of $C$ meets $M$.
\end{lemma}

\begin{proof}
$M\subseteq\bigcup D^{(l)}\subseteq\bigcup C$ since $D^{(l)}$ refines $C$. A link of $C$ is a nonempty union of links of $D^{(l)}$, and
every link of $D^{(l)}$ meets $M$ by \vone{Proposition~3.9}.
\end{proof}

\formal{\lean{HPC.Description.M_subset_consolidation}, \lean{HPC.Description.consolidation_link_inter_M_nonempty},
\lean{HPC.Description.link_inter_M_nonempty}.}

\begin{lemma}\label{lem:level-refining}
Let $C$ be a consolidation of the level $D^{(l)}$ and $\delta>0$. There is $k\ge l$ such that $D^{(k)}$ refines $C$ and
$\varepsilon_k\le\delta$.
\end{lemma}

\begin{proof}
Choose $k\ge l$ with $\varepsilon_k\le\delta$. Then $D^{(k)}$ refines $D^{(l)}$ (\vone{Lemma~3.2(b)}), which refines $C$.
\end{proof}

\formal{\lean{HPC.Description.exists_level_refining}.}

\section{The alternating construction}

Fix the two descriptions of Theorem~\ref{thm:two-desc}; write $n^c_l$ and $\varepsilon^c_l$ for the lengths and mesh bounds
of the description in $X_c$, and $D^{(c,l)}$ for its levels. By Theorem~\ref{thm:pattern-following} both descriptions have the
pattern following property.

\begin{definition}\label{def:rung}
A \emph{rung} consists of a number $N$, chains $C^1$ in $X_1$ and $C^2$ in $X_2$, both of length $N$, and numbers $m>0$,
$m_f>0$, such that
\begin{enumerate}[label=(R\arabic*)]
\item $C^c$ goes from $P_c$ to $Q_c$ and is a consolidation of some level of the description in $X_c$ ($c=1,2$);
\item every link of $C^1$ and of $C^2$ has diameter $\le m$;
\item one of the two chains, the \emph{active} one, is itself a level of its description, and its links have diameter $\le m_f$.
\end{enumerate}
\end{definition}

\formal{\lean{HPC.Rung}. Formally, (R3) asks that the active side be able to follow patterns
(\lean{HPC.Description.CanFollow}), which holds for every level by pattern following
(\lean{HPC.Description.canFollow_level}).}

\begin{lemma}[The first rung]\label{lem:base-rung}
There is a rung with $m_f=1/2$.
\end{lemma}

\begin{proof}
Choose a level $a$ of the first description with $\varepsilon^1_a<1/2$ and $\varepsilon^1_a<d(P_1,Q_1)$. Then $n^1_a\ge2$, since $P_1$
and $Q_1$ cannot lie in one link. By Lemma~\ref{lem:n-tendsto} choose $b$ with $n^2_b\ge n^1_a$, and by
Lemma~\ref{lem:consolidate-length} a consolidation $F$ of $D^{(2,b)}$ of length $n^1_a$ from $P_2$ to $Q_2$. Take $C^1=D^{(1,a)}$
(active, $m_f=1/2$), $C^2=F$, and $m=\max(1/2,\diam X_2+1)$; $X_2$ is compact, hence bounded.
\end{proof}

\formal{\lean{HPC.Rung.exists_base}.}

\begin{lemma}[The next rung]\label{lem:next-rung}
Given a rung $(N,C^1,C^2,m,m_f)$ and $\delta>0$, there are a rung $(N',C'^1,C'^2,m',\delta)$ with $m'=\max(\delta,m_f)$ and a
function $x\colon[N']\to[N]$ such that $C'^1$ follows $x$ in $C^1$ and $C'^2$ follows $x$ in $C^2$.
\end{lemma}

\begin{proof}
Suppose $C^1$ is active, $C^1=D^{(1,a)}$ (the other case is symmetric). The inactive side $C^2$ is a consolidation of some level of
the second description; by Lemma~\ref{lem:level-refining} there is a level $c$ such that $D^{(2,c)}$ refines $C^2$ and
$\varepsilon^2_c\le\delta$. By Lemma~\ref{lem:pattern-of-refines} a refinement map $x\colon[n^2_c]\to[N]$ is a pattern in $C^2$, followed
by $D^{(2,c)}$. Since $C^1=D^{(1,a)}$ is a level, pattern following (Theorem~\ref{thm:pattern-following}) gives $a'>a$ and a
consolidation $C'^1$ of $D^{(1,a')}$ of length $n^2_c$, from $P_1$ to $Q_1$, following $x$ in $C^1$. Let $N'=n^2_c$ and
$C'^2=D^{(2,c)}$, now the active side, with $m_f'=\delta$. The links of $C'^2$ have diameter $\le\varepsilon^2_c\le\delta$; each link of
$C'^1$ lies in a link of $C^1$, of diameter $\le m_f$. So $m'=\max(\delta,m_f)$ bounds all diameters.
\end{proof}

\formal{\lean{HPC.Rung.exists_next}, via \lean{HPC.Description.exists_half_step}.}

\section{Proof of Theorem~\ref{thm:two-desc}}

\begin{proof}
Start with the rung $R_0$ of Lemma~\ref{lem:base-rung} and let $R_{k+1}$ be obtained from $R_k$ by Lemma~\ref{lem:next-rung}
with $\delta=2^{-(k+2)}$. Write $R_i=(N_i,C^{(1,i)},C^{(2,i)},m_i,m_{f,i})$. Then $m_{f,i}=2^{-(i+1)}$ for all $i$, and for
$i\ge1$, $m_i=\max(2^{-(i+1)},2^{-i})=2^{-i}$. Put $\varepsilon_i=2m_i$; these numbers are positive and summable.

We check that the links $x^c_{i,k}=C^{(c,i)}_k$ ($k\in[N_i]$) are matching covers of $M_1$ and $M_2$ for $\varepsilon$
(Definition~\ref{def:matching}). The links are open; \ref{M:cover} and~\ref{M:meet} hold by Lemma~\ref{lem:consolidation-meets};
\ref{M:diam} holds since diameters are at most $m_i<\varepsilon_i$ (links are bounded because $X_c$ is compact). For~\ref{M:match}, by
Lemma~\ref{lem:next-rung} there is $x$ such that $C^{(1,i+1)}$ and $C^{(2,i+1)}$ follow $x$ in $C^{(1,i)}$ and $C^{(2,i)}$
respectively; Lemma~\ref{lem:pattern-matching} gives~\ref{M:match} in both directions, with $\varepsilon=\varepsilon_i$.

Also $M_1$ and $M_2$ are compact (closed subsets of compact spaces, \vone{Lemma~3.2(d)}). Theorem~\ref{thm:matching} gives a
homeomorphism $T\colon M_1\to M_2$ with the tracking property. The constant sequence $a_i=0$ tracks $P_1$ and $P_2$, because each
$C^{(c,i)}$ goes from $P_c$ to $Q_c$; so $T(P_1)=P_2$. Likewise the sequence $a_i=N_i-1$ tracks $Q_1$ and $Q_2$, so $T(Q_1)=Q_2$.
\end{proof}

\section{Descriptions in non-compact spaces}

Theorem~\ref{thm:two-desc} assumes that the ambient spaces are compact. It is enough that the links lie in compact sets.

\begin{corollary}\label{cor:two-desc-K}
Let $D^{(c,i)}$ be descriptions from $P_c$ to $Q_c\ne P_c$ in arbitrary metric spaces $X_c$ ($c=1,2$), and suppose all links of the
$c$-th description lie in a compact set $K_c\subseteq X_c$. Then there is a homeomorphism $T\colon M_1\to M_2$ with $T(P_1)=P_2$ and
$T(Q_1)=Q_2$.
\end{corollary}

\begin{proof}
Regard each link $D^{(c,i)}_t\subseteq K_c$ as a subset of the compact metric space $K_c$. It is open in $K_c$, the chain conditions
and the conditions ``from $P_c$ to $Q_c$'', ``crooked'' and the diameter bounds do not change, and the closure of a link in $K_c$ is its
closure in $X_c$ intersected with $K_c$, so closures of links are still nested. Thus we have descriptions in the compact spaces $K_c$,
with the same described sets. Apply Theorem~\ref{thm:two-desc}.
\end{proof}

\formal{\lean{HPC.Plane.theorem12_of}, using the restriction \lean{HPC.Description.restrict} (in
\lean{HPC/Plane/Transfer.lean}).}

\begin{remark}
Bing states Theorem~12 under the assumption that the described sets are compact. This is not sufficient: \vone{Remark~3.5} gives a
description, in a non-complete space, whose described set is compact but not connected; it is not homeomorphic to the
pseudo-arc described in the plane. The formal development records this as \lean{HPC.Plane.theorem12_of_false}.
\end{remark}

\section*{Exercises}
\addcontentsline{toc}{section}{Exercises}

\begin{exercise}
Show that in Theorem~\ref{thm:two-desc} one may take the two descriptions to be the same. What does the theorem then say?
(Compare Chapter~\ref{ch:main}.)
\end{exercise}

\begin{exercise}
Why does the construction alternate between the two sides? Show that if one only ever refines the first side, the
mesh of the second side need not tend to zero.
\end{exercise}

\begin{exercise}
Show that the homeomorphism $T$ of Theorem~\ref{thm:two-desc} can be chosen so that also $T(R_1)=R_2$ for given points $R_c$ that
lie, for every $i$, in the $k_i$-th link of $C^{(c,i)}$ for the same indices $k_i$. (This is the tracking clause.)
\end{exercise}

%======================================================================
\chapter{The folding theorem}\label{ch:folding}
%======================================================================

This chapter proves the theorem about crooked chains that Part~II of the uniqueness proof needs: a
chain crooked in $D$, a stretch of which reaches both ends of $D$, can be consolidated into a chain
whose first and last blocks are prescribed and whose blocks are small. It is Bing's Theorem~4 (1948),
in the form that the formal proof establishes.

Throughout, $D$ is a chain of length $k$, $E$ is a chain of length $n$ which is crooked in $D$, and
$\lo(t)$, $\hi(t)$ are the least and the greatest index $j$ with $E_t\subseteq D_j$
(Definition~\ref{def:minidx}). For integers $w$ we write $D[w,w+2]$ for the union of those of
$D_w,D_{w+1},D_{w+2}$ that exist. A block map (Definition~\ref{def:block-map}) has \emph{small
blocks} if every block $E^\beta_u$ lies in some $D[w,w+2]$.

\section{The statement}

\begin{theorem}[Folding theorem; Bing 1948, Theorem~4]\label{thm:folding}
Let $k\ge3$, let $E$ be crooked in $D$, and let $r,r',s',s$ be indices of $E$ with $r'=r+1$, $s=s'+1$ and
$r'\le s'$. Suppose that some $E_h$ with $r\le h\le s$ meets $D_0$, and some $E_{h'}$ with $r\le h'\le s$
meets $D_{k-1}$. Then there is a block map $\beta\colon[n]\to[q]$ with small blocks such that
\[
\beta(r)=\beta(r')=0\qquad\text{and}\qquad\beta(s')=\beta(s)=q-1 .
\]
\end{theorem}

\formal{\lean{HPC.Fold.theorem4_pair} (in \lean{HPC/Fold/Theorem4Pair.lean}). The version with single links
$r$ and $s$ in place of the pairs is \lean{HPC.Fold.theorem4}.}

We need the pairs because of the following application. A point lies in at most two links of a chain, and
these are adjacent; for the block chain to go from $R$ to $S$ (Lemma~\ref{lem:block-from}) \emph{every} link
containing $R$ must be in the first block.

\begin{corollary}\label{cor:block-chain}
Let $k\ge5$ and let $E$ be crooked in $D$. Let $a<b$ be indices of $E$ such that the subchain from $a$ to $b$
goes from $R$ to $S$, and suppose some $E_{t_1}$ and some $E_{t_2}$ with $a\le t_1,t_2\le b$ meet $D_0$ and $D_{k-1}$
respectively. Then there is a block map $\beta$ for $E$ with small blocks such that the block chain $E^\beta$ goes
from $R$ to $S$.
\end{corollary}

\begin{proof}
First, $b\ge a+2$: the links $E_{t_1}$ and $E_{t_2}$ lie in links of $D$ with indices $\le1$ and $\ge k-2\ge3$
respectively, so they do not meet, $\abs{t_1-t_2}\ge2$. Next, the links of $E$ containing $R$: $R$ lies in $E_a$, and in
no $E_t$ with $a<t\le b$; since links containing $R$ meet, they are among $E_{a-1},E_a$. Similarly the links containing
$S$ are among $E_b,E_{b+1}$. Apply Theorem~\ref{thm:folding} with $r=a-1$, $r'=a$ if $a\ge1$ (and $r=a$, $r'=a+1$ if
$a=0$), and $s'=b$, $s=b+1$ if $b+1<n$ (and $s'=b-1$, $s=b$ if $b=n-1$). The hypotheses hold since $r\le a\le t_1,t_2\le
b\le s$ and $r'\le a+1\le b-1\le s'$. Every link containing $R$ is $E_r$ or $E_{r'}$, hence in block $0$, and every link
containing $S$ is in the last block; by Lemma~\ref{lem:block-from}, $E^\beta$ goes from $R$ to $S$.
\end{proof}

\formal{\lean{HPC.Chain.exists_blockChain} (in \lean{HPC/PseudoArc/T13Step.lean}).}

\paragraph{Plan of the proof.}
The whole proof is about the integer sequences $\lo$ and $\hi$. Section~\ref{sec:fold-readings} records their
properties: they move by at most one, and $\lo$ is \emph{crooked} in a sense that only involves integers. Sections
\ref{sec:fold-folds}--\ref{sec:fold-zigzag} prove the central result, the \emph{zigzag theorem}: from any base link, a
crooked integer sequence can be ``folded'' onto a walk that goes back and forth, with every fibre of the folding map
inside three consecutive levels. This already gives the version of the theorem with only one prescribed block (Bing's
Theorem~3). To prescribe both the first and the last block, Section~\ref{sec:fold-fake} folds an initial segment of $E$
from $r$ and a final segment from $s$, and Section~\ref{sec:fold-splice} splices the two folds together with a grouping
of the middle part by levels. Section~\ref{sec:fold-proof} finishes the proof.

\section{Readings and crooked sequences}\label{sec:fold-readings}

\begin{lemma}\label{lem:readings}
For all $t$ (and $t+1$ when it is an index):
\begin{enumerate}[label=(\alph*)]
\item $\lo(t)\le\hi(t)\le\lo(t)+1$, and $E_t\subseteq D_j$ exactly for $j\in\set{\lo(t),\hi(t)}$;
\item $\abs{\lo(t+1)-\lo(t)}\le1$, $\abs{\hi(t+1)-\hi(t)}\le1$, $\abs{\hi(t+1)-\lo(t)}\le1$ and $\abs{\hi(t)-\lo(t+1)}\le 1$;
\item if $E_t$ meets $D_j$, then $\abs{\lo(t)-j}\le1$ and $\abs{\hi(t)-j}\le1$;
\item if $\lo(t)\le j\le\hi(t)$, then $E_t\subseteq D_j$.
\end{enumerate}
\end{lemma}

\begin{proof}
(a), (d) and the first two parts of (b) are Lemma~\ref{lem:minidx}. For the others: $E_t$ and $E_{t+1}$ meet at a point
$z$, which lies in $D_{\hi(t+1)}$ and in $D_{\lo(t)}$, so these links meet; likewise $D_{\hi(t)}$ and $D_{\lo(t+1)}$. (c)
$E_t$ lies in $D_{\lo(t)}$ and in $D_{\hi(t)}$ and meets $D_j$.
\end{proof}

\formal{\lean{HPC.Fold.IsCrookedIvF} and \lean{HPC.Fold.isCrookedIvF_minIdx_maxIdx} package these facts with
crookedness; see also \lean{HPC.Fold.B79.seq_facts}.}

\begin{definition}[Crooked sequence]\label{def:crooked-seq}
Let $I=[\ell,h]$ be an interval of integers. A function $f\colon I\to\Z$ with unit steps is \emph{crooked} if for all
$x<y$ in $I$:
\begin{enumerate}[label=(\roman*)]
\item if $f(y)\ge f(x)+3$, there are $x<u<u'<y$ with $f(y)-2\le f(u)\le f(y)-1$ and $f(x)\le f(u')\le f(x)+1$;
\item if $f(y)\le f(x)-3$, there are $x<u<u'<y$ with $f(y)\le f(u)\le f(y)+1$ and $f(x)-2\le f(u')\le f(x)-1$.
\end{enumerate}
\end{definition}

In words: between two places whose values differ by three or more, the sequence first comes close to the later value
and then close to the earlier one. (The asymmetry between (i) and (ii) reflects the fact that $\lo$ is the
\emph{least} index.)

\formal{\lean{HPC.Fold.Z77.CrkOn}.}

\begin{lemma}\label{lem:crooked-seq}
\begin{enumerate}[label=(\alph*)]
\item $\lo$ is crooked on $[0,n-1]$.
\item $\bar\lo:=k-1-\hi$ is crooked on $[0,n-1]$; it is the least-index reading of $E$ in the reversed chain $\overline D$.
\item If $f$ is crooked on $[\ell,h]$, then $x\mapsto f(-x)$ is crooked on $[-h,-\ell]$, and $x\mapsto f(x+c)$ is crooked on
$[\ell-c,h-c]$.
\end{enumerate}
\end{lemma}

\begin{proof}
(a) Let $x<y$ with $\lo(y)\ge\lo(x)+3$. The link $E_x$ meets $D_{\lo(x)}$ and $E_y$ meets $D_{\lo(y)}$, so crookedness of $E$ in $D$
gives $u,u'$ with $\mathrm{Btw}(x,u,u',y)$, that is $x<u<u'<y$, and $E_u\subseteq D_{\lo(y)-1}$, $E_{u'}\subseteq D_{\lo(x)+1}$.
By Lemma~\ref{lem:readings}(a), $\lo(y)-1\in\set{\lo(u),\lo(u)+1}$ and $\lo(x)+1\in\set{\lo(u'),\lo(u')+1}$, which is~(i).
If $\lo(y)\le\lo(x)-3$, crookedness applied to the pair $E_y$ (meeting $D_{\lo(y)}$) and $E_x$ (meeting $D_{\lo(x)}$) gives
$\mathrm{Btw}(y,v,v',x)$, i.e.\ $x<v'<v<y$, with $E_v\subseteq D_{\lo(x)-1}$ and $E_{v'}\subseteq D_{\lo(y)+1}$. Take $u=v'$ and
$u'=v$: then $\lo(u)\in\set{\lo(y),\lo(y)+1}$ and $\lo(u')\in\set{\lo(x)-2,\lo(x)-1}$, which is~(ii).

(b) The link $E_t$ lies in $\overline D_j=D_{k-1-j}$ exactly for $k-1-j\in\set{\lo(t),\hi(t)}$, so its least index in $\overline D$
is $k-1-\hi(t)$. $E$ is crooked in $\overline D$ (Lemma~\ref{lem:reverse}), so (a) applies.

(c) Translation is clear. For reflection let $g(x)=f(-x)$ and $x<y$. If $g(y)\ge g(x)+3$, then $f(-y)\ge f(-x)+3$ with $-y<-x$,
so (ii) for $f$ gives $-y<v<v'<-x$ with $f(-x)\le f(v)\le f(-x)+1$ and $f(-y)-2\le f(v')\le f(-y)-1$. Then $u=-v'$, $u'=-v$
satisfy $x<u<u'<y$, $g(y)-2\le g(u)\le g(y)-1$ and $g(x)\le g(u')\le g(x)+1$. The other case is the same with (i) for $f$.
\end{proof}

\formal{\lean{HPC.Fold.Z77.crkOn_of_isCrookedIvF}, \lean{HPC.Fold.Z77.CrkOn.reverse},
\lean{HPC.Fold.B79.minIdx_reverse}.}

\section{Turning levels and folds}\label{sec:fold-folds}

\begin{definition}[Folds]\label{def:fold}
Let $g\colon[0,m]\to\Z$ have unit steps (an \emph{arm}). A \emph{fold} of $g$ onto a function $c\colon\N\to\Z$ (the \emph{walk}) is a
function $\beta\colon[0,m]\to\N$ with $\beta(0)=0$, unit steps, and $\abs{g(j)-c(\beta(j))}\le1$ for all $j$.
\end{definition}

The walks we use are determined by a list of \emph{turning levels} $T_0,T_1,\dots,T_K$ ($K\ge1$) with
$\abs{T_{i+1}-T_i}\ge3$. The \emph{leg} $i$ goes from $T_i$ to $T_{i+1}$. Its \emph{envelope} is the interval
$\env_i=[\min(T_i,T_{i+1})-1,\max(T_i,T_{i+1})+1]$, and its \emph{arrival level} is $\arr_i=T_{i+1}-1$ if $T_i<T_{i+1}$ and
$\arr_i=T_{i+1}+1$ otherwise: the level one short of the end of the leg.

\begin{definition}[Fitting]\label{def:fit}
An arm $g$ \emph{fits the turning levels $T_0,\dots,T_K$ from time $\tau$} if $K\ge1$, $\abs{T_1-T_0}\ge3$, and either
\begin{enumerate}[label=(F\arabic*)]
\item $g(j)\in\env_0$ for all $j\ge\tau$; or
\item there is $\tau'\ge\tau$ with $g(\tau')=\arr_0$, $g(j)\in\env_0$ for $\tau\le j\le\tau'$, and $g$ fits $T_1,\dots,T_K$ from time $\tau'$.
\end{enumerate}
\end{definition}

So the arm runs along the legs: on each leg it stays within one of the levels between the two turning levels, until it
reaches the level just before the end of the leg, where it passes to the next leg; on the last leg it reaches, it stays
for good.

\formal{\lean{HPC.Fold.Z77.Env}, \lean{HPC.Fold.Z77.arr}, \lean{HPC.Fold.Z77.Fit}.}

\paragraph{The jump walk.} Let $\ell_i=\abs{T_{i+1}-T_i}-2\ge1$ and $O_i=\ell_0+\dots+\ell_{i-1}$. Define $c$ on $[0,O_K]$ by
$c(O_i)=T_i$ for $i\le K$ and, for $0<u<\ell_i$,
\[
c(O_i+u)=T_i+2+u\ \text{ if }T_i<T_{i+1},\qquad c(O_i+u)=T_i-2-u\ \text{ if }T_i>T_{i+1}
\]
(and $c(x)=T_K$ for $x>O_K$). So leg $i$ is run as $T_i,T_i\pm3,T_i\pm4,\dots,T_{i+1}\mp1$, and then the next leg starts at
$T_{i+1}$: the walk skips the two levels next to each turning level. For a level $y$ define the \emph{position on leg $i$}
\[
\pi_i(y)=\begin{cases}\min\bigl(\max(y-T_i-1,0),\,\ell_i\bigr)&\text{if }T_i<T_{i+1},\\
\min\bigl(\max(T_i-1-y,0),\,\ell_i\bigr)&\text{if }T_i>T_{i+1}.\end{cases}
\]

\begin{lemma}\label{lem:legpos}
\begin{enumerate}[label=(\alph*)]
\item $0\le\pi_i(y)\le\ell_i$, and $\abs{\pi_i(y)-\pi_i(y')}\le1$ when $\abs{y-y'}\le1$.
\item If $y\in\env_i$, then $\abs{y-c(O_i+\pi_i(y))}\le1$.
\item $\pi_i(y)=0$ if $\abs{y-T_i}\le1$, and $\pi_i(\arr_i)=\ell_i$.
\end{enumerate}
\end{lemma}

\begin{proof}
Suppose $T_i<T_{i+1}$ (the other case is symmetric); write $a=T_i$, $b=T_{i+1}$, so $\ell_i=b-a-2$ and $\env_i=[a-1,b+1]$.
(a) and (c) are immediate from the formula ($\arr_i=b-1$ gives $\pi_i=\min(b-a-2,\ell_i)=\ell_i$). For (b): if $y\le a+1$,
then $\pi_i(y)=0$, $c(O_i)=a$, and $\abs{y-a}\le1$. If $a+2\le y\le b-2$, then $\pi_i(y)=y-a-1\in(0,\ell_i)$ and
$c(O_i+\pi_i(y))=y+1$. If $y\ge b-1$, then $\pi_i(y)=\ell_i$, $c(O_{i+1})=b$, and $\abs{y-b}\le1$ as $y\le b+1$.
\end{proof}

\formal{\lean{HPC.Fold.Z77.jwalk}, \lean{HPC.Fold.Z77.legPos}, \lean{HPC.Fold.Z77.jwalk_legPos},
\lean{HPC.Fold.Z77.legPos_lip}, \lean{HPC.Fold.Z77.legPos_eq_zero}.}

\begin{lemma}[Folds from fitting]\label{lem:fold-of-fit}
Let the arm $g$ fit $T_0,\dots,T_K$ from time $\tau$, with $\abs{g(\tau)-T_0}\le1$. Then there is $\beta\colon[\tau,m]\to[0,O_K]$ with
$\beta(\tau)=0$, unit steps, and $\abs{g(j)-c(\beta(j))}\le1$ for all $j\in[\tau,m]$, where $c$ is the jump walk of $T_0,\dots,T_K$.
Moreover $\beta(j)=0$ whenever $\abs{g(i)-T_0}\le1$ for all $i\in[\tau,j]$.
\end{lemma}

\begin{proof}
Induction on $K$. In case~(F1), let $\beta(j)=\pi_0(g(j))$; by Lemma~\ref{lem:legpos} it starts at $0$, has unit steps, and
satisfies the estimate. In case~(F2), $\abs{g(\tau')-T_1}=\abs{\arr_0-T_1}=1$, so by induction (applied to $T_1,\dots,T_K$, whose
jump walk is $u\mapsto c(\ell_0+u)$) there is $\beta'$ on $[\tau',m]$. Put $\beta(j)=\pi_0(g(j))$ for $j\le\tau'$ and
$\beta(j)=\ell_0+\beta'(j)$ for $j>\tau'$. At $\tau'$ both formulas give $\ell_0$ (as $\pi_0(\arr_0)=\ell_0$ and $\beta'(\tau')=0$),
so $\beta$ has unit steps; the estimate holds on both parts.

For the last statement: $\arr_0$ differs from $T_0$ by at least $2$, so if $g$ stays within one of $T_0$ on $[\tau,j]$, then in
case~(F2) $\tau'>j$, and in both cases $\beta(j)=\pi_0(g(j))=0$ by Lemma~\ref{lem:legpos}(c).
\end{proof}

\formal{\lean{HPC.Fold.Z77.fold_of_fit}, \lean{HPC.Fold.Z77.fold_of_fit_zero}.}

The point is that the walk depends only on the turning levels. If both arms of a sequence fit the same turning levels, their
folds are folds onto the same walk and can be put together.

\section{One phase of the zigzag}\label{sec:fold-phase}

Fix a sequence $S$ on $[-p,q]$ ($p,q\ge0$) with unit steps and bounds $L\le S\le H$. Its two \emph{arms} are $A_+(j)=S(j)$ for
$0\le j\le q$ and $A_-(j)=S(-j)$ for $0\le j\le p$; we write $m_+=q$, $m_-=p$, and $e$ for one of the two signs. Arm times
are $j\ge0$.

\begin{definition}\label{def:phase-state}
An \emph{up state} consists of integers $b,a$ with $b+2\le a$ (the current \emph{bottom} and the previous \emph{top}) and, for
each arm $e$, either nothing (the arm is \emph{finished}) or an \emph{entry time} $\tau_e\le m_e$ with $\abs{A_e(\tau_e)-b}\le1$.
It is \emph{not stuck} if no unfinished arm $e$ has a time $d\ge\tau_e$ with $A_e(d)=b-2$ and $A_e(j)\le a-1$ for $\tau_e\le
j<d$. \emph{Down states} are defined symmetrically: a top $a$, a previous bottom $b$, entry times with $\abs{A_e(\sigma_e)-a}\le1$;
not stuck means that no unfinished arm rises to $a+2$ while staying $\ge b+1$ from its entry time on.
\end{definition}

\begin{lemma}[Up phase]\label{lem:up-phase}
Let an up state be not stuck. Then there is $t\ge a+1$ such that every unfinished arm $e$ either
\begin{itemize}
\item \emph{finishes}: $b-1\le A_e(j)\le t-2$ for all $j\in[\tau_e,m_e]$; or
\item \emph{passes}: there is $\sigma_e>\tau_e$ with $A_e(\sigma_e)=t-1$ and $b-1\le A_e(j)\le t-2$ for $\tau_e\le j<\sigma_e$.
\end{itemize}
Moreover, either every arm finishes, or $t\le H+1$ and some arm $e$ passes and then drops: there is $d>\sigma_e$ with $A_e(d)=b-2$
and $b-1\le A_e(j)\le t-1$ for $\tau_e\le j<d$.
\end{lemma}

\begin{proof}
For each unfinished arm $e$ define a number $v_e$. If $A_e$ goes below $b-1$ at some time after $\tau_e$, let $d_e$ be the first
such time; then $d_e>\tau_e$ (as $A_e(\tau_e)\ge b-1$), $A_e(d_e)=b-2$ by unit steps, and $A_e\ge b-1$ on $[\tau_e,d_e)$. Let $M_e$
be the maximum of $A_e$ on $[\tau_e,d_e)$ and $v_e=M_e+1$. Since the state is not stuck, $M_e\ge a$, so $v_e\ge a+1$; also $v_e\le
H+1$. If $A_e$ never goes below $b-1$, or $e$ is finished, let $v_e=\max(a+1,H+2)$. Let $t=\min(v_+,v_-)\ge a+1$.

Fix an unfinished arm $e$. \emph{Every drop of $A_e$ below $b-1$ after $\tau_e$ is preceded by a visit to $[t-1,\infty)$}: if
$A_e(j)<b-1$, then $d_e\le j$ exists, and $A_e$ takes the value $M_e=v_e-1\ge t-1$ on $[\tau_e,d_e)$. Now if $A_e$ reaches $t-1$
or more after $\tau_e$, let $\sigma_e$ be the first such time. Before $\sigma_e$, $A_e\le t-2$, and $A_e\ge b-1$ (a drop would be
preceded by a visit to $[t-1,\infty)$ before $\sigma_e$); as $A_e(\tau_e)\le b+1\le t-2$, $\sigma_e>\tau_e$ and $A_e(\sigma_e)=t-1$ by
unit steps. So $e$ passes. If $A_e$ never reaches $t-1$, it never drops below $b-1$ either, and it finishes.

If $t=\max(a+1,H+2)$, then $t-1>H$ and no arm can reach $t-1$, so all arms finish. Otherwise $t=v_e<H+2$ for some arm $e$ which
drops at $d_e$, so $t\le H+1$; its maximum $t-1$ on $[\tau_e,d_e)$ is attained, so $e$ passes at a time $\sigma_e<d_e$, and
$b-1\le A_e\le t-1$ on $[\tau_e,d_e)$.
\end{proof}

\formal{\lean{HPC.Fold.Z77.up_step} (with \lean{HPC.Fold.Z77.transfer} and \lean{HPC.Fold.Z77.exists_armVal}).}

\begin{lemma}[Down phase]\label{lem:down-phase}
Let a down state (top $a$, previous bottom $b$) be not stuck. Then there is $t'\le b-1$ such that every unfinished arm either finishes
($t'+2\le A_e\le a+1$ from $\sigma_e$ on) or passes (at a time $\sigma'_e>\sigma_e$ with $A_e(\sigma'_e)=t'+1$, and $t'+2\le A_e\le a+1$
on $[\sigma_e,\sigma'_e)$). Moreover, either every arm finishes, or $t'\ge L-1$ and some arm passes and then rises: there is
$d>\sigma'_e$ with $A_e(d)=a+2$ and $t'+1\le A_e\le a+1$ on $[\sigma_e,d)$.
\end{lemma}

\begin{proof}
Apply Lemma~\ref{lem:up-phase} to the arms $-A_e$, with bottom $-a$, previous top $-b$, upper bound $-L$, and put $t'=-t$.
\end{proof}

\formal{\lean{HPC.Fold.Z77.down_step}.}

Note that these two lemmas do not use crookedness.

\section{Crookedness across the base link}\label{sec:fold-cores}

From now on $S$ is crooked on $[-p,q]$. For an arm $e$ let $\bar e$ be the other arm. Reading $S$ from the end of arm $e$ to the
end of arm $\bar e$ gives the sequence $x\mapsto A_e(-x)$ for $x\le0$, $x\mapsto A_{\bar e}(x)$ for $x\ge0$, which is $S$ or its
reflection, hence crooked (Lemma~\ref{lem:crooked-seq}(c)). In the next two lemmas, the arm called $\mathcal L$ (left) is read at
negative positions and the arm $\mathcal R$ (right) at positive ones; they apply with $\set{\mathcal L,\mathcal R}=\set{A_+,A_-}$ in
either order.

\begin{lemma}[Up core]\label{lem:up-core}
Let $b\le a$. There are no times $\tau\le d_1$ on $\mathcal L$ and $d_2$ on $\mathcal R$ such that $\mathcal L(j)\ge b$ for $j<\tau$,
$\mathcal L(j)\le a-1$ for $\tau\le j<d_1$, $\mathcal L(d_1)=b-2$, $\mathcal R(j)\ge b$ for $j<d_2$, and $\mathcal R(d_2)=a+2$.
\end{lemma}

\begin{proof}
Suppose there are. The positions $x=-d_1<y=d_2$ have values $b-2$ and $a+2$, which differ by at least $4$. By
Definition~\ref{def:crooked-seq}(i) there are $x<u<u'<y$ with value at $u$ in $[a,a+1]$ and value at $u'$ in $[b-2,b-1]$. If $u'\ge0$,
then $u'<d_2$ and $\mathcal R(u')\ge b$, a contradiction. So $u<u'<0$. The arm time $j'=-u'$ has $\mathcal L(j')\le b-1$, so $j'\ge\tau$;
and $j=-u$ satisfies $j'<j<d_1$, so $\mathcal L(j)\le a-1$, contradicting $\mathcal L(j)\ge a$.
\end{proof}

\formal{\lean{HPC.Fold.Z77.up_core}.}

\begin{lemma}[Down core]\label{lem:down-core}
Let $b\le a$. There are no times $\tau\le\sigma\le\rho_1$ on $\mathcal L$ and $\tau'\le d$ on $\mathcal R$ such that
\[
\mathcal L\ge b\text{ on }[0,\tau),\quad \mathcal L\le a-2\text{ on }[\tau,\sigma),\quad \mathcal L\ge b+1\text{ on }[\sigma,\rho_1),\quad
\mathcal L(\rho_1)=a+2,
\]
\[
\mathcal R\ge b\text{ on }[0,\tau'),\quad\mathcal R\le a-1\text{ on }[\tau',d),\quad\mathcal R(d)=b-2 .
\]
\end{lemma}

\begin{proof}
Suppose there are. The positions $x=-\rho_1<y=d$ have values $a+2$ and $b-2$. By Definition~\ref{def:crooked-seq}(ii) there are
$x<u<u'<y$ with value at $u$ in $[b-2,b-1]$ and value at $u'$ in $[a,a+1]$. If $u\ge0$: $\mathcal R(u)\le b-1$ forces $u\ge\tau'$, and
then $u<u'<d$ gives $\mathcal R(u')\le a-1$, a contradiction. So $u<0$, and the arm time $j=-u$ satisfies $\mathcal L(j)\le b-1$; hence
$j\ge\tau$ and $j<\sigma$.

Apply Definition~\ref{def:crooked-seq}(ii) again, to the positions $-\rho_1<u$, whose values $a+2$ and $\mathcal L(j)\le b-1$ differ by at
least $3$: there are $-\rho_1<v<v'<u$ with value at $v$ in $[\mathcal L(j),\mathcal L(j)+1]$, so at most $b$, and value at $v'$ in
$[a,a+1]$. The arm time $-v'$ lies in $(j,\rho_1)$. If $-v'<\sigma$, then $-v'\in[\tau,\sigma)$ and $\mathcal L(-v')\le a-2$, a
contradiction; so $-v'\ge\sigma$, and $-v\in(-v',\rho_1)\subseteq[\sigma,\rho_1)$, whence $\mathcal L(-v)\ge b+1$, a contradiction.
\end{proof}

\formal{\lean{HPC.Fold.Z77.down_core}.}

We now show that the phases of the zigzag never get stuck. An up state has the \emph{history property} if every unfinished arm
satisfies $A_e(j)\ge b$ for $j<\tau_e$.

\begin{lemma}[Up states are not stuck]\label{lem:nofail-up}
Let an up state (bottom $b$, previous top $a$, $b\le a$) have the history property, and suppose some arm $e_0$ and time $d$ satisfy
$A_{e_0}(d)=a+2$, $A_{e_0}(j)\ge b$ for $j<d$, and $\tau_{e_0}\le d$ if $e_0$ is unfinished. Then the state is not stuck.
\end{lemma}

\begin{proof}
Suppose an unfinished arm $e$ has $d_1\ge\tau_e$ with $A_e(d_1)=b-2$ and $A_e\le a-1$ on $[\tau_e,d_1)$. If $e=e_0$: since
$A_e(j)\ge b$ for $j<d$, we have $d<d_1$, so $d\in[\tau_e,d_1)$ and $A_e(d)\le a-1$, contradicting $A_e(d)=a+2$. If $e\ne e_0$:
Lemma~\ref{lem:up-core} with $\mathcal L=A_e$, $\tau=\tau_e$ and $\mathcal R=A_{e_0}$, $d_2=d$ gives a contradiction.
\end{proof}

\formal{\lean{HPC.Fold.Z77.nofail_up}.}

\begin{lemma}[Down states after an up phase are not stuck]\label{lem:nofail-down}
Let an up state (bottom $b$, previous top) have the history property, and let the up phase of Lemma~\ref{lem:up-phase} give the top
$t$, with not every arm finishing. Then the down state with top $t$, previous bottom $b$, and entry times the passing times $\sigma_e$,
is not stuck.
\end{lemma}

\begin{proof}
Let $e_0$ be the arm of Lemma~\ref{lem:up-phase} which passes and then drops: $\tau_{e_0}<\sigma_{e_0}<d$, $A_{e_0}(d)=b-2$, $A_{e_0}\le
t-1$ on $[\tau_{e_0},d)$. Suppose an arm $e$ (which passed) rises to $t+2$ at a time $\rho_1\ge\sigma_e$ while $A_e\ge b+1$ on
$[\sigma_e,\rho_1)$. If $e=e_0$: $\rho_1\ge d$ since $A_{e}\le t-1$ on $[\sigma_e,d)$; but then $d\in[\sigma_e,\rho_1)$ and $A_e(d)=b-2<b+1$.
If $e\ne e_0$: Lemma~\ref{lem:down-core} applies with $\mathcal L=A_e$ (times $\tau_e\le\sigma_e\le\rho_1$; the bound $A_e\le t-2$ on
$[\tau_e,\sigma_e)$ is part of ``passes'', and $A_e\ge b$ before $\tau_e$ is the history property) and $\mathcal R=A_{e_0}$ (times
$\tau_{e_0}\le d$), with $a=t$.
\end{proof}

\formal{\lean{HPC.Fold.Z77.nofail_down}.}

\begin{lemma}[First direction]\label{lem:first-direction}
Let $P=S(0)$. At least one of the following holds:
\begin{enumerate}
\item[$(\mathrm{OK}^-)$] there is no position $x$ with $S(x)=P-3$ and $S\le P$ at all positions between $0$ and $x$;
\item[$(\mathrm{OK}^+)$] there is no position $x$ with $S(x)=P+3$ and $S\ge P$ at all positions between $0$ and $x$.
\end{enumerate}
\end{lemma}

\begin{proof}
Suppose $x_1$ violates the first and $x_2$ the second. If they lie on the same arm, the one nearer to $0$ lies between $0$ and the
other, which is impossible ($S(x_1)=P-3<P$ and $S(x_2)=P+3>P$). Otherwise apply Lemma~\ref{lem:up-core} with $\mathcal L$ the arm of
$x_1$, $\mathcal R$ the arm of $x_2$, $\tau=0$, $a=P+1$, $b=P-1$.
\end{proof}

\formal{\lean{HPC.Fold.Z77.first_direction}; the two conditions are \lean{HPC.Fold.Z77.OKBelow} and \lean{HPC.Fold.Z77.OKAbove}.}

\section{The zigzag theorem}\label{sec:fold-zigzag}

\begin{proposition}[Main induction]\label{prop:main-up}
Let an up state with bottom $b$ and previous top $a$ satisfy $a\le H+1$, $b\ge L-1$, the history property, and be not stuck. Then
there are turning levels $b=T_0,T_1,\dots,T_K$ which every unfinished arm $e$ fits from its entry time $\tau_e$.
\end{proposition}

\begin{proof}
Induction on $(H+1-a)+(b-L+1)\ge0$. Run the up phase (Lemma~\ref{lem:up-phase}): it gives $t\ge a+1\ge b+3$.

\emph{Leg up.} Suppose the passing arms fit some list $t=T'_0,\dots,T'_{K'}$ from their passing times (or that there are no passing
arms and the list is just $t$). Then every unfinished arm fits $b,t,T'_1,\dots$ from $\tau_e$: the envelope of the leg $b\to t$ is
$[b-1,t+1]$ and its arrival level is $t-1$; a finishing arm stays in $[b-1,t-2]$ for good, which is~(F1); a passing arm stays in
$[b-1,t-1]$ on $[\tau_e,\sigma_e]$ and has $A_e(\sigma_e)=t-1$, which is~(F2).

If every arm finishes, the list $b,t$ works. Otherwise, by Lemma~\ref{lem:nofail-down}, the down state with top $t$ and previous
bottom $b$ is not stuck. Run the down phase (Lemma~\ref{lem:down-phase}); it gives $t'\le b-1\le t-4$. The envelope of the leg
$t\to t'$ is $[t'-1,t+1]$ and its arrival level is $t'+1$. An arm finishing in the down phase stays in $[t'+2,t+1]$ for good,
which is~(F1) for this leg; an arm passing it stays in $[t'+1,t+1]$ up to its passing time $\sigma'_e$ and reaches $t'+1$ there. So
if every arm finishes the down phase, the turning levels $b,t,t'$ work; and in general it suffices to find a list starting with $t'$
which the arms passing the down phase fit from their passing times $\sigma'_e$, and to prepend $b,t$.

Otherwise, consider the up state with bottom $t'$, previous top $t$ and entry times $\sigma'_e$. It satisfies $|A_e(\sigma'_e)-t'|\le1$,
$t'+2\le t$, $t\le H+1$ and $t'\ge L-1$ (both from the ``moreover'' parts of the phase lemmas). It has the history property: before
$\sigma'_e$, an arm was $\ge b$ before $\tau_e$ (history), $\ge b-1$ on $[\tau_e,\sigma_e)$, and $\ge t'+2$ on $[\sigma_e,\sigma'_e)$;
all these are $\ge t'$. It is not stuck by Lemma~\ref{lem:nofail-up}, applied with the arm $e_0$ which passes the down phase and then
rises to $t+2$ at time $d$: before $d$ it is $\ge t'$ (by the same history count, and $\ge t'+1$ on $[\sigma_{e_0},d)$), and
$\sigma'_{e_0}\le d$. The measure has decreased by at least $2$, since $t\ge a+1$ and $t'\le b-1$. By induction there is a list
$t',\dots$ which the arms fit from $\sigma'_e$.
\end{proof}

\formal{\lean{HPC.Fold.Z77.main_up}.}

\begin{theorem}[Zigzag theorem]\label{thm:zigzag}
Let $S$ be crooked on $[-p,q]$ with $L\le S\le H$, and let $\epsilon\in\set{+1,-1}$ be such that $(\mathrm{OK}^\epsilon)$ of
Lemma~\ref{lem:first-direction} holds. Then there are a walk $c$ with $c(0)=S(0)+\epsilon$ and a function $\beta\colon[-p,q]\to\N$ with
unit steps and $\beta(0)=0$ such that $\abs{S(x)-c(\beta(x))}\le1$ for all $x$. Moreover $\beta(x)=0$ whenever $\abs{S(y)-c(0)}\le1$ for
all $y$ between $0$ and $x$.
\end{theorem}

\begin{proof}
Let $P=S(0)$; both arms start at $P$. It suffices to find turning levels $T_0=P+\epsilon,T_1,\dots$ which both arms fit from time
$0$: Lemma~\ref{lem:fold-of-fit} then gives folds $\beta_+$, $\beta_-$ of the two arms onto the same jump walk $c$, and
$\beta(x)=\beta_+(x)$ for $x\ge0$, $\beta(x)=\beta_-(-x)$ for $x\le0$ is as required.

If $\epsilon=-1$: the up state with bottom $P-1$, previous top $P+1$, and both entry times $0$, has the history property (vacuously),
$P+1\le H+1$, $P-1\ge L-1$, and it is not stuck by $(\mathrm{OK}^-)$ (an arm dropping to $(P-1)-2=P-3$ while staying $\le P$).
Proposition~\ref{prop:main-up} gives the turning levels, starting at $P-1$.

If $\epsilon=+1$: the down state with top $P+1$ and previous bottom $P-1$ is not stuck by $(\mathrm{OK}^+)$. Run the down phase; it
gives $t'\le P-2$. As in the proof of Proposition~\ref{prop:main-up}, either all arms finish (and the turning levels $P+1,t'$ work) or the
up state with bottom $t'$, previous top $P+1$, and the passing times as entry times, satisfies the hypotheses of
Proposition~\ref{prop:main-up} (history: before passing, the arms are $\ge t'+2$; not stuck by Lemma~\ref{lem:nofail-up}). Prepend $P+1$
to the turning levels it gives.
\end{proof}

\formal{\lean{HPC.Fold.Z77.greedy_fit_below}, \lean{HPC.Fold.Z77.greedy_fit_above}, and \lean{HPC.Fold.Z77.greedy_fit}
(which uses Lemma~\ref{lem:first-direction} to choose $\epsilon$).}

\begin{corollary}[Bing's Theorem~3, three-link form]\label{cor:theorem3}
Let $E$ be crooked in $D$ and let $i$ be an index of $E$. Then there is a block map for $E$ with small blocks in which $E_i$ lies in the
first block.
\end{corollary}

\begin{proof}
Apply Theorem~\ref{thm:zigzag} to $S(x)=\lo(i+x)$ on $[-i,n-1-i]$ (crooked by Lemma~\ref{lem:crooked-seq}, with $0\le S\le k-1$), with
$\epsilon$ from Lemma~\ref{lem:first-direction}. The function $t\mapsto\beta(t-i)$ is a schedule (Definition~\ref{def:schedule}) with value
$0$ at $i$, hence a block map (Lemma~\ref{lem:schedule}). If $\beta(t-i)=u$ then $\abs{\lo(t)-c(u)}\le1$, so $E_t\subseteq
D_{\lo(t)}\subseteq D[c(u)-1,c(u)+1]$.
\end{proof}

\formal{\lean{HPC.Fold.theorem3}, via \lean{HPC.Fold.Z77.s3_zig} and \lean{HPC.exists_blocks_of_walk}.}

\section{Folds with a fake end link}\label{sec:fold-fake}

To prescribe the last block as well as the first, we fold only an initial segment $[0,a]$ of $E$, where $E_a$ meets the first
link $D_0$, and we need to know that the walk comes back down to level $\le1$ after the end of the segment. This is arranged by
appending a \emph{fake link} at level $0$.

In this section $f\colon[0,n-1]\to\Z$ is crooked (Definition~\ref{def:crooked-seq}) with unit steps and $0\le f\le H$, and $i\le a$
are indices. For a \emph{fake level} $F$ consider the two \emph{events}
\begin{itemize}
\item $\mathrm{Ev}^+_F(f,i,a)$: either there is $t\le a$ with $f(t)=f(i)+3$ and $f\ge f(i)$ at all indices strictly between $i$ and
$t$; or $F=f(i)+3$ and $f\ge f(i)$ on $[i,a]$;
\item $\mathrm{Ev}^-_F(f,i,a)$: the same with $f(i)-3$ in place of $f(i)+3$ and $f\le f(i)$ in place of $f\ge f(i)$.
\end{itemize}
So $\mathrm{Ev}^{\pm}_F(f,i,a)$ says that the sequence $f$ on $[0,a]$, continued by one more value $F$ at $a+1$, violates the condition
$(\mathrm{OK}^\pm)$ of Lemma~\ref{lem:first-direction} at the base $i$.

\formal{\lean{HPC.Fold.B79.EvUp} and \lean{HPC.Fold.B79.EvDown}; reflecting $f$ exchanges them
(\lean{HPC.Fold.B79.evUp_reflect}).}

\begin{lemma}[Crooked extensions]\label{lem:crooked-ext}
\begin{enumerate}[label=(\alph*)]
\item Suppose $f(a)\le1$ and the \emph{low meeting property} holds: for every $x<a$ with $f(x)\ge3$ there are $x<y<y'<a$ with
$f(y)\le1$ and $f(x)-2\le f(y')\le f(x)-1$. Then the sequence $f^\flat$ on $[0,a+1]$ equal to $f$ on $[0,a]$ with $f^\flat(a+1)=0$ is
crooked, has unit steps, and satisfies $0\le f^\flat\le H$.
\item Suppose $f(a)\ge H-1$ and the \emph{high meeting property} holds: for every $x<a$ with $f(x)\le H-3$ there are $x<y<y'<a$ with
$H-2\le f(y)\le H-1$ and $f(x)\le f(y')\le f(x)+1$. Then the sequence $f^\sharp$ equal to $f$ on $[0,a]$ with $f^\sharp(a+1)=H$ is
crooked, has unit steps, and satisfies $0\le f^\sharp\le H$.
\end{enumerate}
\end{lemma}

\begin{proof}
(a) The step from $a$ to $a+1$ is $f(a)-0\in\set{0,1}$, and the bounds are clear. Pairs $x<y\le a$ are covered by crookedness of
$f$. Let $y=a+1$ and $x\le a$. Clause~(i) of Definition~\ref{def:crooked-seq} needs $0\ge f(x)+3$, impossible as $f\ge0$. For
clause~(ii) let $0\le f(x)-3$; then $x\ne a$ (as $f(a)\le1$), so $x<a$, and the low meeting property gives $x<u<u'<a$ with $0\le
f(u)\le1=f^\flat(a+1)+1$ and $f(x)-2\le f(u')\le f(x)-1$, which is clause~(ii) for the pair $(x,a+1)$.

(b) The same with the roles of the two clauses exchanged: clause~(ii) for $(x,a+1)$ needs $H\le f(x)-3$, impossible; clause~(i)
needs $H\ge f(x)+3$, so $x<a$, and the high meeting property supplies the two indices.
\end{proof}

\formal{\lean{HPC.Fold.B79.crkOn_lowExt}, \lean{HPC.Fold.B79.crkOn_highExt} (on sequences re-indexed so that the base is at $0$).}

\begin{lemma}[Fold with a fake low end]\label{lem:fake-low}
Assume the hypotheses of Lemma~\ref{lem:crooked-ext}(a), and let $\epsilon\in\set{+1,-1}$ be such that $\mathrm{Ev}^\epsilon_0(f,i,a)$
fails. Then there are a walk $c\colon\N\to\Z$, a function $\beta\colon[0,a]\to\N$ and a number $U\in\N$ such that
\begin{enumerate}[label=(\alph*)]
\item $c(0)=f(i)+\epsilon$, $\beta(i)=0$, $\beta$ has unit steps, and $\abs{f(t)-c(\beta(t))}\le1$ for $t\le a$;
\item $\abs{U-\beta(a)}\le1$ and $c(U)\le1$;
\item if $t\le a$ is a neighbour of $i$ ($t=i\pm1$) with $\abs{f(t)-c(0)}\le1$, then $\beta(t)=0$;
\item if $a=i$ and $\abs{c(0)}\le1$, then $U=0$.
\end{enumerate}
\end{lemma}

\begin{proof}
Let $S(x)=f^\flat(i+x)$ for $-i\le x\le a-i+1$. By Lemma~\ref{lem:crooked-ext}(a) and Lemma~\ref{lem:crooked-seq}(c), $S$ is crooked
on $[-i,a-i+1]$, with unit steps and $0\le S\le H$; $S(0)=f(i)$. The condition $(\mathrm{OK}^\epsilon)$ of
Lemma~\ref{lem:first-direction} for $S$ is the failure of $\mathrm{Ev}^\epsilon_0(f,i,a)$: a position $x\le a-i$ violating it is an index
$t=i+x\le a$ as in the first alternative of the event, and the position $a-i+1$ (value $0$) violates it exactly when the second
alternative holds, the positions strictly between $0$ and $a-i+1$ being those of $[i,a]$. So Theorem~\ref{thm:zigzag} applies and
gives $c$ and $\tilde\beta$ on $[-i,a-i+1]$. Put $\beta(t)=\tilde\beta(t-i)$ and $U=\tilde\beta(a-i+1)$. Then (a) is the statement of
the theorem; (b) follows from unit steps and $\abs{0-c(U)}=\abs{S(a-i+1)-c(\tilde\beta(a-i+1))}\le1$. For (c), the positions between $0$
and $t-i=\pm1$ are $0$ and $\pm1$, and $\abs{S(0)-c(0)}=1$; so the ``moreover'' part of Theorem~\ref{thm:zigzag} gives
$\tilde\beta(t-i)=0$. For (d), if $a=i$ the positions between $0$ and $1$ are $0$ and $1$, with $S(1)=0$; as $\abs{c(0)}\le1$, again
$U=\tilde\beta(1)=0$.
\end{proof}

\formal{\lean{HPC.Fold.B79.finfold_low}, from \lean{HPC.Fold.B79.fold_low_end_mode} (the version on $\Z$-indexed sequences;
without the choice of $\epsilon$ it is \lean{HPC.Fold.B78.fold_low_end}). With $\epsilon$ chosen by Lemma~\ref{lem:first-direction}
it is \lean{HPC.Fold.B79.finfold_low_any}.}

\begin{lemma}[Fold with a fake high end]\label{lem:fake-high}
Assume the hypotheses of Lemma~\ref{lem:crooked-ext}(b), and let $\epsilon$ be such that $\mathrm{Ev}^\epsilon_H(f,i,a)$ fails. Then
there are $c$, $\beta$, $U$ with properties (a), (c) of Lemma~\ref{lem:fake-low} and $\abs{U-\beta(a)}\le1$, $c(U)\ge H-1$.
\end{lemma}

\begin{proof}
The same, with $f^\sharp$ and Lemma~\ref{lem:crooked-ext}(b): now $\abs{H-c(U)}\le1$.
\end{proof}

\formal{\lean{HPC.Fold.B79.finfold_high}, from \lean{HPC.Fold.B79.fold_high_end_mode}.}

\paragraph{The meeting properties.} They come from crookedness when $E_a$ meets $D_0$.

\begin{lemma}\label{lem:meet-low}
Let $E$ be crooked in $D$ (of length $k\ge2$), let $E_a$ meet $D_0$, let $t<a$, and let $E_t\subseteq D_j$ with $j\ge3$. Then there are
$t<y<y'<a$ with $E_y\subseteq D_1$ and $E_{y'}\subseteq D_{j-1}$. Consequently:
\begin{enumerate}[label=(\alph*)]
\item $\lo$ has the low meeting property at $a$, and $\lo(a)\le1$;
\item $\bar\lo=k-1-\hi$ has the high meeting property at $a$ with $H=k-1$, and $\bar\lo(a)\ge k-2$;
\item for every $t<a$ with $\lo(t)\ge3$ there are $t<y<y'<a$ with $\lo(y)\le1$ and $\hi(y')\ge\lo(t)-1$.
\end{enumerate}
\end{lemma}

\begin{proof}
$E_a$ meets $D_0$, $E_t$ meets $D_j$, and $j-0\ge3$; crookedness gives indices $r,s$ with $\mathrm{Btw}(a,r,s,t)$, $E_r\subseteq D_{j-1}$
and $E_s\subseteq D_1$. Since $t<a$, $\mathrm{Btw}(a,r,s,t)$ means $t<s<r<a$; take $y=s$, $y'=r$.

(a) Apply this with $j=\lo(t)$. Then $\lo(y)\le1$, and $D_{\lo(t)-1}\in\set{D_{\lo(y')},D_{\hi(y')}}$, so $\lo(t)-2\le\lo(y')\le\lo(t)-1$
(Lemma~\ref{lem:readings}(a)). Also $\lo(a)\le1$ and $\hi(a)\le1$ by Lemma~\ref{lem:readings}(c).
(b) Apply it with $j=\hi(t)$, which is $\ge3$ when $\bar\lo(t)\le k-4$. Then $\hi(y)\in\set{1,2}$, i.e.\ $k-3\le\bar\lo(y)\le k-2$, and
$\hi(y')\in\set{\hi(t)-1,\hi(t)}$, i.e.\ $\bar\lo(t)\le\bar\lo(y')\le\bar\lo(t)+1$. And $\hi(a)\le1$ gives $\bar\lo(a)\ge k-2$.
(c) Apply it with $j=\lo(t)$: $\hi(y')\ge\lo(t)-1$ as $E_{y'}\subseteq D_{\lo(t)-1}$.
\end{proof}

\formal{\lean{HPC.Fold.B79.meet_low_gen}, \lean{HPC.Fold.B79.maxIdx_near_zero}; the two sequences and their properties are
collected in \lean{HPC.Fold.B79.seq_facts}.}

\paragraph{Exclusion.} The prefix fold must send \emph{both} links $r,r+1$ to time $0$. When $\lo(r)\ne\lo(r+1)$ this can fail for
the fold read on $\lo$ in the only direction which would work; then we read on $\hi$ instead. The next lemma says that one of the two
readings always works.

\begin{lemma}[Exclusion]\label{lem:exclusion}
Let $E$ be crooked in $D$, let $E_a$ meet $D_0$, and let $\ell,h$ be adjacent indices (that is, $\abs{\ell-h}=1$) with $\ell\le a$ and
\[
x:=\lo(\ell),\qquad \lo(h)=\hi(h)=x+1 .
\]
Then $\mathrm{Ev}^+_0(\lo,\ell,a)$ and $\mathrm{Ev}^-_0(\hi,h,a)$ do not both hold.
\end{lemma}

\begin{proof}
Suppose both hold. The second alternative of $\mathrm{Ev}^+_0(\lo,\ell,a)$ would need $x+3=0$, impossible; so there is $t\le a$ with
$\lo(t)=x+3$ and $\lo\ge x$ strictly between $\ell$ and $t$. Note $t\ne\ell,h$ (their $\lo$ is $x$, $x+1$), and $\hi(t)\ge x+3$. We
use repeatedly that an index strictly between $\ell$ and $t$ has $\lo\ge x$, and the following consequence of adjacency of $\ell,h$:
\emph{an index $w\ne\ell,h$ lies on the same side of $\ell$ as of $h$}.

\emph{Case 1: the first alternative of $\mathrm{Ev}^-_0(\hi,h,a)$}: there is $u\le a$ with $\hi(u)=x-2$ and $\hi\le x+1$ strictly between
$h$ and $u$. Then $u\ne\ell,h$ (their $\hi$ is $\ge x$) and $\lo(u)\le x-2$.
\begin{itemize}
\item $t,u$ both below the pair ($t<\ell$, $u<h$). If $t<u$, then $u$ is strictly between $t$ and $\ell$, so $\lo(u)\ge x$, a contradiction.
If $u<t$, then $t$ is strictly between $u$ and $h$, so $\hi(t)\le x+1$, a contradiction. The case of both above is the same.
\item $t$ below, $u$ above ($t<\ell$, $u>h$). $E_t$ meets $D_{x+3}$ and $E_u$ meets $D_{x-2}$; crookedness gives $t<w<w'<u$ with
$E_w\subseteq D_{x-1}$ and $E_{w'}\subseteq D_{x+2}$. As $\lo(w)\le x-1$, $w$ is not strictly between $t$ and $\ell$ and $w\ne\ell$; so
$w>\ell$, and then $w'>\ell$. As $\hi(w')\ge x+2$, $w'\ne h$, so $w'$ is strictly between $h$ and $u$, and $\hi(w')\le x+1$, a contradiction.
\item $t$ above, $u$ below. Crookedness for $u<t$ ($E_u$ meets $D_{x-2}$, $E_t$ meets $D_{x+3}$) gives $u<w<w'<t$ with $E_w\subseteq
D_{x+2}$ and $E_{w'}\subseteq D_{x-1}$. As $\hi(w)\ge x+2$, $w$ is not strictly between $u$ and $h$ and $w\ne h$, so $w>h$ and $w'>h$; as
$\lo(w')\le x-1$, $w'\ne\ell$, so $w'$ is strictly between $\ell$ and $t$ and $\lo(w')\ge x$, a contradiction.
\end{itemize}

\emph{Case 2: the second alternative}: $x-2=0$ and $\hi\le x+1=3$ on $[h,a]$. If $t>\ell$, then $t\ge h$ (as $t\ne h$) and $t\le a$, so
$\hi(t)\le3<5=\lo(t)$, impossible. If $t<\ell$: $t<a$ and $\lo(t)=5\ge3$, so Lemma~\ref{lem:meet-low}(c) gives $t<y<y'<a$ with $\lo(y)\le1$
and $\hi(y')\ge4$. As $\lo(y)<x$, $y$ is not strictly between $t$ and $\ell$, and $y\ne\ell$; so $y>\ell$ and $y'\ge h$. Then $y'\in[h,a]$ and
$\hi(y')\le3$, a contradiction.
\end{proof}

\formal{\lean{HPC.Fold.B79.pair_exclusion}.}

\begin{proposition}[Prefix fold]\label{prop:prefix-fold}
Let $k\ge4$, let $E$ be crooked in $D$, let $r'=r+1$, and let $a\ge r$ be an index such that $E_a$ meets $D_0$. Then there are a
\emph{reading} $\lambda\colon[0,n-1]\to\Z$ with $\lo\le\lambda\le\hi$, a walk $c$, a function $\beta\colon[0,a]\to\N$ with unit steps
and a number $U$ such that
\begin{enumerate}[label=(\alph*)]
\item $\beta(r)=0$, and either $r'\le a$ and $\beta(r')=0$, or $a=r$ and $U=0$;
\item $\abs{\lambda(t)-c(\beta(t))}\le1$ for $t\le a$;
\item $\abs{U-\beta(a)}\le1$ and $c(U)\le1$.
\end{enumerate}
\end{proposition}

\begin{proof}
By Lemma~\ref{lem:meet-low}(a), Lemma~\ref{lem:fake-low} applies to $f=\lo$ (crooked by Lemma~\ref{lem:crooked-seq}, $0\le\lo\le k$)
with any base $i\le a$.

\emph{Case $a=r$.} Take $\lambda=\lo$, base $r$, and $\epsilon=-1$: the event $\mathrm{Ev}^-_0(\lo,r,a)$ fails, because $\lo(r)=\lo(a)\le1$
and neither a value $\lo(t)=\lo(r)-3<0$ nor $0=\lo(r)-3$ is possible. Then $c(0)=\lo(r)-1\in\set{-1,0}$, so $U=0$ by
Lemma~\ref{lem:fake-low}(d).

From now on $a>r$, so $r'\le a$. \emph{Case $\lo(r')=\lo(r)$.} Take $\lambda=\lo$, base $r$, and $\epsilon$ given by
Lemma~\ref{lem:first-direction} for the extended sequence (one of the two events fails). Then $\abs{\lo(r')-c(0)}=1$, so $\beta(r')=0$ by
Lemma~\ref{lem:fake-low}(c).

\emph{Case $\lo(r')\ne\lo(r)$.} Then they differ by one; let $\ell\in\set{r,r'}$ be the one with the smaller value $x=\lo(\ell)$ and $h$
the other, so $\lo(h)=x+1$. Also $\hi(h)=x+1$: $\hi(h)\ge\lo(h)$, and $\hi(h)\le\lo(\ell)+1$ by Lemma~\ref{lem:readings}(b). It suffices
to find $\lambda,c,\beta,U$ with (b), (c) and $\beta(\ell)=\beta(h)=0$.
\begin{itemize}
\item If $\mathrm{Ev}^+_0(\lo,\ell,a)$ fails, take $\lambda=\lo$, base $\ell$, $\epsilon=+1$. Then $c(0)=x+1=\lo(h)$, so $\beta(h)=0$ by
Lemma~\ref{lem:fake-low}(c).
\item Otherwise $\mathrm{Ev}^-_0(\hi,h,a)$ fails by Lemma~\ref{lem:exclusion}. Reflecting, $\mathrm{Ev}^+_{k-1}(\bar\lo,h,a)$ fails for
$\bar\lo=k-1-\hi$. By Lemma~\ref{lem:meet-low}(b), Lemma~\ref{lem:fake-high} applies to $\bar\lo$ (crooked by
Lemma~\ref{lem:crooked-seq}(b)) with $H=k-1$, base $h$ and $\epsilon=+1$, giving $\bar c,\beta,U$ with $\bar c(0)=\bar\lo(h)+1$ and
$\bar c(U)\ge k-2$. Take $\lambda=\hi$ and $c=k-1-\bar c$. Then (b) holds, $c(U)\le1$, and $c(0)=\hi(h)-1=x$; since
$\hi(\ell)\in\set{x,x+1}$, $\abs{\bar\lo(\ell)-\bar c(0)}=\abs{\hi(\ell)-x}\le1$, so $\beta(\ell)=0$ by Lemma~\ref{lem:fake-low}(c).\qedhere
\end{itemize}
\end{proof}

\formal{\lean{HPC.Fold.B79.pair_prefix_fold}; the conclusion is \lean{HPC.Fold.B79.PrefixFold}.}

\begin{corollary}[Suffix fold]\label{cor:suffix-fold}
Let $k\ge4$, let $E$ be crooked in $D$, let $s=s'+1$, and let $b\le s$ be an index such that $E_b$ meets $D_{k-1}$. Then there are a
reading $\mu$ with $\lo\le\mu\le\hi$, a walk $c'$, a function $\beta'\colon[b,n-1]\to\N$ with unit steps, and $V$, such that
$\beta'(s)=0$; either $b\le s'$ and $\beta'(s')=0$, or $b=s$ and $V=0$; $\abs{(k-1-\mu(t))-c'(\beta'(t))}\le1$ for $t\ge b$; and
$\abs{V-\beta'(b)}\le1$, $c'(V)\le1$.
\end{corollary}

\begin{proof}
Apply Proposition~\ref{prop:prefix-fold} to the reversed chains: $\overline E$ is crooked in $\overline D$ (Lemma~\ref{lem:reverse},
twice), the pair $s,s'$ becomes the pair $n-1-s$, $n-1-s'=(n-1-s)+1$, and $\overline E_{n-1-b}=E_b$ meets $\overline D_0=D_{k-1}$. The
least and greatest indices of $\overline E_{n-1-t}$ in $\overline D$ are $k-1-\hi(t)$ and $k-1-\lo(t)$, so a reading $\bar\lambda$ for the
reversed chains gives the reading $\mu(t)=k-1-\bar\lambda(n-1-t)$ of $E$ in $D$.
\end{proof}

\formal{\lean{HPC.Fold.B79.pair_suffix_fold}.}

\section{The splice}\label{sec:fold-splice}

We now combine a prefix fold on $[0,a]$ and a suffix fold on $[b,n-1]$ with the middle part of $E$, which is grouped by its levels.

\begin{proposition}[Splice]\label{prop:splice}
Let $k\ge6$, let $E$ be crooked in $D$, and let $r\le a<b\le s$ with $\lo(a)\le1$ and $\hi(b)\ge k-2$. Let
$(\lambda,c,\beta_P,U)$ satisfy (b), (c) of Proposition~\ref{prop:prefix-fold} with $\beta_P(r)=0$, and let $(\mu,c',\beta_Q,V)$ satisfy
the conclusion of Corollary~\ref{cor:suffix-fold} except the clause about $s'$. Then there is a schedule $\beta\colon[0,n-1]\to\N$ (unit
steps) such that
\begin{enumerate}[label=(\alph*)]
\item $\beta(r)=0$ and $\beta(t)\le\beta(s)$ for all $t$;
\item every block $\set{t:\beta(t)=v}$ lies in some $D[w,w+2]$: there is $w$ such that for each $t$ with $\beta(t)=v$, $\lo(t)$ or $\hi(t)$
lies in $[w,w+2]$;
\item if $t\le a$ and $\beta_P(t)=0$, then $\beta(t)=0$; if $U=0$, $t<b$ and $\lo(t)\le2$, then $\beta(t)=0$;
\item if $t\ge b$ and $\beta_Q(t)=0$, then $\beta(t)=\beta(s)$; if $V=0$, $t>a$ and $\lo(t)\ge k-3$, then $\beta(t)=\beta(s)$.
\end{enumerate}
\end{proposition}

Note that ``$\lo(t)$ or $\hi(t)$ in $[w,w+2]$'' implies $E_t\subseteq D[w,w+2]$, since $E_t\subseteq D_{\lo(t)}$ and $E_t\subseteq
D_{\hi(t)}$.

\begin{proof}
Let $G=k-5\ge1$ and $\gamma(y)=\min(\max(y-2,0),G)$; $\gamma$ is non-decreasing and moves by at most one when $y$ does. The \emph{middle
schedule} is $g(t)=\gamma(\lo(t))$ for $t<b$ and $g(t)=\gamma(\hi(t))$ for $t\ge b$. Let $K_P=\set{t\le a:\beta_P(t)<U}$ (the part of the
prefix fold kept) and $K_S=\set{t\ge b:\beta_Q(t)<V}$, and define
\[
\beta(t)=\begin{cases}\beta_P(t)&t\in K_P,\\ U+G+V-\beta_Q(t)&t\in K_S,\\ U+g(t)&\text{otherwise.}\end{cases}
\]
We first note $a+3\le b$: by Lemma~\ref{lem:readings}(b), if $b=a+1$ then $\hi(b)\le\lo(a)+1\le2$, and if $b=a+2$ then $\hi(b)\le\lo(a+1)+1
\le\lo(a)+2\le3$; both contradict $\hi(b)\ge k-2\ge4$. So $K_P$ and $K_S$ are disjoint, and far apart.

Two observations do most of the work. \emph{(P)} If $t\le a$ and $\beta_P(t)=U$, then $\abs{\lambda(t)-c(U)}\le1$ and $c(U)\le1$, so
$\lo(t)\le\lambda(t)\le2$ and $g(t)=\gamma(\lo(t))=0$. \emph{(Q)} If $t\ge b$ and $\beta_Q(t)=V$, then $k-1-\mu(t)\le c'(V)+1\le2$, so
$\hi(t)\ge\mu(t)\ge k-3$ and $g(t)=\gamma(\hi(t))=G$ (as $\hi(t)-2\ge k-5$).

\emph{Unit steps.} Let $t'=t+1$.
\begin{itemize}
\item Both in $K_P$, or both in $K_S$: unit steps of $\beta_P$, $\beta_Q$.
\item $t\in K_P$, $t'\notin K_P$ (so $t'\notin K_S$, as $t'\le a+1<b$). If $t'\le a$, then $\beta_P(t)<U\le\beta_P(t')$ forces $\beta_P(t)=U-1$,
$\beta_P(t')=U$, and by (P) $\beta(t')=U$. If $t'=a+1$, then $t=a$, $\beta_P(a)=U-1$ (from $\abs{U-\beta_P(a)}\le1$), and
$\lo(t')\le\lo(a)+1\le2$, so $\beta(t')=U$. In both cases $\beta(t)=U-1$. Entering $K_P$ ($t\notin K_P$, $t'\in K_P$) is the same.
\item $t\in K_S$, $t'\notin K_S$: then $t'\ge b$, $\beta_Q(t)=V-1$, $\beta_Q(t')=V$, and by (Q) $\beta(t')=U+G=\beta(t)-1$. Entering $K_S$ at
$t'$: if $t\ge b$ this is the same; otherwise $t'=b$, $\beta_Q(b)=V-1$ (from $\abs{V-\beta_Q(b)}\le1$), $\beta(t')=U+G+1$, and
$\lo(t)\ge\hi(b)-1\ge k-3$ by Lemma~\ref{lem:readings}(b), so $\beta(t)=U+G$.
\item Neither in $K_P\cup K_S$: $\beta=U+g$, and $g$ has unit steps since $\gamma$ is $1$-Lipschitz and $\lo$, $\hi$ have unit steps;
at the switch $t=b-1$, $t'=b$, use $\abs{\hi(b)-\lo(b-1)}\le1$.
\end{itemize}

\emph{(a)} If $U>0$, then $r\in K_P$ and $\beta(r)=\beta_P(r)=0$. If $U=0$, then $r\notin K_P\cup K_S$, and by (P) (with $\beta_P(r)=0=U$)
$\beta(r)=g(r)=0$. Clearly $\beta\le U+G+V$ everywhere ($\beta<U$ on $K_P$, $g\le G$). And $\beta(s)=U+G+V$: if $V>0$, $s\in K_S$ and
$\beta_Q(s)=0$; if $V=0$, $s\notin K_S$ and by (Q) $\beta(s)=U+G$.

\emph{(b)} Let $v<U$. Then every $t$ with $\beta(t)=v$ lies in $K_P$ (outside $K_P$, $\beta\ge U$) and has $\beta_P(t)=v$, so
$\lambda(t)\in[c(v)-1,c(v)+1]$; since $\lambda(t)$ is $\lo(t)$ or $\hi(t)$ (Lemma~\ref{lem:readings}(a)), take $w=c(v)-1$. Let $U\le
v\le U+G$ and $j=v-U$; then $\beta(t)=v$ forces $t\notin K_P\cup K_S$ and $g(t)=j$, so the relevant reading $y$ ($\lo(t)$ if $t<b$, $\hi(t)$
otherwise) satisfies $\gamma(y)=j$: if $j=0$, $0\le y\le2$ (take $w=0$); if $0<j<G$, $y=j+2$ (take $w=j$); if $j=G$, $k-3\le y\le k-1$ (take
$w=k-3$). Let $v>U+G$; then $t\in K_S$ and $\beta_Q(t)=U+G+V-v$, so $k-1-\mu(t)$ is within one of $c'(U+G+V-v)$; take $w=k-2-c'(U+G+V-v)$.

\emph{(c)} If $t\le a$ and $\beta_P(t)=0$: if $U>0$, $t\in K_P$ and $\beta(t)=0$; if $U=0$, by (P) $\beta(t)=g(t)=0$. If $U=0$, $t<b$ and
$\lo(t)\le2$: $t\notin K_P\cup K_S$ and $g(t)=\gamma(\lo(t))=0$.

\emph{(d)} If $t\ge b$ and $\beta_Q(t)=0$: if $V>0$, $t\in K_S$ and $\beta(t)=U+G+V$; if $V=0$, by (Q) $\beta(t)=U+G=\beta(s)$. If $V=0$,
$t>a$ and $\lo(t)\ge k-3$: $t\notin K_P\cup K_S$, and $g(t)=G$ since $\hi(t)\ge\lo(t)\ge k-3$.
\end{proof}

\formal{\lean{HPC.Fold.B79.splice_schedule_pair} (the version with $\lambda=\mu=\lo$ and without (c), (d) is
\lean{HPC.Fold.B78.splice_schedule}); the passage from levels to links is \lean{HPC.Fold.B78.link_sub_of_level}.}

\section{Proof of the folding theorem}\label{sec:fold-proof}

\begin{lemma}[From a schedule to the block map]\label{lem:assemble}
Let $\beta$ be a schedule for $E$ with $\beta(r)=\beta(r')=0$, $\beta(s')=\beta(s)\ge\beta(t)$ for all $t$, and such that every block of
$\beta$ lies in some $D[w,w+2]$. Then the conclusion of Theorem~\ref{thm:folding} holds.
\end{lemma}

\begin{proof}
By Lemma~\ref{lem:schedule}, $\beta$ is a block map onto $[q]$ with $q=1+\beta(s)$. Its blocks are those of $\beta$, and $\beta(s)=\beta(s')=q-1$.
\end{proof}

\formal{\lean{HPC.Fold.B79.assemble_pair}.}

\begin{lemma}[Two windows]\label{lem:windows}
Let $A,B$ be sets of indices with $A\cup B=[0,n-1]$, such that every $E_t$ with $t\in A$ lies in $D[w_A,w_A+2]$ and every $E_t$ with $t\in
B$ lies in $D[w_B,w_B+2]$. Let $r'=r+1<s'=s-1$ with $r,r'\in A$, and suppose that either $s',s\in B$, or $s',s\in A$ and some $c$ with
$r'<c<s'$ lies in $B$. Then the conclusion of Theorem~\ref{thm:folding} holds.
\end{lemma}

\begin{proof}
In the first case let $\beta(t)=1$ for $t\in\set{s',s}$ and otherwise $\beta(t)=0$ for $t\in A$, $\beta(t)=1$ for $t\notin A$. In the second
case let, for $t<c$, $\beta(t)=0$ if $t\in A$ and $1$ otherwise; $\beta(c)=1$; and, for $t>c$, $\beta(t)=2$ if $t\in A$ and $1$ otherwise.
In both cases $\beta$ has unit steps (in the second case the only way to jump from $0$ to $2$ would be across $c$, where $\beta=1$),
$\beta(r)=\beta(r')=0$, and $\beta(s)=\beta(s')$ is the maximum. Block $1$ consists of indices outside $A$, hence in $B$, together with
$s',s$, resp.\ $c$, which are in $B$; the other blocks consist of indices in $A$. Apply Lemma~\ref{lem:assemble}.
\end{proof}

\formal{\lean{HPC.Fold.B79.pair_windows}.}

\begin{proof}[Proof of Theorem~\ref{thm:folding}]
Let $h_l$ and $h_h$ be the given indices: $E_{h_l}$ meets $D_0$ and $E_{h_h}$ meets $D_{k-1}$, so $\lo(h_l)\le1$ and $\hi(h_h)\ge k-2$
(Lemma~\ref{lem:readings}(c)). Let $A=\set{t:\lo(t)\le2}$ and $B=\set{t:\hi(t)\ge k-3}$; then $E_t\subseteq D[0,2]$ for $t\in A$ and
$E_t\subseteq D[k-3,k-1]$ for $t\in B$.

\emph{$k=3$.} Every link lies in $D[0,2]=\bigcup D$; take the single block, $\beta\equiv0$, $q=1$.

\emph{$k\ge4$: then $r'<s'$.} Otherwise $s=r+2$, and $h_l,h_h\in[r,r+2]$. Crookedness applied to $E_{h_l}$ (meeting $D_0$) and $E_{h_h}$ (meeting
$D_{k-1}$, $k-1\ge3$) gives two indices strictly between $h_l$ and $h_h$, which is impossible within three consecutive indices.

\emph{$k\ge6$.} $h_l\ne h_h$, since $\hi(h_l)\le\lo(h_l)+1\le2<k-2$. If $h_l<h_h$, let $a=h_l$, $b=h_h$, so $r\le a<b\le s$. Take the prefix
fold of Proposition~\ref{prop:prefix-fold} for $(r,r',a)$ and the suffix fold of Corollary~\ref{cor:suffix-fold} for $(s',s,b)$, and splice
them (Proposition~\ref{prop:splice}). The schedule $\beta$ satisfies $\beta(r)=0$ and $\beta(t)\le\beta(s)$, and its blocks are small.
Moreover $\beta(r')=0$: either $r'\le a$ and $\beta_P(r')=0$, and Proposition~\ref{prop:splice}(c) applies; or $a=r$ and $U=0$, and then
$r'=a+1<b$ and $\lo(r')\le\lo(a)+1\le2$, so again (c) applies. Symmetrically $\beta(s')=\beta(s)$ by (d). Lemma~\ref{lem:assemble} finishes
this case. If $h_h<h_l$, apply this to $E$ and the reversed chain $\overline D$ (in which $E_{h_h}$ meets $\overline D_0$ and $E_{h_l}$ meets
$\overline D_{k-1}$; $E$ is crooked in $\overline D$ by Lemma~\ref{lem:reverse}); a block inside $\overline D[j,j+2]$ lies inside
$D[k-3-j,k-1-j]$.

\emph{$k\in\set{4,5}$.} Every index lies in $A\cup B$: if $\lo(t)\ge3$ then $\hi(t)\ge3\ge k-3$. Also, an index $t$ not in $B$ has
$\hi(t)\le k-4\le1$, so its neighbours have $\lo\le2$: \emph{if one of two adjacent indices is not in $B$, both are in $A$}; similarly, an
index not in $A$ has $\lo\ge3$, so its neighbours have $\hi\ge2\ge k-3$: \emph{if one of two adjacent indices is not in $A$, both are in
$B$}. Hence each of the pairs $\set{r,r'}$ and $\set{s',s}$ lies in $A$ or in $B$.
\begin{itemize}
\item If $r,r'\in A$ and $s',s\in B$, apply Lemma~\ref{lem:windows} (first case). If $r,r'\in B$ and $s',s\in A$, apply it with $A$ and $B$
exchanged.
\item Otherwise, suppose $r,r'\in A$. Then $s',s\in A$ (since not both are in $B$), and $h_h$ lies strictly between $r'$ and $s'$: if
$h_h\in\set{r,r'}$, then $\hi\ge k-3$ on both $r,r'$ (as $\hi(h_h)\ge k-2$), so $r,r'\in B$ and we would be in the previous case; if
$h_h\in\set{s',s}$, similarly $s',s\in B$. Since $h_h\in B$, Lemma~\ref{lem:windows} (second case) applies with $c=h_h$.
\item Otherwise $r,r'$ are not both in $A$, so both are in $B$, and then $s',s$ are not both in $A$, so both are in $B$. Now $h_l\in A$
lies strictly between $r'$ and $s'$: if $h_l\in\set{r,r'}$, then $\lo\le2$ on both $r,r'$, i.e.\ $r,r'\in A$; if $h_l\in\set{s',s}$, then
similarly $s',s\in A$, a contradiction. Apply Lemma~\ref{lem:windows} (second case) with $A$ and $B$ exchanged and $c=h_l$.\qedhere
\end{itemize}
\end{proof}

\formal{\lean{HPC.Fold.theorem4_pair}; the main case is \lean{HPC.Fold.B79.theorem4_pair_core}, and the reflection uses
\lean{HPC.Fold.B78.unionRange_reverse}.}


\section*{Exercises}
\addcontentsline{toc}{section}{Exercises}

\begin{exercise}
Show that no sequence beginning with the values $0,1,2,3$ is crooked. Show that every sequence with unit steps whose values lie in an
interval of length at most $2$ is crooked.
\end{exercise}

\begin{exercise}
Compute the jump walk of the turning levels $0,4,1,5$ and the leg positions $\pi_0(y)$ for $-1\le y\le5$. Check
Lemma~\ref{lem:legpos} directly in this case.
\end{exercise}

\begin{exercise}
Let $S(x)=x$ on $[-3,3]$. Show that both conditions $(\mathrm{OK}^+)$ and $(\mathrm{OK}^-)$ of Lemma~\ref{lem:first-direction} fail, and
verify directly that $S$ is not crooked.
\end{exercise}

\begin{exercise}
In Lemma~\ref{lem:up-phase}, explain why the minimum $t=\min(v_+,v_-)$ cannot be replaced by $\max(v_+,v_-)$.
\end{exercise}

\begin{exercise}
Let $\ell,h$ be adjacent indices with $\lo(h)=\lo(\ell)+1$. Show that $\hi(h)=\lo(h)$ and $\hi(\ell)\in\set{\lo(\ell),\lo(\ell)+1}$. Give an
example with chains of open intervals in $\R$ in which $\hi(\ell)=\lo(\ell)+1$.
\end{exercise}

\begin{exercise}
Where exactly is the hypothesis $k\ge6$ used in the proof of Proposition~\ref{prop:splice}? What goes wrong with the middle schedule
$g$ for $k=5$?
\end{exercise}

\begin{exercise}
Deduce from Theorem~\ref{thm:folding} the single-link version: if $r\le s$ and some $E_h$, $E_{h'}$ with $r\le h,h'\le s$ meet $D_0$ and
$D_{k-1}$, then there is a block map with small blocks sending $r$ to the first and $s$ to the last block. (Hint: use the pairs
$(r,r+1)$ and $(s-1,s)$; treat $s\le r+1$ separately.)
\end{exercise}

%======================================================================
\chapter{Every pseudo-arc has a description}\label{ch:descriptions}
%======================================================================

This chapter proves Part~II of the uniqueness theorem.

\begin{theorem}[Bing 1951, Theorem~1, body of the proof]\label{thm:desc-exists}
Let $M\subseteq X$ be a pseudo-arc, and let $p,q\in M$ with $q\notin\cmp Mp$. Then:
\begin{enumerate}[label=(\alph*)]
\item in the compact metric space $M$ there is a description from $p$ to $q$ whose described set is all of $M$;
\item in $X$ there is a description from $p$ to $q$ whose described set is $M$.
\end{enumerate}
\end{theorem}

\formal{(a) is \lean{HPC.Plane.PA51.exists_description_sub}; (b) is
\lean{HPC.Plane.exists_description_of_chainable_HI} (in \lean{HPC/Plane/PseudoArc51.lean}).}

Only (a) is needed for the uniqueness theorem; (b) is proved in Section~\ref{sec:lift} for completeness. The proof of (a) takes place
in a compact metric space $X$ which is itself the continuum, and has three steps:
\begin{enumerate}
\item fine chains covering $X$ are eventually crooked in a given coarse chain covering $X$ (Section~\ref{sec:crooked-eventually});
\item hence there is a sequence of chains covering $X$, each crooked in the previous one, with nested closures and mesh tending to
zero, but not yet going from $p$ to $q$ (Section~\ref{sec:predesc});
\item the folding theorem turns deep levels of this sequence into chains from $p$ to $q$ (Section~\ref{sec:desc-univ}).
\end{enumerate}
Throughout Sections~\ref{sec:crooked-eventually}--\ref{sec:desc-univ}, $X$ is a compact metric space.

\section{Crookedness eventually}\label{sec:crooked-eventually}

\begin{theorem}[Bing 1951]\label{thm:eventually-crooked}
Suppose every subcontinuum of $X$ is indecomposable. Let $C$ be a chain of length $n$ covering $X$, and for $m\in\N$ let $E^{(m)}$ be a
chain refining $C$ whose links are bounded with diameter $\le\varepsilon_m$, where $\varepsilon_m\to0$. Then $E^{(m)}$ is crooked in $C$
for all sufficiently large $m$.
\end{theorem}

\formal{\lean{HPC.Plane.PA51.eventually_crookedIn} (in \lean{HPC/Plane/PA51Crook.lean}).}

The idea: a failure of crookedness of a fine chain $E$ at a pair $(h,k)$ produces two long pieces of $E$, one avoiding $C_{k-1}$ and
the other avoiding $C_{h+1}$, which meet. Passing to the limit along infinitely many failures produces two subcontinua of $X$ with a
common point, neither containing the other, and their union is a decomposable continuum.

\subsection{The data of a failure}

\begin{lemma}\label{lem:fail-data}
Let $E$ be a chain refining $C$, whose links are bounded with diameter $\le\varepsilon$, and suppose $E$ is not crooked in $C$. Then
there are $h,k$ with $k>h+2$ and sets $W,V\subseteq X$ such that
\begin{enumerate}[label=(\alph*)]
\item $W$ and $V$ are $\varepsilon$-chained (\vone{Definition~1.5});
\item there are $x\in W$ and $x'\in V$ with $d(x,x')\le\varepsilon$;
\item $W$ meets $C_h$ and $V$ meets $C_k$;
\item $W$ is a union of links of $E$ none of which lies in $C_{k-1}$, and $V$ is a union of links of $E$ none of which lies in $C_{h+1}$.
\end{enumerate}
\end{lemma}

\begin{proof}
Since $E$ refines $C$ but is not crooked in it, there are $h,k$ with $k>h+2$ and indices $i,j$ with $E_i$ meeting $C_h$ and $E_j$
meeting $C_k$, such that there are no $r,s$ with $\mathrm{Btw}(i,r,s,j)$, $E_r\subseteq C_{k-1}$ and $E_s\subseteq C_{h+1}$. Let $\varphi$
be a refinement map. Then $i\ne j$: otherwise $C_{\varphi(i)}\supseteq E_i$ would meet both $C_h$ and $C_k$, so $\abs{\varphi(i)-h}\le1$ and
$\abs{\varphi(i)-k}\le1$, contradicting $k>h+2$. Reversing $E$ if necessary (which preserves the condition $\mathrm{Btw}$, and the
conclusion), we may assume $i<j$.

Since $C_{\varphi(i)}$ meets $C_h$ and $C_{\varphi(j)}$ meets $C_k$, we have $\varphi(i)\le h+1<k-1\le\varphi(j)$. By the discrete
intermediate value theorem (\vone{Lemma~2.5}) there is $t$ with $i<t\le j$ and $\varphi(t)=k-1$, so $E_t\subseteq C_{k-1}$. Let $r$ be the
least $t>i$ with $E_t\subseteq C_{k-1}$; then $i<r\le j$. Put
\[
W=E_i\cup\dots\cup E_{r-1},\qquad V=E_c\cup\dots\cup E_j,\quad c=\min(r+1,j).
\]
(a) holds by Lemma~\ref{lem:span-chained}. (b): pick $x\in E_{r-1}\cap E_r$ and $x'\in E_r\cap E_c$ (note $c\in\set{r,r+1}$); both lie in
$E_r$, so $d(x,x')\le\varepsilon$. (c) is clear. (d): $E_i$ meets $C_h$, which is disjoint from $C_{k-1}$ (as $k-1>h+1$), so $E_i\not\subseteq
C_{k-1}$; and $E_t\not\subseteq C_{k-1}$ for $i<t<r$ by the choice of $r$. Similarly $E_j$ meets $C_k$, which is disjoint from $C_{h+1}$. Finally,
if $r<t<j$, then $\mathrm{Btw}(i,r,t,j)$ and $E_r\subseteq C_{k-1}$, so $E_t\not\subseteq C_{h+1}$ by the choice of $(h,k,i,j)$.
\end{proof}

\formal{\lean{HPC.Plane.PA51.FailData} (which records the property (d) through small sets, the links, none inside $C_{k-1}$, resp.\
$C_{h+1}$), \lean{HPC.Plane.PA51.failData_of_fail_lt} (the case $i<j$) and \lean{HPC.Plane.PA51.failData_of_not_crooked}.}

\subsection{Tools for the limit}

Recall the upper limit $\limsup_jA_j$ (\vone{Definition~1.6}) and that, in a compact space, the upper limit of $\delta_j$-chained sets
with $\delta_j\to0$ which all contain a point $R$ is a continuum containing $R$ (\vone{Lemma~1.7}).

\begin{lemma}\label{lem:limit-tools}
\begin{enumerate}[label=(\alph*)]
\item If $A$ is $\varepsilon$-chained and $a\in A$, then $A\cup\set R$ is $(\varepsilon+d(a,R))$-chained.
\item Let $A_j\subseteq X$, let $a_j\in A_j$ with $a_j\to R$, let $\varepsilon_j\to0$, and let $O$ be open. Suppose every $u\in A_j$ lies in a
bounded set $S$ with $\diam S\le\varepsilon_j$ and $S\not\subseteq O$. Then $\limsup_j(A_j\cup\set R)$ does not meet $O$.
\item Let $A_j\subseteq X$, $R\in X$, and $y_j\in A_j\cap O$ for all $j$. Then there is a point $Y\in\limsup_j(A_j\cup\set R)$ every open
neighbourhood of which meets $O$.
\end{enumerate}
\end{lemma}

\begin{proof}
(a) Let $A\cup\set R\subseteq U\cup V$ with $d(u,v)>\varepsilon+d(a,R)$ for $u\in U$, $v\in V$. In particular $U,V$ are more than
$\varepsilon$ apart, so $A\subseteq U$, say. If $R\in V$, then $d(a,R)>\varepsilon+d(a,R)$, which is absurd. So $A\cup\set R\subseteq U$.

(b) Suppose $w\in\limsup_j(A_j\cup\set R)$ lies in $O$, and let $B(w,\rho)\subseteq O$. Choose $N$ with $\varepsilon_j<\rho/2$ and
$d(a_j,R)<\rho/4$ for $j\ge N$. As $w\in\cl{\bigcup_{j\ge N}(A_j\cup\set R)}$, there are $j\ge N$ and $v\in A_j\cup\set R$ with
$d(v,w)<\rho/4$. Let $u=v$ if $v\in A_j$ and $u=a_j$ if $v=R$; in both cases $u\in A_j$ and $d(u,w)<\rho/2$. The set $S\ni u$ given by
the hypothesis satisfies $d(s,w)\le d(s,u)+d(u,w)<\varepsilon_j+\rho/2<\rho$ for $s\in S$, so $S\subseteq B(w,\rho)\subseteq O$, a
contradiction.

(c) By compactness some subsequence $y_{j_1},y_{j_2},\dots$ converges to a point $Y$. For each $N$, all but finitely many of its terms lie in
$\bigcup_{j\ge N}A_j$, so $Y\in\limsup_j(A_j\cup\set R)$. Every open neighbourhood of $Y$ contains some $y_{j_l}\in O$.
\end{proof}

\formal{\lean{HPC.Plane.PA51.chainedAt_union_singleton}, \lean{HPC.Plane.PA51.limSup_miss} (with
\lean{HPC.Plane.PA51.approx_of_mem_limSup}), \lean{HPC.Plane.PA51.exists_limSup_near}.}

\begin{lemma}\label{lem:side}
Let $C$ be a chain covering $X$, let $S$ be a preconnected set missing the link $C_c$, and let $y\in S\cap C_l$. If $l<c$, then $S$ lies
in the union $C_{<c}$ of the links of index $<c$; if $l>c$, then $S$ lies in the union $C_{>c}$ of the links of index $>c$.
\end{lemma}

\begin{proof}
$C_{<c}$ and $C_{>c}$ are open, and disjoint (a link of index $<c$ and a link of index $>c$ have indices differing by at least $2$). They
cover $S$, as every point of $S$ lies in a link other than $C_c$. Since $S$ is preconnected, it lies in one of them, and $y$ decides which.
\end{proof}

\formal{\lean{HPC.Plane.PA51.subset_side_of_miss}, with \lean{HPC.Plane.PA51.lowU} and \lean{HPC.Plane.PA51.highU}.}

\subsection{The limit argument}

\begin{proposition}\label{prop:limit-false}
Suppose every subcontinuum of $X$ is indecomposable, and let $C$ be a chain covering $X$. Let $k>h+2$ and $\varepsilon_j\to0$. Then it is
impossible that for every $j$ there are a chain $E_j$ refining $C$, with links bounded of diameter $\le\varepsilon_j$, and sets $W_j,V_j$
satisfying (a)--(d) of Lemma~\ref{lem:fail-data} for $E_j$, $(h,k)$ and $\varepsilon_j$.
\end{proposition}

\begin{proof}
Suppose there are, with points $x_j\in W_j$, $x'_j\in V_j$, $d(x_j,x'_j)\le\varepsilon_j$, and $y_j\in W_j\cap C_h$, $z_j\in V_j\cap C_k$.
Passing to a subsequence we may assume $x_j\to R$; then also $x'_j\to R$. Let $A_j=W_j\cup\set R$ and $B_j=V_j\cup\set R$. By
Lemma~\ref{lem:limit-tools}(a), $A_j$ and $B_j$ are $\delta_j$-chained with $\delta_j=2\varepsilon_j+d(x_j,R)\to0$, and both contain $R$. By
\vone{Lemma~1.7},
\[
\mathcal W=\limsup_jA_j\qquad\text{and}\qquad\mathcal V=\limsup_jB_j
\]
are continua containing $R$.

\emph{$\mathcal W$ misses $C_{k-1}$ and $\mathcal V$ misses $C_{h+1}$}, by Lemma~\ref{lem:limit-tools}(b), with the links of $E_j$ (of diameter
$\le\varepsilon_j$) as the sets $S$, and the points $x_j$, resp.\ $x'_j$, tending to $R$.

\emph{Points near the ends.} By Lemma~\ref{lem:limit-tools}(c) there are $Y\in\mathcal W$ whose neighbourhoods meet $C_h$ and $Z\in\mathcal V$
whose neighbourhoods meet $C_k$. Let $Y\in C_{l_1}$ and $Z\in C_{l_2}$. Then $C_{l_1}$ meets $C_h$, so $l_1\le h+1$, and $C_{l_2}$ meets
$C_k$, so $l_2\ge k-1$.

\emph{$\mathcal W$ meets $C_{h+1}$.} Otherwise $l_1\ne h+1$, so $l_1<h+1$, and Lemma~\ref{lem:side} gives $\mathcal W\subseteq C_{<h+1}$. Also $\mathcal
V$ misses $C_{h+1}$ and $l_2\ge k-1>h+1$, so $\mathcal V\subseteq C_{>h+1}$. But $R$ lies in both, and $C_{<h+1}\cap C_{>h+1}=\emptyset$.

\emph{$\mathcal V$ meets $C_{k-1}$.} Otherwise $l_2>k-1$ and $\mathcal V\subseteq C_{>k-1}$; and $\mathcal W$ misses $C_{k-1}$ with $l_1\le h+1<k-1$, so
$\mathcal W\subseteq C_{<k-1}$. Again $R$ lies in both.

\emph{Conclusion.} $\mathcal W\cup\mathcal V$ is a continuum (two continua with the common point $R$), so it is indecomposable. Hence
$\mathcal W=\mathcal W\cup\mathcal V$ or $\mathcal V=\mathcal W\cup\mathcal V$. In the first case $\mathcal V\subseteq\mathcal W$; but $\mathcal V$ meets
$C_{k-1}$ and $\mathcal W$ does not. In the second case $\mathcal W\subseteq\mathcal V$; but $\mathcal W$ meets $C_{h+1}$ and $\mathcal V$ does not.
\end{proof}

\formal{\lean{HPC.Plane.PA51.limit_false}.}

\begin{proof}[Proof of Theorem~\ref{thm:eventually-crooked}]
Suppose $E^{(m)}$ is not crooked in $C$ for infinitely many $m$. For each such $m$, Lemma~\ref{lem:fail-data} gives a pair $(h,k)$ and sets
$W,V$ for $\varepsilon_m$. There are finitely many pairs, so one pair $(h,k)$ occurs for infinitely many $m$, say $m_0<m_1<\cdots$. Then
the data for $\varepsilon_{m_j}\to0$ contradict Proposition~\ref{prop:limit-false}.
\end{proof}

\begin{remark}
The formal proof follows Bing's argument (1951, first paragraph of the proof of Theorem~1), with the pieces of $E$ chosen explicitly. The
choice of $r$ as the \emph{first} index after $i$ with $E_r\subseteq C_{k-1}$ is what makes property~(d) of Lemma~\ref{lem:fail-data} hold
for $W$, and hence makes the two limit continua visibly not nested.
\end{remark}

\section{Bing's sequence of chains}\label{sec:predesc}

\begin{definition}\label{def:predesc}
A \emph{pre-description} of $X$ is a sequence of chains $D^{(0)},D^{(1)},\dots$ in $X$, $D^{(i)}$ of length $n_i$, such that for every $i$:
$D^{(i)}$ covers $X$; $D^{(i+1)}$ is crooked in $D^{(i)}$; the closure of each link of $D^{(i+1)}$ lies in some link of $D^{(i)}$; and
the links of $D^{(i)}$ are bounded with diameter $\le1/(i+1)$.
\end{definition}

So a pre-description is a description (\vone{Definition~3.1}) except that the chains need not go from a fixed point to another.

\formal{\lean{HPC.Plane.PA51.PreDesc}.}

In the rest of this section $X$ is \emph{snake-like} (for every $\varepsilon>0$ some $\varepsilon$-chain covers $X$) and every subcontinuum
of $X$ is indecomposable.

\begin{lemma}\label{lem:crooked-next}
Let $C$ be a chain covering $X$ and $\eta>0$. Then there is a chain $C'$ covering $X$, crooked in $C$, such that the closure of each link
of $C'$ lies in a link of $C$ and the links of $C'$ are bounded with diameter $\le\eta$.
\end{lemma}

\begin{proof}
By the Lebesgue number lemma there is $\lambda>0$ such that every ball $B(y,\lambda)$ lies in some link of $C$. For each $m$ choose an
$\eta_m$-chain $E^{(m)}$ covering $X$, where $\eta_m=\min(\lambda/2,1/(m+1))$. If $L$ is a link of $E^{(m)}$ and $y\in L$, then $\cl
L\subseteq\set{x:d(x,y)\le\diam L}\subseteq B(y,\lambda)$, which lies in a link of $C$. In particular $E^{(m)}$ refines $C$, and its mesh
is $\le1/(m+1)\to0$. By Theorem~\ref{thm:eventually-crooked}, $E^{(m)}$ is crooked in $C$ for all large $m$; take such an $m$ with also
$1/(m+1)\le\eta$, and $C'=E^{(m)}$.
\end{proof}

\formal{\lean{HPC.Plane.PA51.exists_crooked_next}.}

\begin{proposition}\label{prop:predesc-exists}
$X$ has a pre-description.
\end{proposition}

\begin{proof}
Let $D^{(0)}$ be a $1$-chain covering $X$, and given $D^{(i)}$ let $D^{(i+1)}$ be the chain $C'$ of Lemma~\ref{lem:crooked-next} for
$C=D^{(i)}$ and $\eta=1/(i+2)$.
\end{proof}

\formal{\lean{HPC.Plane.PA51.exists_preDesc}.}

\begin{lemma}\label{lem:predesc-basic}
Let $D^{(i)}$ be a pre-description of $X$.
\begin{enumerate}[label=(\alph*)]
\item If $i\le j$, then $D^{(j)}$ refines $D^{(i)}$; if $i<j$, then $D^{(j)}$ is crooked in $D^{(i)}$.
\item If $X$ has two distinct points, then $n_m\ge5$ for all sufficiently large $m$.
\end{enumerate}
\end{lemma}

\begin{proof}
(a) The first claim is clear by induction. For the second, $D^{(j)}$ refines $D^{(i+1)}$, which is crooked in $D^{(i)}$; apply
\vone{Lemma~2.12}. (b) Let $p\ne q$ and suppose $n_m\le4$. Since $D^{(m)}$ covers $X$, $p\in D^{(m)}_a$ and $q\in D^{(m)}_b$ with $\abs{a-b}\le3$,
so $d(p,q)\le4/(m+1)$ by \vone{Lemma~2.9}. This fails once $4/(m+1)<d(p,q)$.
\end{proof}

\formal{\lean{HPC.Plane.PA51.PreDesc.refines_of_le}, \lean{HPC.Plane.PA51.PreDesc.crooked_of_lt},
\lean{HPC.Plane.PA51.PreDesc.eventually_five_le}.}

\section{From a pre-description to a description}\label{sec:desc-univ}

In this section $D^{(i)}$ is a pre-description of $X$, and $R\ne S$ are points of $X$ with $S\notin\cmp XR$.

\begin{lemma}[Reaching both ends]\label{lem:meets-ends}
For every $j$ there are $k>j$ and indices $a<b$ of $D^{(k)}$ such that the subchain of $D^{(k)}$ from $a$ to $b$ joins $R$ and $S$
(Definition~\ref{def:subchain-from}), and some links $D^{(k)}_{t_1}$, $D^{(k)}_{t_2}$ with $a\le t_1,t_2\le b$ meet the first and the last link
of $D^{(j)}$, respectively.
\end{lemma}

\begin{proof}
Choose $k_0$ with $1/(k+1)<d(R,S)$ for $k\ge k_0$, and let $k\ge\max(k_0,j+1)$. No link of $D^{(k)}$ contains both $R$ and $S$, so by
Lemma~\ref{lem:subchain-exists} there are $a_k<b_k$ such that the subchain from $a_k$ to $b_k$ joins $R$ and $S$; let $A_k$ be its span. If the
lemma fails, then for each such $k$ the set $A_k$ misses $O_1=D^{(j)}_0$ or misses $O_2=D^{(j)}_{n_j-1}$: if $A_k$ met both, at points of
links $D^{(k)}_{t_1}$ and $D^{(k)}_{t_2}$ with indices in $[a_k,b_k]$, these would be the required links. So for some $c\in\set{1,2}$ the set
$I$ of indices $k$ with $A_k\cap O_c=\emptyset$ is infinite.

Let $Z$ be the upper limit of the sets $A_k$, $k\in I$ (enumerated increasingly). Each $A_k$ is $1/(k+1)$-chained (Lemma~\ref{lem:span-chained})
and contains $R$ and $S$, so $Z$ is a continuum containing $R$ and $S$ by \vone{Lemma~1.7}. Moreover $Z\ne X$: the nonempty open set $O_c$ misses
$\bigcup_{k\in I}A_k$, hence also its closure, hence $Z$. So $Z$ is a proper subcontinuum of $X$ containing $R$ and $S$, and $S\in\cmp XR$, a
contradiction.
\end{proof}

\formal{\lean{HPC.Plane.PA51.PreDesc.exists_subchain_meets_ends}, with the limit step
\lean{HPC.Plane.PA51.PreDesc.false_of_disjoint_limit}.}

\begin{lemma}[One step]\label{lem:desc-step}
For every $L$ there are $k>L$ and a chain $F$ from $R$ to $S$ which is a consolidation of $D^{(k)}$, such that the closure of each link of
$F$ lies in a link of $D^{(L)}$ and every link of $F$ has diameter $\le1/(L+1)$.
\end{lemma}

\begin{proof}
By the Lebesgue number lemma there is $\delta>0$ such that every ball $B(x,\delta)$ lies in a link of $D^{(L)}$. Let
$\eta=\min(\delta,1/(L+1))$, and choose $m\ge L+1$ with $1/(m+1)<\eta/3$ and $n_m\ge5$ (Lemma~\ref{lem:predesc-basic}(b)). Lemma~\ref{lem:meets-ends}
for $j=m$ gives $k>m$, indices $a<b$, and $t_1,t_2$. By Lemma~\ref{lem:predesc-basic}(a), $D^{(k)}$ is crooked in $D^{(m)}$.

If the subchain from $a$ to $b$ goes from $R$ to $S$, Corollary~\ref{cor:block-chain} gives a block map $\beta$ with small blocks such that
$F=(D^{(k)})^\beta$ goes from $R$ to $S$. If it goes from $S$ to $R$, the corollary gives a block chain from $S$ to $R$, and we let $F$ be its
reversal. In both cases $F$ is a consolidation of $D^{(k)}$ (Lemma~\ref{lem:block-chain}), and every link of $F$ lies in the union of three
consecutive links of $D^{(m)}$; by \vone{Lemma~2.9}, any two of its points have distance $\le3/(m+1)<\eta$. So $\diam F_u<\eta\le1/(L+1)$,
and if $x_0\in F_u$, then $\cl{F_u}\subseteq\set{x:d(x,x_0)\le3/(m+1)}\subseteq B(x_0,\delta)$, which lies in a link of $D^{(L)}$.
\end{proof}

\formal{\lean{HPC.Plane.PA51.PreDesc.exists_step}; the analogous step for descriptions is \lean{HPC.Description.exists_step} (in
\lean{HPC/PseudoArc/T13Step.lean}).}

\begin{theorem}\label{thm:desc-univ}
There is a description from $R$ to $S$ in $X$ whose described set is $X$.
\end{theorem}

\begin{proof}
For each $L$ choose $k(L)>L$ and $F^{(L)}$ as in Lemma~\ref{lem:desc-step}. Define levels $\ell_0=0$ and $\ell_{i+1}=k(\ell_i)+1$; then
$\ell_i\ge i$. Let $G^{(i)}=F^{(\ell_i)}$ and $\varepsilon_i=1/(i+1)$. We check \vone{Definition~3.1}.
\begin{itemize}
\item $G^{(i)}$ goes from $R$ to $S$.
\item \emph{Crookedness.} Write $k=k(\ell_i)$. The chain $D^{(k+1)}$ is crooked in $D^{(k)}$, and $G^{(i)}$ is a consolidation of $D^{(k)}$,
so $D^{(k+1)}$ is crooked in $G^{(i)}$ by Theorem~\ref{thm:bing1}. The closures of the links of $G^{(i+1)}=F^{(k+1)}$ lie in links of
$D^{(k+1)}$, so $G^{(i+1)}$ refines $D^{(k+1)}$ and is crooked in $G^{(i)}$ by \vone{Lemma~2.12}.
\item \emph{Closures.} For a link $G^{(i+1)}_t$: its closure lies in a link $D^{(k+1)}_{s_1}$, whose closure lies in a link $D^{(k)}_{s_2}$,
which lies in a link of its consolidation $G^{(i)}$.
\item \emph{Mesh.} The links of $G^{(i)}$ have diameter $\le1/(\ell_i+1)\le1/(i+1)$.
\item \emph{Described set.} The union of a consolidation is the union of the chain consolidated, so $\bigcup G^{(i)}=\bigcup D^{(k(\ell_i))}=X$
for every $i$, and the described set is $X$.\qedhere
\end{itemize}
\end{proof}

\formal{\lean{HPC.Plane.PA51.PreDesc.exists_description}, and with Proposition~\ref{prop:predesc-exists},
\lean{HPC.Plane.PA51.exists_description_univ}.}

\begin{proof}[Proof of Theorem~\ref{thm:desc-exists}(a)]
Regard $M$ as a compact metric space. Every subcontinuum of the space $M$ is indecomposable, $q$ is not in the composant of $p$ in the space
$M$ (Lemma~\ref{lem:sub-HI}), and the space $M$ is snake-like (Corollary~\ref{cor:sub-snake}). Also $p\ne q$ (Lemma~\ref{lem:ne-of-composant}).
By Proposition~\ref{prop:predesc-exists} the space $M$ has a pre-description, and Theorem~\ref{thm:desc-univ} gives the description.
\end{proof}

\section{Back to the ambient space}\label{sec:lift}

Let $M\subseteq X$ be compact. For a relatively open set $U\subseteq M$ (that is, $U=O\cap M$ with $O$ open in $X$) and $r>0$ let
\[
\mathrm{Ext}_r(U)=\set{x\in X: d(x,U)<d(x,M\setminus U)\text{ and }d(x,M)<r},
\]
with the convention $d(x,\emptyset)=\infty$. It is open in $X$, since distances to sets are continuous.

\begin{lemma}\label{lem:ext}
Let $U,V\subseteq M$ be relatively open and $r,r'>0$.
\begin{enumerate}[label=(\alph*)]
\item $\mathrm{Ext}_r(U)\cap M=U$.
\item If $U\cap V=\emptyset$, then $\mathrm{Ext}_r(U)\cap\mathrm{Ext}_{r'}(V)=\emptyset$.
\item If $U\subseteq V$ and $r\le r'$, then $\mathrm{Ext}_r(U)\subseteq\mathrm{Ext}_{r'}(V)$.
\item Every point of $\mathrm{Ext}_r(U)$ is at distance $<r$ from $U$; so if $U$ is bounded, $\diam\mathrm{Ext}_r(U)\le\diam U+2r$.
\item If the closure of $U$ lies in $V$ and $\rho>0$, then there is $r\in(0,\rho)$ with $\cl{\mathrm{Ext}_r(U)}\subseteq\mathrm{Ext}_\rho(V)$.
\end{enumerate}
\end{lemma}

\begin{proof}
Note that $M\setminus U$ is closed, as $M$ is closed and $U$ is relatively open.
(a) If $z\in M\setminus U$, then $d(z,M\setminus U)=0$, so $z\notin\mathrm{Ext}_r(U)$. If $z\in U$, then $d(z,U)=0<d(z,M\setminus U)$, as $z$ is not in
the closed set $M\setminus U$, and $d(z,M)=0<r$.
(b) Now $V\subseteq M\setminus U$ and $U\subseteq M\setminus V$. A common point $x$ would satisfy $d(x,U)<d(x,M\setminus U)\le d(x,V)<d(x,M\setminus
V)\le d(x,U)$.
(c) $d(x,V)\le d(x,U)<d(x,M\setminus U)\le d(x,M\setminus V)$, and $d(x,M)<r\le r'$.
(d) $d(x,M)=\min(d(x,U),d(x,M\setminus U))=d(x,U)$ for $x\in\mathrm{Ext}_r(U)$.
(e) The closure $K$ of $U$ is compact and disjoint from the closed set $M\setminus V$, so there is $\delta>0$ such that no point is at distance
$<\delta$ from both. Let $r=\min(\delta,\rho)/2$. If $x\in\cl{\mathrm{Ext}_r(U)}$, then $d(x,U)\le r<\delta$ by (d), so $d(x,K)<\delta$ and
$d(x,M\setminus V)\ge\delta>d(x,U)\ge d(x,V)$; and $d(x,M)\le d(x,U)\le r<\rho$.
\end{proof}

\formal{\lean{HPC.Plane.PA51.ext}, \lean{HPC.Plane.PA51.mem_ext_iff}, \lean{HPC.Plane.PA51.ext_disjoint},
\lean{HPC.Plane.PA51.ext_mono}, \lean{HPC.Plane.PA51.ext_subset_thickening}, \lean{HPC.Plane.PA51.diam_ext_le},
\lean{HPC.Plane.PA51.exists_closure_ext_subset}.}

By (a) and (b), if $C$ is a chain of the space $M$, then $\mathrm{Ext}_r(C)=(\mathrm{Ext}_r(C_0),\dots,\mathrm{Ext}_r(C_{n-1}))$ is a chain in $X$:
its links are open and nonempty, and two of them meet exactly when the corresponding links of $C$ do
(\lean{HPC.Plane.PA51.liftChain}).

\begin{proof}[Proof of Theorem~\ref{thm:desc-exists}(b)]
Let $C^{(i)}$ be the description of part~(a), in the space $M$, with constants $\varepsilon_i$. Choose radii $r_0=1$ and, inductively,
$0<r_{i+1}<r_i$ with $r_{i+1}\le1/(i+2)$ such that for every link $C^{(i+1)}_t$, with $\cl{C^{(i+1)}_t}\subseteq C^{(i)}_{s(t)}$,
\[
\cl{\mathrm{Ext}_{r_{i+1}}(C^{(i+1)}_t)}\subseteq\mathrm{Ext}_{r_i}(C^{(i)}_{s(t)})
\]
(Lemma~\ref{lem:ext}(e) for each of the finitely many $t$, and the minimum of the radii obtained, using (c)). Let $D^{(i)}=\mathrm{Ext}_{r_i}(C^{(i)})$
with constants $\varepsilon_i+2r_i\to0$.
\begin{itemize}
\item By Lemma~\ref{lem:ext}(a), $p$ and $q$ lie in the same links of $D^{(i)}$ as of $C^{(i)}$; so $D^{(i)}$ goes from $p$ to $q$.
\item \emph{Crookedness.} If $C^{(i+1)}_a\subseteq C^{(i)}_s$, then $D^{(i+1)}_a\subseteq D^{(i)}_s$ by Lemma~\ref{lem:ext}(c), since $r_{i+1}\le r_i$; so
$D^{(i+1)}$ refines $D^{(i)}$. If $D^{(i+1)}_a$ meets $D^{(i)}_h$, then $C^{(i+1)}_a$ meets $C^{(i)}_h$ by (b). So the hypotheses of crookedness
for $D$ give those for $C$, and the conclusions ($C^{(i+1)}_u\subseteq C^{(i)}_{k-1}$, $C^{(i+1)}_v\subseteq C^{(i)}_{h+1}$) give the conclusions
for $D$ by (c).
\item Closures are nested by the choice of the radii, and $\diam D^{(i)}_t\le\varepsilon_i+2r_i$ by (d).
\item \emph{Described set.} Every point of $M$ lies in $\bigcup C^{(i)}$ for all $i$ (the described set of $C$ is $M$), hence in $\bigcup
D^{(i)}$. Conversely, a point in all $\bigcup D^{(i)}$ has $d(x,M)<r_i\to0$, so $x\in M$ as $M$ is closed.\qedhere
\end{itemize}
\end{proof}

\formal{\lean{HPC.Plane.PA51.lift_description}.}

\section*{Exercises}
\addcontentsline{toc}{section}{Exercises}

\begin{exercise}
Show that Theorem~\ref{thm:eventually-crooked} fails for $X=[0,1]$: find chains $E^{(m)}$ of intervals with mesh tending to zero which
refine the chain $C=([0,0.3),(0.2,0.5),(0.4,0.7),(0.6,0.9),(0.8,1])$ and are never crooked in it. Which subcontinuum of $[0,1]$ does the proof
produce?
\end{exercise}

\begin{exercise}
Let $A_k\subseteq X$ for $k$ in an infinite set $I\subseteq\N$, and let $O$ be an open set meeting no $A_k$. Show that $O$ does not meet
$\limsup_{k\in I}A_k$. Show by an example that this fails if $O$ is only assumed to be closed.
\end{exercise}

\begin{exercise}
In the proof of Theorem~\ref{thm:desc-univ}, why is the level $\ell_{i+1}$ taken to be $k(\ell_i)+1$ rather than $k(\ell_i)$?
\end{exercise}

\begin{exercise}
Show by an example that Lemma~\ref{lem:ext}(e) fails if $\cl U\not\subseteq V$.
\end{exercise}

\begin{exercise}
Deduce from Theorem~\ref{thm:desc-exists} that if $M$ is a pseudo-arc and $q\notin\cmp Mp$, then for every $\varepsilon>0$ there is an
$\varepsilon$-chain covering $M$ which goes from $p$ to $q$. (Points with this property are called \emph{opposite end points} of $M$.)
\end{exercise}

%======================================================================
\chapter{The main theorems}\label{ch:main}
%======================================================================

This chapter assembles Parts~I--III and proves the theorems announced in the introduction:
Theorems~\ref{thm:intro-unique} and~\ref{thm:intro-homog} and Corollary~\ref{cor:intro}. All the
work has been done; what remains is to put the pieces together, and to be careful about the
ambient spaces. Throughout, $X$ and $X'$ are metric spaces.

Recall (\vone{Definition~2.14}) that a \emph{pseudo-arc} is a nondegenerate, snake-like,
hereditarily indecomposable continuum $M\subseteq X$.

\formal{\lean{HPC.Plane.IsPseudoArc}: the conjunction of \lean{HPC.Plane.IsSnakeLike} (which
contains ``$M$ is a continuum''), \lean{HPC.Plane.IsHereditarilyIndecomposable} and
\lean{Set.Nontrivial}.}

\section{Uniqueness}\label{sec:main-unique}

\begin{theorem}[Bing 1951, Theorem~1, the ``in fact'' clause]\label{thm:unique-points}
Let $M\subseteq X$ and $M'\subseteq X'$ be pseudo-arcs. Let $p,q\in M$ with $q\notin\cmp Mp$, and
let $p',q'\in M'$ with $q'\notin\cmp{M'}{p'}$. Then there is a homeomorphism $h\colon M\to M'$
with $h(p)=p'$ and $h(q)=q'$.
\end{theorem}

\begin{proof}
Regard $M$ and $M'$ as compact metric spaces (with the metrics of $X$ and $X'$).

\emph{Step 1: descriptions.} By Theorem~\ref{thm:desc-exists}(a) there is a description
$C^{(0)},C^{(1)},\dots$ from $p$ to $q$ \emph{in the space $M$} whose described set is the whole
space $M$, and likewise a description from $p'$ to $q'$ in the space $M'$ whose described set is
all of $M'$.

\emph{Step 2: the end points are distinct.} A pseudo-arc has at least two points, so $p\ne q$
and $p'\ne q'$ by Lemma~\ref{lem:ne-of-composant}.

\emph{Step 3: two descriptions.} The spaces $M$ and $M'$ are compact, and the two descriptions go
between distinct points. By Theorem~\ref{thm:two-desc} there is a homeomorphism $T$ from the
described set of the first description onto the described set of the second, with $T(p)=p'$ and
$T(q)=q'$. These described sets are $M$ and $M'$, so $h=T$ is the required homeomorphism.
\end{proof}

\formal{\lean{HPC.Plane.theorem1_HIC}. Formally the described sets are subsets
$\Delta_1.M\subseteq M$ and $\Delta_2.M\subseteq M'$ of the spaces $M$, $M'$, equal to the whole
spaces; the homeomorphism of \lean{HPC.theorem12} is composed with the identifications
$M\cong\Delta_1.M$ and $\Delta_2.M\cong M'$.}

\begin{remark}
Note how little of Theorem~\ref{thm:desc-exists} is used: only part~(a), the description
inside $M$ itself. Working inside $M$ has the advantage that the ambient space is compact, which
is exactly what Theorem~\ref{thm:two-desc} asks for. Part~(b) would give descriptions in $X$ and
$X'$; to use those one would need their links to lie in compact sets
(Corollary~\ref{cor:two-desc-K}), which is true for the lifted descriptions of
Section~\ref{sec:lift} only when $X$ is, for example, a proper space.
\end{remark}

\begin{theorem}[Uniqueness of the pseudo-arc]\label{thm:unique}
Any two pseudo-arcs $M\subseteq X$ and $M'\subseteq X'$ are homeomorphic.
\end{theorem}

\begin{proof}
$M$ is nondegenerate, so it has a point $p$; by Corollary~\ref{cor:other-composant} there is
$q\in M$ with $q\notin\cmp Mp$. Likewise choose $p'\in M'$ and $q'\in M'$ with
$q'\notin\cmp{M'}{p'}$. Apply Theorem~\ref{thm:unique-points}.
\end{proof}

\formal{\lean{HPC.Plane.theorem1_HIC'}; the same statement with the three defining properties
written out is \lean{HPC.Plane.theorem1_snake} (Bing, \emph{Snake-like continua}, Theorem~1).}

Theorem~\ref{thm:intro-unique} of the introduction is Theorems~\ref{thm:unique-points}
and~\ref{thm:unique} together.

\section{Homogeneity}\label{sec:main-homog}

\begin{theorem}[Homogeneity of the pseudo-arc]\label{thm:homog}
Let $M\subseteq X$ be a pseudo-arc and $p,q\in M$. There is a homeomorphism $h\colon M\to M$ with
$h(p)=q$.
\end{theorem}

\begin{proof}
By Corollary~\ref{cor:other-composant} choose $r\in M$ with $r\notin\cmp Mp$ and $s\in M$ with
$s\notin\cmp Mq$. Apply Theorem~\ref{thm:unique-points} with $M'=M$, to the pairs $(p,r)$ and
$(q,s)$: there is a homeomorphism $h\colon M\to M$ with $h(p)=q$ (and $h(r)=s$).
\end{proof}

\formal{\lean{HPC.Plane.homogeneous_of_isPseudoArc}.}

This is Theorem~\ref{thm:intro-homog}. The proof shows more: for any two pairs $(p,r)$ and
$(q,s)$ of points in different composants, some homeomorphism of $M$ onto itself takes $p$ to $q$
and $r$ to $s$ at the same time.

\begin{remark}
It is instructive to see why the argument does not prove that the arc $[0,1]$ is homogeneous,
which is false. The arc is snake-like but decomposable, so Chapter~\ref{ch:descriptions} does not
apply: fine chains covering $[0,1]$ need not be crooked in coarse ones
(Theorem~\ref{thm:eventually-crooked} fails). And even the input of Chapter~\ref{ch:composants}
fails: for $0<p<1$ the composant of $p$ is all of $[0,1]$, so there is no partner point $r$.
\end{remark}

\section{Consequences}\label{sec:main-cons}

\begin{corollary}\label{cor:plane-unique}
\begin{enumerate}[label=(\alph*)]
\item There is a pseudo-arc $M_0$ in the plane $\R^2$.
\item Every pseudo-arc, in any metric space, is homeomorphic to $M_0$.
\end{enumerate}
\end{corollary}

\begin{proof}
(a) is \vone{Theorem~5.18}: the plane description of \vone{Theorem~5.17} goes between two distinct
points and all its links lie in a compact rectangle, so its described set is a pseudo-arc by
\vone{Theorem~3.8}. (b) follows from Theorem~\ref{thm:unique}.
\end{proof}

\formal{\lean{HPC.Plane.isPseudoArc_plane} and \lean{HPC.Plane.pseudoArc_unique}.}

So ``the pseudo-arc'' is well defined up to homeomorphism, and it can be drawn in the plane. This
is Corollary~\ref{cor:intro}(a).

\begin{lemma}\label{lem:sub-pseudoarc}
Every nondegenerate subcontinuum $K$ of a pseudo-arc $M$ is a pseudo-arc.
\end{lemma}

\begin{proof}
$K$ is a nondegenerate continuum by assumption. Every subcontinuum of $K$ is a subcontinuum of
$M$, hence indecomposable; so $K$ is hereditarily indecomposable. Given $\varepsilon>0$, an
$\varepsilon$-chain covering $M$ also covers $K$; so $K$ is snake-like.
\end{proof}

\formal{\lean{HPC.Plane.isPseudoArc_of_subcontinuum}.}

\begin{corollary}\label{cor:sub-homeo}
Every nondegenerate subcontinuum of a pseudo-arc $M$ is homeomorphic to $M$.
\end{corollary}

\begin{proof}
Lemma~\ref{lem:sub-pseudoarc} and Theorem~\ref{thm:unique}.
\end{proof}

\formal{\lean{HPC.Plane.homeomorphic_subcontinuum}.}

This is Corollary~\ref{cor:intro}(b). Together with Corollary~\ref{cor:plane-unique} it shows
that the pseudo-arc is homeomorphic to each of its nondegenerate subcontinua, although it
contains no arc (an arc is decomposable).

\section{End points}\label{sec:main-endpoints}

Bing's paper \emph{Snake-like continua} characterizes the pseudo-arc among snake-like continua by
a property of its points. We prove this characterization, since it follows quickly from what we
have.

\begin{definition}\label{def:endpoint}
Let $M\subseteq X$ be a continuum and $p\in M$.
\begin{enumerate}[label=(\alph*)]
\item $p$ is an \emph{end point} of $M$ if for every $\varepsilon>0$ there is an $\varepsilon$-chain $E$ of
length $n\ge1$ covering $M$ such that $p$ lies in the first link $E_0$ and in no other link.
\item Points $p,q\in M$ are \emph{opposite end points} of $M$ if for every $\varepsilon>0$ there is an
$\varepsilon$-chain covering $M$ which goes from $p$ to $q$.
\end{enumerate}
\end{definition}

\formal{\lean{HPC.Plane.IsEndPoint}, \lean{HPC.Plane.OppositeEndPoints} (in
\lean{HPC/Plane/EndPoints51.lean}).}

\begin{example}
The points $0$ and $1$ are opposite end points of the arc $[0,1]$: cover it by the chain of
intervals $(-\delta,2\delta),(\delta,3\delta),\dots$ of length $3\delta$ and step $\delta$, where $\delta=1/m$
is small, cut down to the links meeting $[0,1]$. A point $p$ with $0<p<1$ is not an end point of
$[0,1]$; this follows from Lemma~\ref{lem:endpoint-nested} below, since $[0,p]$ and $[p,1]$ are
subcontinua containing $p$, neither of which contains the other.
\end{example}

\begin{proposition}\label{prop:opposite}
Let $M$ be a pseudo-arc and $p,q\in M$ with $q\notin\cmp Mp$. Then $p$ and $q$ are opposite end
points of $M$. In particular every point of a pseudo-arc is an end point of it.
\end{proposition}

\begin{proof}
By Theorem~\ref{thm:desc-exists}(b) there is a description $D^{(0)},D^{(1)},\dots$ in $X$ from $p$
to $q$ with described set $M$. Given $\varepsilon>0$, choose $i$ with $\varepsilon_i<\varepsilon$. The
chain $D^{(i)}$ is an $\varepsilon$-chain by~\vone{(D4)}, it covers $M$ because $M$ is the described
set, and it goes from $p$ to $q$ by~\vone{(D1)}.

For the second statement let $p\in M$; by Corollary~\ref{cor:other-composant} there is $q\in M$
with $q\notin\cmp Mp$, and the chains just found have $p$ in the first link only.
\end{proof}

\formal{\lean{HPC.Plane.oppositeEndPoints_of_composants}.}

Here part~(b) of Theorem~\ref{thm:desc-exists} is really used: the chains must be chains of $X$.

\begin{lemma}\label{lem:endpoint-nested}
Let $p$ be an end point of a continuum $M$, and let $A,B\subseteq M$ be continua containing $p$.
Then $A\subseteq B$ or $B\subseteq A$.
\end{lemma}

\begin{proof}
Suppose not: there are $a\in A\setminus B$ and $b\in B\setminus A$. Since $A$ and $B$ are compact,
$\delta=\min(\dist(a,B),\dist(b,A))>0$. Let $E$ be a $\delta$-chain of length $n$ covering $M$ with
$p$ in $E_0$ only.

The links of $E$ that meet $A$ form an interval of indices $[0,i_A]$: it is an interval by
Lemma~\ref{lem:cover-meeting}(a) (applied to the connected set $A$, which is covered by $E$), and
it contains $0$ because $p\in A\cap E_0$. Likewise the links meeting $B$ form an interval
$[0,i_B]$. Suppose $i_A\le i_B$ (the other case is symmetric, exchanging the roles of $a$ and $b$).
The point $a$ lies in some link $E_j$; this link meets $A$, so $j\le i_A\le i_B$, and therefore $E_j$
also meets $B$, say at $z$. Then $\dist(a,B)\le d(a,z)\le\diam E_j<\delta$, a contradiction.
\end{proof}

\formal{\lean{HPC.Plane.condB_of_endPoint}.}

\begin{lemma}\label{lem:not-HI-not-end}
If a continuum $M$ is not hereditarily indecomposable, then some point of $M$ is not an end point
of $M$.
\end{lemma}

\begin{proof}
Let $K\subseteq M$ be a decomposable subcontinuum, $K=A\cup B$ with $A,B$ subcontinua of $K$
different from $K$. Since $K$ is connected and $A,B$ are nonempty closed sets covering it, they
meet; let $p\in A\cap B$. If $p$ were an end point of $M$, Lemma~\ref{lem:endpoint-nested} would
give, say, $A\subseteq B$, and then $K=A\cup B=B$, a contradiction.
\end{proof}

\formal{\lean{HPC.Plane.exists_not_endPoint_of_not_HI}.}

\begin{theorem}[Bing 1951, \emph{Snake-like continua}, Theorem~16]\label{thm:endpoints}
A nondegenerate snake-like continuum $M$ is a pseudo-arc if and only if every point of $M$ is an
end point of $M$.
\end{theorem}

\begin{proof}
If $M$ is a pseudo-arc, every point is an end point by Proposition~\ref{prop:opposite}.
Conversely, if every point is an end point, $M$ is hereditarily indecomposable by
Lemma~\ref{lem:not-HI-not-end}, and hence a pseudo-arc.
\end{proof}

\formal{\lean{HPC.Plane.isPseudoArc_iff_endPoints}.}

\section{What the classification uses}\label{sec:main-uses}

The later volumes use the results of this volume in the following forms.
\begin{itemize}
\item \emph{Uniqueness} (Theorem~\ref{thm:unique}, Corollary~\ref{cor:plane-unique}): to recognize
a continuum as ``the'' pseudo-arc it suffices to check the three defining properties. This is how
continua arising later (for instance the indecomposable homogeneous plane continua) are
identified with the pseudo-arc: one proves that they are snake-like and hereditarily
indecomposable.
\item \emph{Homogeneity} (Theorem~\ref{thm:homog}): the pseudo-arc appears in the list of
homogeneous plane continua, so it must be shown to be homogeneous.
\item \emph{Subcontinua} (Lemma~\ref{lem:sub-pseudoarc}, Corollary~\ref{cor:sub-homeo}).
\end{itemize}

\begin{remark}[On the hypotheses]
Two statements in the literature need an extra hypothesis, which the formal proof adds.
\begin{enumerate}[label=(\alph*)]
\item Bing's Theorem~12 for descriptions whose described sets are merely compact is false
(Remark after Corollary~\ref{cor:two-desc-K}); one needs the links to lie in a compact set.
\item In the same way, ``the described set of a description from $P$ to $Q\ne P$ is a pseudo-arc''
needs the links to lie in a compact set (\vone{Theorem~3.8} and \vone{Remark~3.5}); the formal
record of the counterexample is \lean{HPC.Plane.isPseudoArc_of_description_false}.
\end{enumerate}
Neither affects the theorems of this chapter, whose hypotheses only involve pseudo-arcs, which
are compact by definition.
\end{remark}

\section*{Exercises}
\addcontentsline{toc}{section}{Exercises}

\begin{exercise}\label{ex:arc-endpoints}
Show directly from Definition~\ref{def:endpoint} that the end points of $[0,1]$ are exactly $0$
and $1$. (Hint: use Lemma~\ref{lem:endpoint-nested} with $A=[0,p]$ and $B=[p,1]$.)
\end{exercise}

\begin{exercise}
Let $p$ be an end point of a continuum $M$ and let $K\subseteq M$ be a continuum containing $p$. Show
that $p$ is an end point of $K$. (Hint: Lemma~\ref{lem:cover-meeting}(a).) Deduce again that every
point of a pseudo-arc is an end point of each subcontinuum containing it.
\end{exercise}

\begin{exercise}
Let $M$ be a pseudo-arc and $p\in M$. Show that the homeomorphisms $h\colon M\to M$ with $h(p)=p$
act transitively on $M\setminus\cmp Mp$. (Hint: Theorem~\ref{thm:unique-points} with
$p'=p$.)
\end{exercise}

\begin{exercise}
Show that a pseudo-arc has uncountably many composants. (Hint: each composant is a countable
union of proper subcontinua, each of which is nowhere dense by Theorem~\ref{thm:thin}; use the
Baire category theorem, Lemma~\ref{lem:baire}.)
\end{exercise}

\begin{exercise}
Using Corollary~\ref{cor:sub-homeo}, show that a pseudo-arc $M$ contains a strictly decreasing
sequence $M=K_0\supsetneq K_1\supsetneq K_2\supsetneq\cdots$ of subcontinua, each homeomorphic to
$M$, whose intersection is a single point.
\end{exercise}


\appendix
%======================================================================
\chapter{Correspondence with the formal proof}\label{app:formal}
%======================================================================

The table lists the definitions and results of this volume that carry a ``Formal counterpart''
line, together with the corresponding declarations of the Lean development. All declaration
names start with \lean{HPC.}; file names are relative to the root of the repository. The formal
statements are sometimes slightly more general, or are split differently into lemmas; the table
points to the place where the content of each result is proved. Rows are in the order of the
text.

The main theorems of the volume are, in the Lean development (prefix \lean{HPC.} omitted):
\begin{center}
\begin{tabular}{@{}p{0.42\textwidth}p{0.54\textwidth}@{}}
\toprule
Uniqueness, with two pairs of points (Theorem~\ref{thm:unique-points}) & \lean{Plane.theorem1_HIC}\\
Uniqueness (Theorem~\ref{thm:unique}) & \lean{Plane.theorem1_HIC'}\\
Homogeneity (Theorem~\ref{thm:homog}) & \lean{Plane.homogeneous_of_isPseudoArc}\\
Every pseudo-arc is the plane one (Corollary~\ref{cor:plane-unique}) & \lean{Plane.pseudoArc_unique}\\
Subcontinua (Corollary~\ref{cor:sub-homeo}) & \lean{Plane.homeomorphic_subcontinuum}\\
Two descriptions (Theorem~\ref{thm:two-desc}) & \lean{theorem12}\\
Every pseudo-arc has a description (Theorem~\ref{thm:desc-exists}) &
\lean{Plane.exists_description_of_chainable_HI}\\
\bottomrule
\end{tabular}
\end{center}
All of them are proved in the Lean development without \lean{sorry} and without axioms beyond the
standard ones of Lean and Mathlib.

\begin{longtable}{@{}p{0.40\textwidth}p{0.56\textwidth}@{}}
\toprule
In these notes & In the Lean development\\
\midrule
\endhead
\bottomrule
\endfoot
\multicolumn{2}{@{}l}{\emph{Chapter~\ref{ch:composants}}}\\
Definition~\ref{def:composant} & \lean{HPC.composant} (in \lean{HPC/PseudoArc/Statements79.lean}).\\
Lemma~\ref{lem:ne-of-composant} & \lean{HPC.Plane.PA51.ne_of_not_mem_composant}.\\
Boundary bumping, clopen form (Lemma~\ref{lem:bumping}) & \lean{HPC.T13.clopen_meets_closure_diff}.\\
Lemma~\ref{lem:union-clopen} & \lean{HPC.T13.isPreconnected_union_clopen}.\\
Theorem~\ref{thm:thin} & \lean{HPC.T13.false_of_interior}.\\
Theorem~\ref{thm:composant-proper} & \lean{HPC.T13.exists_not_mem_proper_subcontinua} (in a compact ambient space) and \lean{HPC.composant_ne_of_indecomposable}.\\
Corollary~\ref{cor:other-composant} & \lean{HPC.Plane.exists_not_mem_composant}. Formally the theorem is applied in the compact space $M$, and the subcontinuum found there is transported back to $X$.\\
Lemma~\ref{lem:sub-HI} & \lean{HPC.Plane.PA51.hI_sub}, \lean{HPC.Plane.PA51.not_mem_composant_sub}.\\
Lemma~\ref{lem:cover-meeting} & (a) is \lean{HPC.Plane.exists_cover_meeting}; (b) is \lean{HPC.Plane.PA51.inter_inter_nonempty} and \lean{HPC.Plane.PA51.subChain}.\\
Corollary~\ref{cor:sub-snake} & \lean{HPC.Plane.PA51.snake_sub}.\\[4pt]
\multicolumn{2}{@{}l}{\emph{Chapter~\ref{ch:consolidations}}}\\
Definition~\ref{def:consolidation} & \lean{HPC.Chain.IsConsolidationOf}.\\
Lemma~\ref{lem:block-chain} & \lean{HPC.Chain.IsBlockMap}, \lean{HPC.Chain.ofBlocks}, \lean{HPC.Chain.blockLink_adjacent}, \lean{HPC.Chain.ofBlocks_isConsolidationOf}.\\
Lemma~\ref{lem:block-from} & \lean{HPC.Fold.Q78.isFrom_ofBlocks_iff}.\\
Lemma~\ref{lem:schedule} & \lean{HPC.IsSchedule}, \lean{HPC.IsSchedule.exists_blockMap}.\\
Lemma~\ref{lem:consolidate-length} & \lean{HPC.Chain.exists_consolidation_isFrom}.\\
Lemma~\ref{lem:index-dist} & \lean{HPC.Chain.index_dist_le_of_subset}.\\
Bing 1948, Theorem~1 (Theorem~\ref{thm:bing1}) & \lean{HPC.Chain.isCrookedIn'_of_isConsolidationOf}; the interpolation of Step~4 is \lean{HPC.Chain.exists_value_between}.\\
Lemma~\ref{lem:reverse} & \lean{HPC.Chain.reverse}, \lean{HPC.Chain.IsCrookedIn'.reverse_left}, \lean{HPC.Chain.IsCrookedIn'.reverse_right}.\\
Lemma~\ref{lem:minidx} & \lean{HPC.Chain.Refines.minIdx}, \lean{HPC.Chain.Refines.maxIdx}, \lean{HPC.Chain.subset_adjacent_of_subset}, \lean{HPC.Chain.minIdx_lip}, \lean{HPC.Chain.maxIdx_lip}, \lean{HPC.Chain.minIdx_near_of_inter}.\\
Lemma~\ref{lem:span-chained} & \lean{HPC.Chain.span}, \lean{HPC.Chain.chainedAt_span}.\\
Definition~\ref{def:subchain-from} & \lean{HPC.Chain.IsSubchainFrom}.\\
Lemma~\ref{lem:subchain-exists} & \lean{HPC.Chain.exists_isSubchainFrom}.\\[4pt]
\multicolumn{2}{@{}l}{\emph{Chapter~\ref{ch:patterns}}}\\
Pattern following; Bing 1948, Theorem~6 (Theorem~\ref{thm:pattern-following}) & \lean{HPC.Description.patternFollowing}; the property itself is \lean{HPC.Description.PatternFollowing}, and patterns are \lean{HPC.Description.IsPattern}; the proof is \lean{HPC.Description.theorem6Pos_of_ne}, with \lean{HPC.Description.patternFollowing_of_theorem6Pos} for the passage to arbitrary $i$ (via the shifted description \lean{HPC.Description.shift}).\\
Remark~\ref{rem:pf-false} & \lean{HPC.Description.unitDescription}, \lean{HPC.Description.not_patternFollowing_unconditional}.\\
Lemma~\ref{lem:allowed} & \lean{HPC.PatCore.Setting}, \lean{HPC.PatCore.Slow}, \lean{HPC.PatCore.Al}, \lean{HPC.PatCore.al_canon}, \lean{HPC.PatCore.al_move}, \lean{HPC.PatCore.al_near}, \lean{HPC.PatCore.al_far}, \lean{HPC.PatCore.time_ge}; \lean{HPC.PatCore.ivt_between}.\\
Definition~\ref{def:lift} & \lean{HPC.PatCore.Real}.\\
Operations on lifts (Lemma~\ref{lem:lift-ops}) & \lean{HPC.PatCore.real_glue}, \lean{HPC.PatCore.real_widen}, \lean{HPC.PatCore.real_of_neg}, \lean{HPC.PatCore.real_reflect}, \lean{HPC.PatCore.real_step}, \lean{HPC.PatCore.real_ramp}.\\
Definition~\ref{def:good} & \lean{HPC.PatCore.Good}.\\
Main induction (Theorem~\ref{thm:pf-main}) & \lean{HPC.PatCore.main_step} with the cases \lean{HPC.PatCore.case_const}, \lean{HPC.PatCore.case_plateau}, \lean{HPC.PatCore.case_plateau_end}, \lean{HPC.PatCore.case_amp1}, \lean{HPC.PatCore.case_return}, \lean{HPC.PatCore.case_noreturn}; the induction is \lean{HPC.PatCore.ih_all}.\\
Combinatorial core of pattern following (Theorem~\ref{thm:realize}) & \lean{HPC.PatCore.realize}.\\
Lemma~\ref{lem:n-tendsto} & \lean{HPC.Description.n_tendsto_atTop} (via the more general \lean{HPC.Chain.diam_le_length_mul}).\\
Thick readings (Lemma~\ref{lem:thick}) & \lean{HPC.Description.eventually_thick}.\\
Lemma~\ref{lem:setting-of-crooked} & \lean{HPC.Chain.setting_of_crooked}.\\
Lemma~\ref{lem:consolidation-of-lift} & \lean{HPC.Chain.exists_consolidation_of_lift}.\\[4pt]
\multicolumn{2}{@{}l}{\emph{Chapter~\ref{ch:matching}}}\\
Definition~\ref{def:matching} & \lean{HPC.MatchingCovers}, \leansub{HPC.MatchingCovers.Tracks}{1}, \leansub{HPC.MatchingCovers.Tracks}{2}.\\
Matching covers; Bing 1948, Theorem~11 (Theorem~\ref{thm:matching}) & \lean{HPC.MatchingCovers.exists_homeomorph}.\\
Fattening & \lean{HPC.tail}, \lean{HPC.rad}, \lean{HPC.mesh}, \lean{HPC.fat}, \leansub{HPC.MatchingCovers.y}{1}, \leansub{HPC.MatchingCovers.y}{2}, \lean{HPC.diam_fat_le_mesh}, \lean{HPC.mesh_tendsto_zero}.\\
Nesting (Lemma~\ref{lem:nesting}) & \lean{HPC.closure_fat_subset_fat}.\\
Step 1 (Lemma~\ref{lem:step1}) & \leansub{HPC.MatchingCovers.step1}{2}.\\
Step 2 (Lemma~\ref{lem:step2}) & \leansub{HPC.MatchingCovers.step2}{2}.\\
Step 3 (Lemma~\ref{lem:limit-point}) & \lean{HPC.MatchingCovers.exists_limit_point}.\\
Step 4: independence of choices (Lemma~\ref{lem:T-indep}) & \lean{HPC.MatchingCovers.toFun}, \leansub{HPC.MatchingCovers.toFun_mem_y}{2}.\\
Step 5 (Lemma~\ref{lem:T-cont}) & \lean{HPC.MatchingCovers.continuous_toMap}.\\
Step 6 (Lemma~\ref{lem:T-inj}) & \lean{HPC.MatchingCovers.injective_toMap}.\\
Step 7 (Lemma~\ref{lem:T-surj}) & \lean{HPC.MatchingCovers.surjective_toMap}.\\[4pt]
\multicolumn{2}{@{}l}{\emph{Chapter~\ref{ch:two}}}\\
Bing 1948, Theorem~12 (Theorem~\ref{thm:two-desc}) & \lean{HPC.theorem12}, which combines \lean{HPC.Description.exists_homeomorph} (the theorem assuming pattern following) with \lean{HPC.Description.patternFollowing}.\\
Lemma~\ref{lem:pattern-of-refines} & \lean{HPC.Chain.exists_pattern_of_refines}.\\
Matching (Lemma~\ref{lem:pattern-matching}) & \lean{HPC.Chain.pattern_matching}.\\
Lemma~\ref{lem:consolidation-meets} & \lean{HPC.Description.M_subset_consolidation}, \lean{HPC.Description.consolidation_link_inter_M_nonempty}, \lean{HPC.Description.link_inter_M_nonempty}.\\
Lemma~\ref{lem:level-refining} & \lean{HPC.Description.exists_level_refining}.\\
Definition~\ref{def:rung} & \lean{HPC.Rung}. Formally, (R3) asks that the active side be able to follow patterns (\lean{HPC.Description.CanFollow}), which holds for every level by pattern following (\lean{HPC.Description.canFollow_level}).\\
The first rung (Lemma~\ref{lem:base-rung}) & \lean{HPC.Rung.exists_base}.\\
The next rung (Lemma~\ref{lem:next-rung}) & \lean{HPC.Rung.exists_next}, via \lean{HPC.Description.exists_half_step}.\\
Corollary~\ref{cor:two-desc-K} & \lean{HPC.Plane.theorem12_of}, using the restriction \lean{HPC.Description.restrict} (in \lean{HPC/Plane/Transfer.lean}).\\[4pt]
\multicolumn{2}{@{}l}{\emph{Chapter~\ref{ch:folding}}}\\
Folding theorem; Bing 1948, Theorem~4 (Theorem~\ref{thm:folding}) & \lean{HPC.Fold.theorem4_pair} (in \lean{HPC/Fold/Theorem4Pair.lean}). The version with single links $r$ and $s$ in place of the pairs is \lean{HPC.Fold.theorem4}. In the proof, the main case is \lean{HPC.Fold.B79.theorem4_pair_core}, and the reflection uses \lean{HPC.Fold.B78.unionRange_reverse}.\\
Corollary~\ref{cor:block-chain} & \lean{HPC.Chain.exists_blockChain} (in \lean{HPC/PseudoArc/T13Step.lean}).\\
Lemma~\ref{lem:readings} & \lean{HPC.Fold.IsCrookedIvF} and \lean{HPC.Fold.isCrookedIvF_minIdx_maxIdx} package these facts with crookedness; see also \lean{HPC.Fold.B79.seq_facts}.\\
Crooked sequence (Definition~\ref{def:crooked-seq}) & \lean{HPC.Fold.Z77.CrkOn}.\\
Lemma~\ref{lem:crooked-seq} & \lean{HPC.Fold.Z77.crkOn_of_isCrookedIvF}, \lean{HPC.Fold.Z77.CrkOn.reverse}, \lean{HPC.Fold.B79.minIdx_reverse}.\\
Fitting (Definition~\ref{def:fit}) & \lean{HPC.Fold.Z77.Env}, \lean{HPC.Fold.Z77.arr}, \lean{HPC.Fold.Z77.Fit}.\\
Lemma~\ref{lem:legpos} & \lean{HPC.Fold.Z77.jwalk}, \lean{HPC.Fold.Z77.legPos}, \lean{HPC.Fold.Z77.jwalk_legPos}, \lean{HPC.Fold.Z77.legPos_lip}, \lean{HPC.Fold.Z77.legPos_eq_zero}.\\
Folds from fitting (Lemma~\ref{lem:fold-of-fit}) & \lean{HPC.Fold.Z77.fold_of_fit}, \lean{HPC.Fold.Z77.fold_of_fit_zero}.\\
Up phase (Lemma~\ref{lem:up-phase}) & \lean{HPC.Fold.Z77.up_step} (with \lean{HPC.Fold.Z77.transfer} and \lean{HPC.Fold.Z77.exists_armVal}).\\
Down phase (Lemma~\ref{lem:down-phase}) & \lean{HPC.Fold.Z77.down_step}.\\
Up core (Lemma~\ref{lem:up-core}) & \lean{HPC.Fold.Z77.up_core}.\\
Down core (Lemma~\ref{lem:down-core}) & \lean{HPC.Fold.Z77.down_core}.\\
Up states are not stuck (Lemma~\ref{lem:nofail-up}) & \lean{HPC.Fold.Z77.nofail_up}.\\
Down states after an up phase are not stuck (Lemma~\ref{lem:nofail-down}) & \lean{HPC.Fold.Z77.nofail_down}.\\
First direction (Lemma~\ref{lem:first-direction}) & \lean{HPC.Fold.Z77.first_direction}; the two conditions are \lean{HPC.Fold.Z77.OKBelow} and \lean{HPC.Fold.Z77.OKAbove}.\\
Main induction (Proposition~\ref{prop:main-up}) & \lean{HPC.Fold.Z77.main_up}.\\
Zigzag theorem (Theorem~\ref{thm:zigzag}) & \lean{HPC.Fold.Z77.greedy_fit_below}, \lean{HPC.Fold.Z77.greedy_fit_above}, and \lean{HPC.Fold.Z77.greedy_fit} (which uses Lemma~\ref{lem:first-direction} to choose $\epsilon$).\\
Bing's Theorem~3, three-link form (Corollary~\ref{cor:theorem3}) & \lean{HPC.Fold.theorem3}, via \lean{HPC.Fold.Z77.s3_zig} and \lean{HPC.exists_blocks_of_walk}.\\
Folds with a fake end link (Section~\ref{sec:fold-fake}) & \lean{HPC.Fold.B79.EvUp} and \lean{HPC.Fold.B79.EvDown}; reflecting $f$ exchanges them (\lean{HPC.Fold.B79.evUp_reflect}).\\
Crooked extensions (Lemma~\ref{lem:crooked-ext}) & \lean{HPC.Fold.B79.crkOn_lowExt}, \lean{HPC.Fold.B79.crkOn_highExt} (on sequences re-indexed so that the base is at $0$).\\
Fold with a fake low end (Lemma~\ref{lem:fake-low}) & \lean{HPC.Fold.B79.finfold_low}, from \lean{HPC.Fold.B79.fold_low_end_mode} (the version on $\Z$-indexed sequences; without the choice of $\epsilon$ it is \lean{HPC.Fold.B78.fold_low_end}). With $\epsilon$ chosen by Lemma~\ref{lem:first-direction} it is \lean{HPC.Fold.B79.finfold_low_any}.\\
Fold with a fake high end (Lemma~\ref{lem:fake-high}) & \lean{HPC.Fold.B79.finfold_high}, from \lean{HPC.Fold.B79.fold_high_end_mode}.\\
Lemma~\ref{lem:meet-low} & \lean{HPC.Fold.B79.meet_low_gen}, \lean{HPC.Fold.B79.maxIdx_near_zero}; the two sequences and their properties are collected in \lean{HPC.Fold.B79.seq_facts}.\\
Exclusion (Lemma~\ref{lem:exclusion}) & \lean{HPC.Fold.B79.pair_exclusion}.\\
Prefix fold (Proposition~\ref{prop:prefix-fold}) & \lean{HPC.Fold.B79.pair_prefix_fold}; the conclusion is \lean{HPC.Fold.B79.PrefixFold}.\\
Suffix fold (Corollary~\ref{cor:suffix-fold}) & \lean{HPC.Fold.B79.pair_suffix_fold}.\\
Splice (Proposition~\ref{prop:splice}) & \lean{HPC.Fold.B79.splice_schedule_pair} (the version with $\lambda=\mu=\lo$ and without (c), (d) is \lean{HPC.Fold.B78.splice_schedule}); the passage from levels to links is \lean{HPC.Fold.B78.link_sub_of_level}.\\
From a schedule to the block map (Lemma~\ref{lem:assemble}) & \lean{HPC.Fold.B79.assemble_pair}.\\
Two windows (Lemma~\ref{lem:windows}) & \lean{HPC.Fold.B79.pair_windows}.\\[4pt]
\multicolumn{2}{@{}l}{\emph{Chapter~\ref{ch:descriptions}}}\\
Bing 1951, Theorem~1, body of the proof (Theorem~\ref{thm:desc-exists}) & (a) is \lean{HPC.Plane.PA51.exists_description_sub}; (b) is \lean{HPC.Plane.exists_description_of_chainable_HI} (in \lean{HPC/Plane/PseudoArc51.lean}), via \lean{HPC.Plane.PA51.lift_description}.\\
Bing 1951 (Theorem~\ref{thm:eventually-crooked}) & \lean{HPC.Plane.PA51.eventually_crookedIn} (in \lean{HPC/Plane/PA51Crook.lean}).\\
Lemma~\ref{lem:fail-data} & \lean{HPC.Plane.PA51.FailData} (which records the property (d) through small sets, the links, none inside $C_{k-1}$, resp.\ $C_{h+1}$), \lean{HPC.Plane.PA51.failData_of_fail_lt} (the case $i<j$) and \lean{HPC.Plane.PA51.failData_of_not_crooked}.\\
Lemma~\ref{lem:limit-tools} & \lean{HPC.Plane.PA51.chainedAt_union_singleton}, \lean{HPC.Plane.PA51.limSup_miss} (with \lean{HPC.Plane.PA51.approx_of_mem_limSup}), \lean{HPC.Plane.PA51.exists_limSup_near}.\\
Lemma~\ref{lem:side} & \lean{HPC.Plane.PA51.subset_side_of_miss}, with \lean{HPC.Plane.PA51.lowU} and \lean{HPC.Plane.PA51.highU}.\\
Proposition~\ref{prop:limit-false} & \lean{HPC.Plane.PA51.limit_false}.\\
Definition~\ref{def:predesc} & \lean{HPC.Plane.PA51.PreDesc}.\\
Lemma~\ref{lem:crooked-next} & \lean{HPC.Plane.PA51.exists_crooked_next}.\\
Proposition~\ref{prop:predesc-exists} & \lean{HPC.Plane.PA51.exists_preDesc}.\\
Lemma~\ref{lem:predesc-basic} & \lean{HPC.Plane.PA51.PreDesc.refines_of_le}, \lean{HPC.Plane.PA51.PreDesc.crooked_of_lt}, \lean{HPC.Plane.PA51.PreDesc.eventually_five_le}.\\
Reaching both ends (Lemma~\ref{lem:meets-ends}) & \lean{HPC.Plane.PA51.PreDesc.exists_subchain_meets_ends}, with the limit step \lean{HPC.Plane.PA51.PreDesc.false_of_disjoint_limit}.\\
One step (Lemma~\ref{lem:desc-step}) & \lean{HPC.Plane.PA51.PreDesc.exists_step}; the analogous step for descriptions is \lean{HPC.Description.exists_step} (in \lean{HPC/PseudoArc/T13Step.lean}).\\
Theorem~\ref{thm:desc-univ} & \lean{HPC.Plane.PA51.PreDesc.exists_description}, and with Proposition~\ref{prop:predesc-exists}, \lean{HPC.Plane.PA51.exists_description_univ}.\\
Lemma~\ref{lem:ext} & \lean{HPC.Plane.PA51.ext}, \lean{HPC.Plane.PA51.mem_ext_iff}, \lean{HPC.Plane.PA51.ext_disjoint}, \lean{HPC.Plane.PA51.ext_mono}, \lean{HPC.Plane.PA51.ext_subset_thickening}, \lean{HPC.Plane.PA51.diam_ext_le}, \lean{HPC.Plane.PA51.exists_closure_ext_subset}.\\[4pt]
\multicolumn{2}{@{}l}{\emph{Chapter~\ref{ch:main}}}\\
Pseudo-arc (\vone{Definition~2.14}) & \lean{HPC.Plane.IsPseudoArc}: the conjunction of \lean{HPC.Plane.IsSnakeLike} (which contains ``$M$ is a continuum''), \lean{HPC.Plane.IsHereditarilyIndecomposable} and \lean{Set.Nontrivial}.\\
Bing 1951, Theorem~1, the ``in fact'' clause (Theorem~\ref{thm:unique-points}) & \lean{HPC.Plane.theorem1_HIC}. Formally the described sets are subsets $\Delta_1.M\subseteq M$ and $\Delta_2.M\subseteq M'$ of the spaces $M$, $M'$, equal to the whole spaces; the homeomorphism of \lean{HPC.theorem12} is composed with the identifications $M\cong\Delta_1.M$ and $\Delta_2.M\cong M'$.\\
Uniqueness of the pseudo-arc (Theorem~\ref{thm:unique}) & \lean{HPC.Plane.theorem1_HIC'}; the same statement with the three defining properties written out is \lean{HPC.Plane.theorem1_snake} (Bing, \emph{Snake-like continua}, Theorem~1).\\
Homogeneity of the pseudo-arc (Theorem~\ref{thm:homog}) & \lean{HPC.Plane.homogeneous_of_isPseudoArc}.\\
Corollary~\ref{cor:plane-unique} & \lean{HPC.Plane.isPseudoArc_plane} and \lean{HPC.Plane.pseudoArc_unique}.\\
Lemma~\ref{lem:sub-pseudoarc} & \lean{HPC.Plane.isPseudoArc_of_subcontinuum}.\\
Corollary~\ref{cor:sub-homeo} & \lean{HPC.Plane.homeomorphic_subcontinuum}.\\
Definition~\ref{def:endpoint} & \lean{HPC.Plane.IsEndPoint}, \lean{HPC.Plane.OppositeEndPoints} (in \lean{HPC/Plane/EndPoints51.lean}).\\
Proposition~\ref{prop:opposite} & \lean{HPC.Plane.oppositeEndPoints_of_composants}.\\
Lemma~\ref{lem:endpoint-nested} & \lean{HPC.Plane.condB_of_endPoint}.\\
Lemma~\ref{lem:not-HI-not-end} & \lean{HPC.Plane.exists_not_endPoint_of_not_HI}.\\
Bing 1951, \emph{Snake-like continua}, Theorem~16 (Theorem~\ref{thm:endpoints}) & \lean{HPC.Plane.isPseudoArc_iff_endPoints}.\\
\end{longtable}


\end{document}
